% Statistiques — Mathématiques 1ère A — Cours signé OnBuch+
% Programme officiel MINESEC. Compiler avec : tectonic NA06-statistiques-a.tex
\newcommand{\DOCMATIERE}{Mathématiques}
\newcommand{\DOCNIVEAU}{1ère A}
\newcommand{\DOCMODULE}{Organisation et gestion des données}
\newcommand{\DOCLECON}{NA06}
\newcommand{\DOCTITRE}{Statistiques}
% OnBuch+ — préambule « cours signés » ENRICHI (compilé avec tectonic / XeLaTeX)
% Système visuel « Pop » d'OnBuch+ : fond crème (jamais blanc), bordures encre
% épaisses, ombres « dures » décalées, titres Archivo Black en capitales.
% Version MATHÉMATIQUES : graphes (pgfplots/tikz), boîtes pédagogiques
% (définition, propriété, méthode, attention, le savais-tu, exercice résolu,
% exercice, TP, bilan), page de garde et signature OnBuch+ ancrée en bas.
%
% Macros attendues dans main.tex AVANT \input{preamble} :
%   \DOCMATIERE  \DOCNIVEAU  \DOCMODULE  \DOCLECON (numéro)  \DOCTITRE
\documentclass[11pt]{article}

\usepackage{fontspec}
\usepackage[french]{babel}
\usepackage[a4paper,margin=2cm,top=3.4cm,bottom=2.8cm,headheight=1.4cm,footskip=1.3cm]{geometry}
\usepackage{amsmath,amssymb,amsfonts}
\usepackage{mathtools} % \xmapsto, \coloneqq... souvent utilisés par les modèles
\usepackage{cancel} % simplifications barrées
% \abs{z} (module, valeur absolue) et \norm{u} (norme), utilisés par les modèles
\providecommand{\abs}{}\let\abs\relax\DeclarePairedDelimiter{\abs}{\lvert}{\rvert}
\providecommand{\norm}{}\let\norm\relax\DeclarePairedDelimiter{\norm}{\lVert}{\rVert}
\usepackage[table]{xcolor} % [table] : fournit aussi \rowcolors (alternance de lignes)
\usepackage{pagecolor}
\usepackage{tikz}
\usetikzlibrary{arrows.meta,positioning,calc,shapes.geometric,shapes.misc,shapes.arrows,decorations.pathmorphing,decorations.pathreplacing,decorations.markings,patterns,shadows,mindmap,trees,fit,backgrounds,matrix,intersections,angles,quotes}
\usepackage{pgfplots}
\usepackage{pgf-pie} % diagrammes circulaires \pie (statistiques)
\pgfplotsset{compat=1.18}
\usepgfplotslibrary{fillbetween,groupplots} % aires entre courbes (calcul intégral)
\usepackage{enumitem}
\usepackage{fancyhdr}
% [most] charge toutes les bibliothèques tcolorbox (listings, minted, raster,
% documentation...) et gonfle inutilement la mémoire TeX sur un document long
% (~300 boîtes) : on ne charge que ce qui est réellement utilisé.
\usepackage[skins,breakable]{tcolorbox}
\usepackage{graphicx}
\usepackage{colortbl}
\usepackage{booktabs}
\usepackage{tabularx}
\usepackage{diagbox} % case d'en-tête diagonale des tableaux à double entrée
% Colonnes extensibles centrée / gauche / droite (C, L, R), souvent utilisées par les modèles
\newcolumntype{C}{>{\centering\arraybackslash}X}
\newcolumntype{L}{>{\raggedright\arraybackslash}X}
\newcolumntype{R}{>{\raggedleft\arraybackslash}X}
\usepackage{array}
\usepackage{multicol}
\usepackage{siunitx}
% parse-numbers=false : accepte aussi \SI{2,5 \times 10^{-3}}{...}
\sisetup{locale=FR,output-decimal-marker={,},group-minimum-digits=4,parse-numbers=false}
\DeclareSIUnit\ohm{\ensuremath{\symup{\Omega}}} % U+2126 absent de la police
\usepackage{qrcode}
% \degree : commande fréquemment utilisée par les modèles pour un angle ou une
% température en dehors de \SI{}{} (non fournie par les paquets chargés).
\providecommand{\degree}{^{\circ}}
% \qed (fin de démonstration), utilisé par les modèles sans amsthm
\providecommand{\qed}{\hfill$\square$}
% \barre : double barre des valeurs interdites dans un tableau de variations (tkz-tab)
\providecommand{\barre}{\ensuremath{\|}}
% environnement proof (amsthm non chargé) utilisé par les modèles
\ifdefined\proof\else\newenvironment{proof}[1][Démonstration]{\par\noindent\textit{#1.}\ }{\qed\par}\fi
\usepackage{lastpage}
\usepackage{hyperref}
\hypersetup{hidelinks}

% ============================================================
% Palette « Pop » OnBuch+ — mêmes hex que la classe P (plus_ui.dart)
% ============================================================
\definecolor{poporange}{HTML}{F59321}
\definecolor{popdark}{HTML}{E07A0C}
\definecolor{popcream}{HTML}{FFF6E8}
\definecolor{popink}{HTML}{1C1714}
\definecolor{poppurple}{HTML}{7A5AE0}
\definecolor{popgreen}{HTML}{2E9E6A}
\definecolor{poppink}{HTML}{E0457B}
\definecolor{popblue}{HTML}{2F7DE1}
\definecolor{popgold}{HTML}{C98A12}
\definecolor{popmuted}{HTML}{8A7F76}
\definecolor{popline}{HTML}{EADFD2}
\definecolor{popgray}{HTML}{8A7F76}
\colorlet{popgrey}{popgray}
% Alias tolérants (le modèle nomme parfois les couleurs différemment)
\colorlet{popwhite}{white}
\colorlet{popblack}{popink}
\colorlet{popred}{poppink}
\colorlet{popyellow}{popgold}
% Teintes claires pour fonds de tableaux / zones
\colorlet{popblueL}{popblue!10!white}
\colorlet{poporangeL}{poporange!14!white}
\colorlet{popgreenL}{popgreen!12!white}
\colorlet{poppurpleL}{poppurple!12!white}
\colorlet{poppinkL}{poppink!12!white}
\colorlet{popgoldL}{popgold!14!white}

\pagecolor{popcream}

% ============================================================
% Typographie
% ============================================================
\defaultfontfeatures{Ligatures=TeX}
\setmainfont{PlusJakartaSans-Medium.ttf}[Path=../../../fonts/,BoldFont=PlusJakartaSans-Bold.ttf]
% Maths en sans-serif (Fira Math) avec chiffres et lettres droites de Plus
% Jakarta, pour que les formules se fondent dans le texte.
\usepackage[math-style=ISO,bold-style=ISO]{unicode-math}
\setmathfont{FiraMath-Regular.otf}[Path=../../../fonts/]
\setmathfont{PlusJakartaSans-Medium.ttf}[Path=../../../fonts/,range={up/num,up/{Latin,latin},\mathup}]
\usepackage{icomma} % virgule décimale sans espace en mode math
% Caractères mathématiques Unicode absents des polices de texte (Plus Jakarta,
% Archivo Black) : rendus par leur équivalent mathématique, en texte comme en
% formule — sinon ils s'affichent en carré vide (ex. « Arithmétique dans ℤ »).
\usepackage{newunicodechar}
\newunicodechar{ℝ}{\ensuremath{\mathbb{R}}}
\newunicodechar{ℤ}{\ensuremath{\mathbb{Z}}}
\newunicodechar{ℂ}{\ensuremath{\mathbb{C}}}
\newunicodechar{ℕ}{\ensuremath{\mathbb{N}}}
\newunicodechar{ℚ}{\ensuremath{\mathbb{Q}}}
\newunicodechar{𝔻}{\ensuremath{\mathbb{D}}}
\newunicodechar{α}{\ensuremath{\alpha}}
\newunicodechar{β}{\ensuremath{\beta}}
\newunicodechar{γ}{\ensuremath{\gamma}}
\newunicodechar{δ}{\ensuremath{\delta}}
\newunicodechar{ε}{\ensuremath{\varepsilon}}
\newunicodechar{θ}{\ensuremath{\theta}}
\newunicodechar{λ}{\ensuremath{\lambda}}
\newunicodechar{μ}{\ensuremath{\mu}}
\newunicodechar{σ}{\ensuremath{\sigma}}
\newunicodechar{τ}{\ensuremath{\tau}}
\newunicodechar{φ}{\ensuremath{\varphi}}
\newunicodechar{ψ}{\ensuremath{\psi}}
\newunicodechar{ω}{\ensuremath{\omega}}
\newunicodechar{Σ}{\ensuremath{\Sigma}}
% majuscules produites par \MakeUppercase dans les titres (α-aminés -> Α)
\newunicodechar{Α}{\ensuremath{\alpha}}
\newunicodechar{Β}{\ensuremath{\beta}}
\newunicodechar{Γ}{\ensuremath{\Gamma}}
\newunicodechar{Δ}{\ensuremath{\Delta}}
\newunicodechar{Ε}{\ensuremath{\varepsilon}}
\newunicodechar{Θ}{\ensuremath{\Theta}}
\newunicodechar{Λ}{\ensuremath{\Lambda}}
\newunicodechar{Μ}{\ensuremath{\mu}}
\newunicodechar{Τ}{\ensuremath{\tau}}
\newunicodechar{Φ}{\ensuremath{\Phi}}
\newunicodechar{Ψ}{\ensuremath{\Psi}}
\newunicodechar{Ω}{\ensuremath{\Omega}}
\newunicodechar{↦}{\ensuremath{\mapsto}}
\newunicodechar{∈}{\ensuremath{\in}}
\newunicodechar{∉}{\ensuremath{\notin}}
\newunicodechar{∧}{\ensuremath{\wedge}}
\newunicodechar{⇔}{\ensuremath{\Leftrightarrow}}
\newunicodechar{⟺}{\ensuremath{\Longleftrightarrow}}
\newunicodechar{⇒}{\ensuremath{\Rightarrow}}
\newunicodechar{⟹}{\ensuremath{\Longrightarrow}}
\newunicodechar{∘}{\ensuremath{\circ}}
\newunicodechar{∪}{\ensuremath{\cup}}
\newunicodechar{∩}{\ensuremath{\cap}}
\newunicodechar{≡}{\ensuremath{\equiv}}
\newunicodechar{⊂}{\ensuremath{\subset}}
\newunicodechar{⊥}{\ensuremath{\perp}}
\newunicodechar{∃}{\ensuremath{\exists}}
\newunicodechar{∀}{\ensuremath{\forall}}
\newunicodechar{∛}{\ensuremath{\sqrt[3]{\,}}}
\newunicodechar{∜}{\ensuremath{\sqrt[4]{\,}}}
\newunicodechar{⃗}{\ensuremath{{}^{\to}}}
\newunicodechar{ⁿ}{\ifmmode^{n}\else\textsuperscript{\ensuremath{n}}\fi}
\newunicodechar{ˣ}{\ifmmode^{x}\else\textsuperscript{\ensuremath{x}}\fi}
\newunicodechar{ᵇ}{\ifmmode^{b}\else\textsuperscript{\ensuremath{b}}\fi}
\newunicodechar{ʳ}{\ifmmode^{r}\else\textsuperscript{\ensuremath{r}}\fi}
\newunicodechar{ᵘ}{\ifmmode^{u}\else\textsuperscript{\ensuremath{u}}\fi}
\newunicodechar{ʸ}{\ifmmode^{y}\else\textsuperscript{\ensuremath{y}}\fi}
\newunicodechar{ᵃ}{\ifmmode^{a}\else\textsuperscript{\ensuremath{a}}\fi}
\newunicodechar{ᵉ}{\ifmmode^{e}\else\textsuperscript{\ensuremath{e}}\fi}
\newunicodechar{ᵏ}{\ifmmode^{k}\else\textsuperscript{\ensuremath{k}}\fi}
\newunicodechar{ᵖ}{\ifmmode^{p}\else\textsuperscript{\ensuremath{p}}\fi}
\newunicodechar{ᴺ}{\ifmmode^{N}\else\textsuperscript{\ensuremath{N}}\fi}
\newunicodechar{ᶜ}{\ifmmode^{c}\else\textsuperscript{\ensuremath{c}}\fi}
\newunicodechar{ʰ}{\ifmmode^{h}\else\textsuperscript{\ensuremath{h}}\fi}
\newunicodechar{ᵠ}{\ifmmode^{q}\else\textsuperscript{\ensuremath{q}}\fi}
\newunicodechar{ₙ}{\ifmmode_{n}\else\textsubscript{\ensuremath{n}}\fi}
\newunicodechar{ₐ}{\ifmmode_{a}\else\textsubscript{\ensuremath{a}}\fi}
\newunicodechar{ᵢ}{\ifmmode_{i}\else\textsubscript{\ensuremath{i}}\fi}
\newunicodechar{ᵣ}{\ifmmode_{r}\else\textsubscript{\ensuremath{r}}\fi}
\newunicodechar{ₖ}{\ifmmode_{k}\else\textsubscript{\ensuremath{k}}\fi}
\newunicodechar{ᵦ}{\ifmmode_{\beta}\else\textsubscript{\ensuremath{\beta}}\fi}
\newunicodechar{ₓ}{\ifmmode_{x}\else\textsubscript{\ensuremath{x}}\fi}
\newfontfamily\popdisplay{ArchivoBlack-Regular.ttf}[Path=../../../fonts/]
\newfontfamily\poplogo{SpaceGrotesk-Bold.ttf}[Path=../../../fonts/]
\newfontfamily\popsemi{PlusJakartaSans-SemiBold.ttf}[Path=../../../fonts/]
\newfontfamily\popxbold{PlusJakartaSans-ExtraBold.ttf}[Path=../../../fonts/]
\newfontfamily\popxboldls{PlusJakartaSans-ExtraBold.ttf}[Path=../../../fonts/,LetterSpace=3.0]

\setlist[enumerate]{leftmargin=1.7em,itemsep=5pt,topsep=4pt}
\setlist[itemize]{leftmargin=1.7em,itemsep=4pt,topsep=4pt,label={\color{poporange}\textbullet}}
\setlength{\parindent}{0pt}
\setlength{\parskip}{5pt}
\renewcommand{\arraystretch}{1.25}

% pgfplots : style Pop par défaut
\pgfplotsset{
  every axis/.append style={axis line style={popink,line width=1pt},tick style={popink},
    grid style={popline,line width=0.5pt},label style={font=\small},tick label style={font=\footnotesize}},
  popcurve/.style={poporange,line width=1.8pt,no markers,smooth},
}
% Même style disponible dans un \draw TikZ simple (hors pgfplots)
\tikzset{popcurve/.style={poporange,line width=1.8pt}}

% ============================================================
% Pill / squiggle
% ============================================================
\newcommand{\poppill}[3]{%
  \tikz[baseline=-0.55ex]\node[draw=popink,line width=1.2pt,rounded corners=2.6pt,%
    fill=#2,inner xsep=6.5pt,inner ysep=3pt]{%
    \popxboldls\fontsize{7}{7}\selectfont\color{#3}\MakeUppercase{#1}};%
}
\newcommand{\popsquiggle}[1]{%
  \tikz[baseline=-0.4ex]{
    \draw[#1,line width=2.6pt,line cap=round]
      plot [smooth] coordinates {(0,0.09) (0.34,0.01) (0.68,0.21) (1.02,0.01) (1.36,0.10)};
  }%
}
% Mot-clé surligné (marqueur orange sous le texte)
\newcommand{\cle}[1]{{\popxbold\color{popdark}#1}}

% ============================================================
% En-tête / pied
% ============================================================
\pagestyle{fancy}
\fancyhf{}
\renewcommand{\headrulewidth}{0pt}
\renewcommand{\footrulewidth}{0pt}
\lhead{\raisebox{-0.15ex}{\poplogo\fontsize{13}{13}\selectfont\color{popink}OnBuch\color{poporange}+}}
\rhead{\poppill{\DOCMATIERE\ \textbullet\ \DOCNIVEAU\ \textbullet\ Leçon \DOCLECON}{white}{popink}}
\lfoot{\popxbold\fontsize{7.6}{7.6}\selectfont\color{popmuted}\MakeUppercase{Cours signé OnBuch+}\ \textbullet\ {\popsemi\DOCTITRE}}
\rfoot{\tikz[baseline=-0.6ex]\node[rounded corners=8.5pt,fill=popink,minimum height=17pt,minimum width=17pt,inner xsep=5pt,inner sep=0]{\popxbold\fontsize{8}{8}\selectfont\color{popcream}\thepage\,/\,\pageref*{LastPage}};}

% ============================================================
% Titres
% ============================================================
\newcommand{\chaptitre}[2]{%
  \begin{tcolorbox}[enhanced,colback=poporange,colframe=popink,boxrule=2.6pt,arc=11pt,
      drop shadow=popink,left=15pt,right=15pt,top=12pt,bottom=11pt,before skip=3pt,after skip=18pt]
    \poppill{#2}{popink}{popcream}\par\vspace{9pt}
    {\raggedright\hyphenpenalty=10000\exhyphenpenalty=10000\popdisplay\fontsize{22}{24}\selectfont\color{popcream}\MakeUppercase{#1}\par}
  \end{tcolorbox}
}
% Section numérotée : pastille orange avec le numéro + titre Archivo
\newcounter{popsec}
\newcommand{\coursec}[1]{%
  \stepcounter{popsec}\setcounter{popsub}{0}%
  \par\vspace{16pt}\noindent
  \begin{minipage}{\linewidth}
  \tikz[baseline=-0.7ex]\node[draw=popink,line width=1.6pt,fill=poporange,rounded corners=4pt,minimum size=22pt,inner sep=0]{\popdisplay\fontsize{12}{12}\selectfont\color{popcream}\thepopsec};%
  \hspace{8pt}{\popdisplay\fontsize{15}{17}\selectfont\color{popink}\MakeUppercase{#1}}\par
  \vspace{4pt}\noindent{\color{poporange}\rule{36pt}{3.6pt}}
  \end{minipage}\par\nopagebreak\vspace{8pt}
}
\newcounter{popsub}
\newcommand{\courssub}[1]{%
  \stepcounter{popsub}%
  \par\vspace{9pt}\noindent{\popxbold\fontsize{11.5}{13}\selectfont\color{popink}{\color{poporange}\thepopsec.\thepopsub}\hspace{6pt}#1}\par\nopagebreak\vspace{4pt}
}

% ============================================================
% Boîtes pédagogiques — même recette : fond blanc, bordure épaisse colorée,
% ombre dure encre, badge plein. Titre optionnel [#1].
% ============================================================
\tcbset{popbase/.style={breakable,enhanced,colback=white,boxrule=2.3pt,arc=9pt,
  shadow={3pt}{-3pt}{0pt}{popink},left=14pt,right=14pt,top=11pt,bottom=11pt,before skip=11pt,after skip=15pt,
  fontupper=\popsemi\fontsize{10.6}{14.5}\selectfont\color{popink},
  fontlower=\popsemi\fontsize{10.6}{14.5}\selectfont\color{popink}}}
% Coche dessinée (le glyphe ✓ n'existe pas dans les polices du document)
\newcommand{\popcheck}[1]{\tikz[baseline=-0.5ex]\draw[#1,line width=1.6pt,line cap=round,line join=round](-0.13,0.0)--(-0.04,-0.09)--(0.13,0.1);}
\newcommand{\popboxhead}[3]{\poppill{#1}{#2}{white}\ifx\relax#3\relax\else\hspace{6pt}{\popxbold\fontsize{10.6}{12}\selectfont\color{popink}#3}\fi\par\vspace{7pt}}

\newtcolorbox{aretenir}[1][]{popbase,colframe=popgreen,before upper={\popboxhead{À retenir}{popgreen}{#1}}}
\newtcolorbox{exemplebox}[1][]{popbase,colframe=poporange,before upper={\popboxhead{Exemple}{poporange}{#1}}}
\newtcolorbox{definition}[1][]{popbase,colframe=popblue,before upper={\popboxhead{Définition}{popblue}{#1}}}
\newtcolorbox{propriete}[1][]{popbase,colframe=poppurple,before upper={\popboxhead{Propriété}{poppurple}{#1}}}
\newtcolorbox{methode}[1][]{popbase,colframe=popgold,before upper={\popboxhead{Méthode}{popgold}{#1}}}
\newtcolorbox{attention}[1][]{popbase,colframe=poppink,before upper={\popboxhead{Attention}{poppink}{#1}}}
\newtcolorbox{experience}[1][]{popbase,colframe=popblue,colback=popblueL,before upper={\popboxhead{Expérience}{popblue}{#1}}}
% « Le savais-tu ? » : carte violette pleine (contexte, culture, Cameroun)
\newtcolorbox{savaistu}[1][]{popbase,colback=poppurple,colframe=popink,
  fontupper=\popsemi\fontsize{10.4}{14.2}\selectfont\color{white},
  before upper={\poppill{Le savais-tu\,?}{white}{poppurple}\ifx\relax#1\relax\else\hspace{6pt}{\popxbold\color{white}#1}\fi\par\vspace{7pt}}}
% Exercice résolu : énoncé en haut, \tcblower puis la solution sur fond crème
\newcounter{popexo}\newcounter{popexr}
\newtcolorbox{exoresolu}[1][]{popbase,colframe=poporange,colbacklower=popcream,
  segmentation style={popink,line width=1.2pt,dashed},
  before upper={\stepcounter{popexr}\popboxhead{Exercice résolu \thepopexr}{poporange}{#1}},
  before lower={\poppill{Solution}{popink}{popcream}\par\vspace{6pt}}}
% Exercice (non résolu en place) — la correction est dans la partie Corrigés
\newtcolorbox{exercice}[1][]{popbase,colframe=popink,
  before upper={\stepcounter{popexo}\popboxhead{Exercice \thepopexo}{popink}{#1}}}
% Corrigé
\newtcolorbox{corrige}[1][]{popbase,colframe=popgreen,colback=popgreenL,
  before upper={\popboxhead{Corrigé}{popgreen}{#1}}}
% Objectifs / prérequis / bilan
\newtcolorbox{objectifs}{popbase,colframe=popink,colback=poporangeL,
  before upper={\popboxhead{Objectifs de la leçon}{popink}{}}}
\newtcolorbox{prerequis}{popbase,colframe=popink,colback=popblueL,
  before upper={\popboxhead{Prérequis}{popblue}{}}}
\newtcolorbox{bilan}{popbase,colframe=popink,colback=white,boxrule=2.8pt,
  before upper={\popboxhead{Fiche bilan}{poporange}{L'essentiel à connaître pour l'examen}}}
% Figure encadrée avec légende
\newtcolorbox{popfigure}[1][]{enhanced,breakable=false,colback=white,colframe=popline,boxrule=1.4pt,arc=8pt,
  left=10pt,right=10pt,top=10pt,bottom=8pt,before skip=10pt,after skip=12pt,halign=center,
  lower separated=false,halign lower=center,
  fontlower=\popsemi\fontsize{9}{11.5}\selectfont\color{popmuted},
  #1}
\newcommand{\legende}[1]{\par\vspace{6pt}{\popsemi\fontsize{9}{11.5}\selectfont\color{popmuted}#1}}

% ============================================================
% Signature OnBuch+
% ============================================================
\newcommand{\poplogoinline}[1]{{\poplogo\fontsize{#1}{#1}\selectfont\color{popink}OnBuch\color{poporange}+}}
\newcommand{\signatureonbuch}{%
  \begin{tcolorbox}[enhanced,colback=popink,colframe=popink,boxrule=2.2pt,arc=10pt,
      drop shadow=poporange,left=14pt,right=14pt,top=10pt,bottom=10pt,before skip=0pt,after skip=0pt]
    \tikz[baseline=-0.7ex]\node[circle,fill=popgreen,draw=popcream,line width=1.3pt,minimum size=22pt,inner sep=0]{\popcheck{white}};%
    \hspace{9pt}\begin{minipage}[c]{0.62\linewidth}
      {\popxbold\fontsize{12}{13}\selectfont\color{popcream}Signé\ \textbullet\ {\poplogo OnBuch\color{poporange}+}}\par\vspace{2pt}
      {\raggedright\popsemi\fontsize{8}{10.5}\selectfont\color{popline}\MakeUppercase{Équipe pédagogique OnBuch+}\\\MakeUppercase{Conforme au programme officiel MINESEC}\par}
    \end{minipage}\hfill
    {\poplogo\fontsize{22}{22}\selectfont\color{popcream}OnBuch\color{poporange}+}
  \end{tcolorbox}%
}

% ============================================================
% Page de garde
% #1 titre · #2 matière · #3 niveau · #4 module · #5 n° leçon · #6 durée ·
% #7 description / objectifs (texte ou itemize)
% ============================================================
\newcommand{\pagegarde}[7]{%
  \thispagestyle{empty}
  \newgeometry{a4paper,margin=1.7cm,top=1.6cm,bottom=1.4cm}%
  \begin{tikzpicture}[remember picture,overlay]
    \fill[poporange!18] ([xshift=-3.2cm,yshift=-3.6cm]current page.north east) circle (4.6cm);
    \fill[poppurple!14] ([xshift=2.4cm,yshift=7.6cm]current page.south west) circle (3.8cm);
  \end{tikzpicture}%
  {\poplogo\fontsize{30}{30}\selectfont\color{popink}OnBuch\color{poporange}+}\hfill
  \raisebox{0.6ex}{\poppill{Cours signé \textbullet\ Premium}{popink}{popcream}}\par
  \vspace{1pt}\popsquiggle{poppurple}\par
  \vspace{8pt}
  \begin{tcolorbox}[enhanced,colback=poporange,colframe=popink,boxrule=2.8pt,arc=13pt,
      drop shadow=popink,left=18pt,right=18pt,top=12pt,bottom=13pt,after skip=10pt]
    \poppill{#2 \textbullet\ #3}{popink}{popcream}\hspace{5pt}\poppill{#4}{white}{popink}\par\vspace{8pt}
    {\popdisplay\fontsize{14}{14}\selectfont\color{popink}LEÇON #5}\par\vspace{4pt}
    {\raggedright\hyphenpenalty=10000\exhyphenpenalty=10000\popdisplay\fontsize{25}{27}\selectfont\color{popcream}\MakeUppercase{#1}\par}
    \vspace{8pt}
    {\popxbold\fontsize{9.5}{9.5}\selectfont\color{popink}\MakeUppercase{Durée conseillée : #6}}
  \end{tcolorbox}
  \begin{tcolorbox}[enhanced,colback=white,colframe=popink,boxrule=2.2pt,arc=10pt,
      drop shadow=popink,left=16pt,right=16pt,top=10pt,bottom=10pt,after skip=8pt]
    \poppill{Dans cette leçon}{poppurple}{white}\par\vspace{5pt}
    \popsemi\fontsize{9.3}{12.6}\selectfont\color{popink}#7
  \end{tcolorbox}
  \noindent\begin{minipage}[t]{0.487\linewidth}
    \begin{tcolorbox}[enhanced,colback=popgreen,colframe=popink,boxrule=2.2pt,arc=10pt,
        drop shadow=popink,left=11pt,right=11pt,top=10pt,bottom=10pt,height=3.1cm,valign=center]
      \tcbox[colback=white,colframe=popink,boxrule=1.3pt,arc=5pt,boxsep=0pt,nobeforeafter,
        left=3pt,right=3pt,top=3pt,bottom=3pt,valign=center]{\qrcode[height=1.9cm]{https://play.google.com/store/apps/details?id=com.luvvix.onbuch&hl=fr}}%
      \hspace{8pt}\begin{minipage}[c]{0.5\linewidth}
        {\popdisplay\fontsize{11}{12.5}\selectfont\color{white}\MakeUppercase{Télécharge OnBuch}}\par\vspace{4pt}
        {\raggedright\popsemi\fontsize{8.4}{11}\selectfont\color{white}Gratuit \textbullet\ Google Play\\Tuteur IA, annales, concours, résultats\par}
      \end{minipage}
    \end{tcolorbox}
  \end{minipage}\hfill
  \begin{minipage}[t]{0.487\linewidth}
    \begin{tcolorbox}[enhanced,colback=poppurple,colframe=popink,boxrule=2.2pt,arc=10pt,
        drop shadow=popink,left=11pt,right=11pt,top=10pt,bottom=10pt,height=3.1cm,valign=center]
      \tikz[baseline=-0.7ex]\node[circle,fill=white,draw=popink,line width=1.3pt,minimum size=1.6cm,inner sep=0]{\popdisplay\fontsize{24}{24}\selectfont\color{poppurple}+};%
      \hspace{8pt}\begin{minipage}[c]{0.6\linewidth}
        {\popdisplay\fontsize{11}{12.5}\selectfont\color{white}\MakeUppercase{Passe à OnBuch+}}\par\vspace{4pt}
        {\raggedright\popsemi\fontsize{8.4}{11}\selectfont\color{white}Tous les cours signés, Tuteur IA illimité, fiches PDF — onglet OnBuch+ de l'app\par}
      \end{minipage}
    \end{tcolorbox}
  \end{minipage}
  \vspace{8pt}
  \popcontenu
  \vfill
  \signatureonbuch
  \restoregeometry
  \newpage
}

% Rangée « Ce cours contient » (page de garde)
\newcommand{\poptile}[3]{% #1 couleur #2 gros texte #3 légende
  \begin{tcolorbox}[enhanced,colback=#1,colframe=popink,boxrule=2pt,arc=9pt,shadow={2.5pt}{-2.5pt}{0pt}{popink},
      left=6pt,right=6pt,top=9pt,bottom=9pt,halign=center,valign=center,height=1.75cm,width=\linewidth,nobeforeafter]
    {\popdisplay\fontsize{12.5}{13}\selectfont\color{white}\MakeUppercase{#2}}\par\vspace{3pt}
    {\centering\popsemi\fontsize{7.8}{9.5}\selectfont\color{white}#3\par}
  \end{tcolorbox}}
\newcommand{\popcontenu}{%
  \par\noindent{\popxboldls\fontsize{7.5}{7.5}\selectfont\color{popmuted}CE COURS SIGNÉ CONTIENT}\par\vspace{7pt}
  \noindent\begin{minipage}[t]{0.235\linewidth}\poptile{popblue}{Cours}{complet, illustré, pas à pas}\end{minipage}\hfill
  \begin{minipage}[t]{0.235\linewidth}\poptile{poporange}{Méthodes}{et exercices résolus}\end{minipage}\hfill
  \begin{minipage}[t]{0.235\linewidth}\poptile{poppink}{Exercices}{type examen + corrigés}\end{minipage}\hfill
  \begin{minipage}[t]{0.235\linewidth}\poptile{popgreen}{Fiche bilan}{l'essentiel à retenir}\end{minipage}\par}

% Sommaire
\newtcolorbox{sommaire}{enhanced,colback=white,colframe=popink,boxrule=2.2pt,arc=10pt,
  drop shadow=popink,left=18pt,right=18pt,top=13pt,bottom=13pt,before skip=6pt,after skip=20pt,
  before upper={\poppill{Sommaire}{poppurple}{white}\par\vspace{9pt}\popsemi\fontsize{10.6}{14.5}\selectfont\color{popink}}}

% Signature de fin de cours — poussée tout en bas de la dernière page
\newcommand{\findecours}{%
  \par\vfill\nopagebreak
  \begin{center}\popsquiggle{poporange}\end{center}\vspace{4pt}
  \signatureonbuch
}
\AtBeginDocument{\def\llbracket{\mathopen{[\mkern-3.5mu[}}\def\rrbracket{\mathclose{]\mkern-3.5mu]}}} % intervalles d'entiers (glyphe ⟦ absent de la police)
% Glyphes absents de Fira Math : redéfinis à partir de glyphes disponibles
\AtBeginDocument{%
  \def\setminus{\mathbin{\backslash}}\let\smallsetminus\setminus
  \def\vdots{\mathord{\vbox{\baselineskip3.2pt\lineskiplimit0pt\kern4pt\hbox{.}\hbox{.}\hbox{.}}}}%
  \def\ddots{\mathinner{\mkern1mu\raise6.5pt\hbox{.}\mkern2mu\raise3.6pt\hbox{.}\mkern2mu\raise0.7pt\hbox{.}\mkern1mu}}%
  \def\triangle{\mathord{\scalebox{0.9}{$\symup{\Delta}$}}}%
  \def\nabla{\mathord{\raisebox{\depth}{\scalebox{1}[-1]{$\symup{\Delta}$}}}}%
  \def\checkmark{\popcheck{popdark}}%
  \def\star{\mathbin{\ast}}\def\bigstar{\mathord{\ast}}%
  \def\diamond{\mathbin{\scalebox{0.6}{$\Diamond$}}}%
  \def\bigcup{\mathop{\mathchoice{\raisebox{-0.3em}{\scalebox{1.7}{$\cup$}}}{\scalebox{1.25}{$\cup$}}{\cup}{\cup}}\displaylimits}%
  \def\bigcap{\mathop{\mathchoice{\raisebox{-0.3em}{\scalebox{1.7}{$\cap$}}}{\scalebox{1.25}{$\cap$}}{\cap}{\cap}}\displaylimits}%
  \def\lesseqqgtr{\mathrel{\lessgtr}}\def\bigoplus{\mathop{\oplus}\displaylimits}%
}
\providecommand{\ssi}{si et seulement si} % abréviation fréquente
\AtBeginDocument{\providecommand{\wideparen}{\overparen}} % arc de cercle

%% FIGFIX-BEGIN v3 — correctifs globaux des figures (docs/onbuch-generation/tools/figfix)
\usepackage{adjustbox}\usepackage{etoolbox}
% toute figure TikZ/pgfplots plus large que la ligne est réduite à la largeur disponible
\BeforeBeginEnvironment{tikzpicture}{\begin{adjustbox}{max width=\linewidth}}
\AfterEndEnvironment{tikzpicture}{\end{adjustbox}}
% légendes pgfplots lisibles par-dessus les courbes
\pgfplotsset{every axis/.append style={legend style={fill=white,fill opacity=0.92,text opacity=1,draw=black!15,inner sep=3pt}}}
% étiquettes d'arêtes (syntaxe edge["..."]) avec fond blanc
\tikzset{every edge quotes/.append style={fill=white,inner sep=1.2pt}}
% bulles de cartes mentales : texte à la ligne dans la bulle, bulle un peu plus grande
\tikzset{every concept/.append style={text width=2.0cm,align=center,minimum size=2.3cm,inner sep=2pt}}
%% FIGFIX-END
\begin{document}

\pagegarde{Statistiques}{Mathématiques}{1ère A}{Organisation et gestion des données}{NA06}{8 h}{Cette leçon aborde l'organisation de données discrètes en classes d'égales amplitudes, le calcul et l'interprétation de la moyenne, de la médiane, de la variance et de l'écart-type, ainsi que la construction d'histogrammes et de polygones des effectifs cumulés. Elle prépare l'élève à mener de petites enquêtes statistiques concrètes.
\vspace{4pt}
\begin{multicols}{2}\small
\begin{itemize}
\item À la fin de cette leçon, je sais regrouper une population en classes d'égales amplitudes et construire le tableau des effectifs et des fréquences cumulés croissants et décroissants.
\item Je sais calculer la moyenne d'une série à modalités simples et d'une série regroupée en classes.
\item Je sais déterminer la médiane à partir des tableaux des effectifs cumulés et à partir des polygones des effectifs cumulés.
\item Je sais calculer la variance et l'écart-type d'une série statistique.
\item Je sais construire un histogramme et les polygones des effectifs (ou fréquences) cumulés croissants et décroissants.
\item Je sais interpréter la moyenne, la médiane, la variance et l'écart-type dans le contexte d'une enquête réelle.
\end{itemize}\end{multicols}}

\begin{sommaire}
\begin{multicols}{2}
\begin{enumerate}
\item 1. Organisation des données : classes d'égales amplitudes et tableaux de fréquences
\item 2. Caractéristique de position : la moyenne
\item 3. Caractéristique de position : la médiane
\item 4. Caractéristiques de dispersion : variance et écart-type
\item 5. Représentations graphiques : histogramme et polygones des effectifs cumulés
\item 6. Interprétation et communication des résultats d'une enquête
\item 7. Compléments utiles pour la Terminale (optionnel)
\item Activité d'intégration
\item Méthodes et exercices résolus
\item Exercices
\item Corrigés des exercices
\item Fiche bilan
\end{enumerate}
\end{multicols}
\end{sommaire}

\begin{objectifs}
\begin{itemize}
\item À la fin de cette leçon, je sais regrouper une population en classes d'égales amplitudes et construire le tableau des effectifs et des fréquences cumulés croissants et décroissants.
\item Je sais calculer la moyenne d'une série à modalités simples et d'une série regroupée en classes.
\item Je sais déterminer la médiane à partir des tableaux des effectifs cumulés et à partir des polygones des effectifs cumulés.
\item Je sais calculer la variance et l'écart-type d'une série statistique.
\item Je sais construire un histogramme et les polygones des effectifs (ou fréquences) cumulés croissants et décroissants.
\item Je sais interpréter la moyenne, la médiane, la variance et l'écart-type dans le contexte d'une enquête réelle.
\item Je sais choisir la représentation graphique la plus adaptée pour communiquer les résultats d'une étude statistique.
\item Je sais identifier les erreurs fréquentes lors du calcul des caractéristiques de position et de dispersion.
\end{itemize}
\end{objectifs}

\begin{prerequis}
\begin{itemize}
\item Notion de variable statistique discrète et continue (Seconde).
\item Calcul de moyennes arithmétiques simples (Seconde).
\item Lecture et construction de tableaux à double entrée (Seconde).
\item Notions de base sur les graphiques cartésiens (abscisses, ordonnées, échelles).
\item Manipulation des pourcentages et des fréquences (Seconde).
\end{itemize}
\end{prerequis}

\newpage

\coursec{Organisation des données : classes d'égales amplitudes et tableaux de fréquences}

\textbf{Situation d'accroche.} Au marché Mokolo à Yaoundé, Maman Béa vend des plantains. Chaque soir, pendant trente jours, elle note soigneusement dans son cahier le nombre de kilogrammes vendus : $45$, $62$, $38$, $71$, $55$, $49$, $66$, $58$, $42$, $60$, $53$, $47$, $69$, $51$, $64$, $57$, $44$, $73$, $50$, $59$, $63$, $48$, $56$, $70$, $41$, $65$, $52$, $61$, $46$, $54$ (en kg). Trente nombres éparpillés : impossible d'y voir clair à l'œil nu ! Pourtant, Maman Béa a des questions très concrètes : \emph{quelle quantité vend-elle en moyenne par jour ? Quelle vente sépare ses « petits jours » de ses « grands jours » ? Ses ventes sont-elles régulières ou très irrégulières ?} Pour répondre, il faut d'abord \textbf{organiser} les données : les trier, les regrouper en classes, compter les effectifs et calculer des fréquences. C'est l'objet de cette première section, fondation de toute la leçon.

\textbf{Problématique.} Comment transformer une liste brute de données en un tableau clair, exploitable et prêt pour les calculs (moyenne, médiane, écart-type) et les représentations graphiques ?

\courssub{Population, caractère et modalités}

Avant de calculer quoi que ce soit, il faut savoir \emph{de quoi} et \emph{de qui} l'on parle. Le vocabulaire de la statistique est précis ; prenons le temps de le fixer une fois pour toutes.

\begin{definition}[Population, individu, caractère, modalité]
\begin{itemize}
\item Une \cle{population statistique} est l'ensemble des objets (personnes, jours, marchandises, élèves\ldots) sur lesquels porte l'étude. Chaque élément de la population est un \cle{individu} (ou \emph{unité statistique}).
\item Un \cle{caractère} (ou \emph{variable statistique}) est une propriété que l'on observe sur chaque individu.
\item Les différentes valeurs prises par le caractère sont ses \cle{modalités}.
\item Un caractère est \cle{quantitatif discret} s'il ne prend que des valeurs numériques isolées (nombre d'enfants, nombre de sacs vendus) ; il est \cle{quantitatif continu} s'il peut prendre n'importe quelle valeur d'un intervalle (taille, masse, durée).
\end{itemize}
\end{definition}

Dans notre situation : la population est l'ensemble des $30$ jours d'observation ; chaque jour est un individu ; le caractère est la masse de plantains vendue, exprimée en kg ; les modalités observées sont les nombres $38, 41, 42, \ldots, 73$. Bien que la masse soit en théorie un caractère continu, les valeurs relevées ici sont des entiers : on parle de série à caractère discret \emph{à nombreuses modalités}, ce qui justifiera le regroupement en classes.

\begin{attention}[Bien délimiter la population]
Toute conclusion statistique n'est valable que pour la population étudiée. Si Maman Béa a relevé ses ventes en \emph{saison sèche}, elle ne peut pas en déduire ses ventes de saison des pluies. Au Probatoire, précise toujours la population et le caractère dans ta rédaction : c'est un point de méthode apprécié.
\end{attention}

\courssub{Pourquoi regrouper en classes ?}

Lorsque le caractère prend beaucoup de valeurs distinctes (ici $20$ modalités différentes pour $30$ observations), un tableau « valeur par valeur » serait long et peu lisible : presque chaque valeur n'apparaîtrait qu'une ou deux fois. L'idée — naturelle — est de \textbf{regrouper} les valeurs proches dans des « cases » appelées \cle{classes}. C'est exactement ce que fait un commerçant qui range sa monnaie par coupures, ou le service des examens qui annonce : « tant de candidats entre $0$ et $5$, tant entre $5$ et $10$\ldots ».

\begin{definition}[Classe d'égales amplitudes, centre de classe]
Regrouper une série en \cle{classes d'égales amplitudes}, c'est partager l'intervalle $[\min ; \max]$ des valeurs en $k$ intervalles de même longueur $h$, appelée \cle{amplitude} :
\[ [a_0 ; a_1[,\ [a_1 ; a_2[,\ \ldots,\ [a_{k-1} ; a_k], \qquad h = a_{i+1} - a_i \ \text{pour tout } i. \]
Le \cle{centre} de la classe $[a_i ; a_{i+1}[$ est le nombre $c_i = \dfrac{a_i + a_{i+1}}{2}$. C'est lui qui représentera la classe dans les calculs de moyenne et de variance.
\end{definition}

\begin{attention}[Classes disjointes et exhaustives]
Deux règles d'or :
\begin{itemize}
\item les classes ne doivent \textbf{jamais se chevaucher} : on utilise des intervalles \textbf{semi-ouverts} $[a ; b[$, de sorte que la borne $b$ appartient à la classe suivante ;
\item la \textbf{dernière} classe est fermée des deux côtés : $[a_{k-1} ; a_k]$, pour inclure la valeur maximale.
\end{itemize}
Écrire $[40;50]$ puis $[50;60]$ est une erreur classique : où rangerait-on la valeur $50$ ? La convention correcte est $[40;50[$, $[50;60[$, \ldots, et la dernière classe fermée.
\end{attention}

\courssub{Combien de classes ? La règle de Sturges}

Trop peu de classes écrasent l'information ; trop de classes la dispersent. Le statisticien Herbert Sturges a proposé en 1926 une règle simple, encore enseignée partout dans le monde, qui donne un bon compromis pour les échantillons de taille modérée.

\begin{methode}[Choisir le nombre de classes et l'amplitude]
Pour une série de $N$ observations dont le minimum est $m$ et le maximum $M$ :
\begin{enumerate}
\item Calculer le nombre de classes par la \cle{règle de Sturges} : $k \approx 1 + 3{,}3 \log_{10}(N)$, arrondi à l'entier le plus proche.
\item Calculer l'\cle{amplitude} : $h = \dfrac{M - m}{k}$, \textbf{arrondie au supérieur} (à l'unité ou au dixième près selon les données) pour être sûr de couvrir toute l'étendue.
\item Construire les classes en partant d'une borne inférieure $a_0 \leq m$ (souvent $m$ elle-même) : $[a_0 ; a_0+h[$, $[a_0+h ; a_0+2h[$, \ldots
\item Vérifier que la dernière classe contient bien $M$.
\end{enumerate}
\end{methode}

\begin{exemplebox}[Application aux ventes de plantains]
Ici $N = 30$, $m = 38$, $M = 73$.
\begin{align*}
k &\approx 1 + 3{,}3 \log_{10}(30) \approx 1 + 3{,}3 \times 1{,}477 \approx 5{,}87 \quad \Longrightarrow \quad k = 6 \text{ classes}.\\
h &= \frac{73 - 38}{6} = \frac{35}{6} \approx 5{,}83 \quad \Longrightarrow \quad h = 6 \text{ (arrondi au supérieur)}.
\end{align*}
En partant de $38$, les classes sont : $[38;44[$, $[44;50[$, $[50;56[$, $[56;62[$, $[62;68[$, $[68;74]$. La dernière contient bien $73$. Les centres sont $41$, $47$, $53$, $59$, $65$, $71$.
\end{exemplebox}

\begin{savaistu}[Sturges et au-delà]
La règle de Sturges suppose que les données suivent à peu près une répartition « en cloche » et donne de bons résultats pour $N < 200$. Pour de grands échantillons, les statisticiens préfèrent la \emph{règle de Freedman–Diaconis}, qui utilise l'écart interquartile et résiste mieux aux valeurs extrêmes. Tu la rencontreras peut-être en études supérieures — en Première, Sturges suffit amplement et est exigible.
\end{savaistu}

\courssub{Effectifs, fréquences et cumuls}

Une fois les classes posées, on \emph{dépouille} les données : pour chaque classe, on compte combien d'observations tombent dedans. C'est le travail le plus long, et il exige de la méthode (barre chaque donnée traitée dans la liste brute !).

\begin{definition}[Effectif, fréquence, cumuls]
Soit une série de $N$ observations réparties en $k$ classes.
\begin{itemize}
\item L'\cle{effectif} $n_i$ d'une classe est le nombre d'observations appartenant à cette classe.
\item La \cle{fréquence} de la classe est $f_i = \dfrac{n_i}{N}$ (nombre entre $0$ et $1$, souvent exprimé en pourcentage).
\item L'\cle{effectif cumulé croissant} $N_i^{\uparrow}$ est le nombre d'observations \textbf{strictement inférieures à la borne supérieure} de la classe : $N_i^{\uparrow} = n_1 + n_2 + \cdots + n_i$.
\item L'\cle{effectif cumulé décroissant} $N_i^{\downarrow}$ est le nombre d'observations \textbf{supérieures ou égales à la borne inférieure} de la classe : $N_i^{\downarrow} = n_i + n_{i+1} + \cdots + n_k$.
\item On définit de même les \cle{fréquences cumulées} croissantes $F_i^{\uparrow} = \dfrac{N_i^{\uparrow}}{N}$ et décroissantes $F_i^{\downarrow} = \dfrac{N_i^{\downarrow}}{N}$.
\end{itemize}
\end{definition}

\begin{exoresolu}[Tableau complet des ventes de plantains]
\textbf{Énoncé.} Voici les ventes quotidiennes (kg) de Maman Béa sur $30$ jours :
$45, 62, 38, 71, 55, 49, 66, 58, 42, 60, 53, 47, 69, 51, 64, 57, 44, 73, 50, 59, 63, 48, 56, 70, 41, 65, 52, 61, 46, 54$.
Construire le tableau complet : classes, centres, effectifs, fréquences (en \%), effectifs cumulés croissants et décroissants.

\tcblower
\textbf{Solution détaillée.} Les classes ont été déterminées plus haut (Sturges : $k=6$, $h=6$). Dépouillons classe par classe.

\emph{Classe $[38;44[$} : valeurs $38, 42, 41$ $\Rightarrow n_1 = 3$.
\emph{Classe $[44;50[$} : valeurs $45, 49, 47, 44, 48, 46$ $\Rightarrow n_2 = 6$.
\emph{Classe $[50;56[$} : valeurs $55, 53, 51, 50, 52, 54$ $\Rightarrow n_3 = 6$.
\emph{Classe $[56;62[$} : valeurs $58, 60, 57, 59, 56, 61$ $\Rightarrow n_4 = 6$.
\emph{Classe $[62;68[$} : valeurs $62, 66, 64, 63, 65$ $\Rightarrow n_5 = 5$.
\emph{Classe $[68;74]$} : valeurs $71, 69, 73, 70$ $\Rightarrow n_6 = 4$.

\emph{Cumuls croissants} : $3$ ; $3+6=9$ ; $9+6=15$ ; $15+6=21$ ; $21+5=26$ ; $26+4=30$.
\emph{Cumuls décroissants} : $30$ ; $30-3=27$ ; $27-6=21$ ; $21-6=15$ ; $15-6=9$ ; $9-5=4$.

\begin{center}
\begin{tabular}{|c|c|c|c|c|c|}
\hline
\rowcolor{popblueL} Classe & Centre $c_i$ & Effectif $n_i$ & Fréquence $f_i$ (\%) & Eff. cum. croissant & Eff. cum. décroissant \\
\hline
$[38;44[$ & $41$ & $3$ & $10,0$ & $3$ & $30$ \\
\hline
$[44;50[$ & $47$ & $6$ & $20,0$ & $9$ & $27$ \\
\hline
$[50;56[$ & $53$ & $6$ & $20,0$ & $15$ & $21$ \\
\hline
$[56;62[$ & $59$ & $6$ & $20,0$ & $21$ & $15$ \\
\hline
$[62;68[$ & $65$ & $5$ & $16,7$ & $26$ & $9$ \\
\hline
$[68;74]$ & $71$ & $4$ & $13,3$ & $30$ & $4$ \\
\hline
\rowcolor{popgoldL} \textbf{Total} & & $\mathbf{30}$ & $\mathbf{100}$ & & \\
\hline
\end{tabular}
\end{center}

\emph{Lecture} : la fréquence cumulée croissante de la classe $[56;62[$ vaut $\frac{21}{30} = 70\,\%$ : Maman Béa vend \textbf{moins de $62$ kg} dans $70\,\%$ des journées. Symétriquement, l'effectif cumulé décroissant de la classe $[50;56[$ vaut $21$ : il y a $21$ journées où elle vend \textbf{au moins $50$ kg}. Voilà déjà une information précieuse pour l'approvisionnement !
\end{exoresolu}

\begin{popfigure}
\begin{center}
\begin{tikzpicture}[scale=1]
\draw[fill=popblueL, draw=popblue] (0,0) rectangle (12,1);
\node[text=popdark, font=\bfseries] at (6,0.5) {Lecture du tableau des fréquences cumulées};
\draw[fill=popcream, draw=popline] (0,-1.2) rectangle (5.8,0);
\node[text=popdark, align=center, font=\small] at (2.9,-0.6) {Cumul \textbf{croissant} en $62$ : $26$\\ $\Rightarrow$ $26$ jours sur $30$ : ventes $< 62$ kg};
\draw[fill=popgreenL, draw=popgreen] (6.2,-1.2) rectangle (12,0);
\node[text=popdark, align=center, font=\small] at (9.1,-0.6) {Cumul \textbf{décroissant} en $50$ : $21$\\ $\Rightarrow$ $21$ jours sur $30$ : ventes $\geq 50$ kg};
\draw[-{Stealth}, thick, popblue] (2.9,-1.3) -- (2.9,-2.1);
\draw[-{Stealth}, thick, popgreen] (9.1,-1.3) -- (9.1,-2.1);
\draw[thick, popdark] (1,-2.3) -- (11,-2.3);
\foreach \x/\lab in {1/38, 2.67/44, 4.33/50, 6/56, 7.67/62, 9.33/68, 11/74}
  \draw[popdark] (\x,-2.4) -- (\x,-2.2) node[above, font=\scriptsize, text=popdark] {\lab};
\node[font=\scriptsize, text=popmuted] at (6,-2.75) {Masse de plantains vendue (kg)};
\draw[fill=popblue!40, draw=popblue] (1,-2.3) rectangle (7.67,-2.05);
\draw[fill=popgreen!40, draw=popgreen] (4.33,-2.3) rectangle (11,-2.55);
\end{tikzpicture}
\end{center}
\legende{Les deux lectures complémentaires des effectifs cumulés : « moins de $62$ kg » (croissant, $26$ jours) et « au moins $50$ kg » (décroissant, $21$ jours).}
\end{popfigure}

\courssub{Vérifier la cohérence du tableau}

Un tableau de fréquences contient des \emph{tests de cohérence intégrés}. Avant de passer aux calculs, vérifie systématiquement :

\begin{aretenir}[Contrôles obligatoires]
\begin{itemize}
\item $\displaystyle\sum_{i=1}^{k} n_i = N$ : la somme des effectifs égale l'effectif total (ici $3+6+6+6+5+4 = 30$).
\item $\displaystyle\sum_{i=1}^{k} f_i = 1$ (ou $100\,\%$), aux arrondis près (ici $10 + 20 + 20 + 20 + 16{,}7 + 13{,}3 = 100$).
\item Le \textbf{dernier} effectif cumulé croissant vaut $N$ ; la \textbf{dernière} fréquence cumulée croissante vaut $1$ (ou $100\,\%$).
\item Le \textbf{premier} effectif cumulé décroissant vaut $N$.
\item Pour chaque classe : $N_i^{\uparrow} + N_{i+1}^{\downarrow} = N$ (croissant d'une classe + décroissant de la suivante).
\end{itemize}
\end{aretenir}

\begin{attention}[Pièges fréquents au Probatoire]
\begin{itemize}
\item \textbf{Oublier d'arrondir l'amplitude au supérieur} : avec $h = 5{,}83$ non arrondi, la dernière classe s'arrêterait à $38 + 6 \times 5{,}83 \approx 73$, mais les classes « propres » à valeurs entières deviennent impossibles ; arrondir à $6$ garantit la couverture.
\item \textbf{Confondre cumul croissant et décroissant} : le cumul croissant répond à « moins de\ldots », le décroissant à « au moins\ldots ». Lis la question deux fois.
\item \textbf{Arrondis des fréquences} : si la somme des fréquences arrondies donne $99{,}9\,\%$ ou $100{,}1\,\%$, ce n'est pas une erreur de calcul mais un effet d'arrondi ; signale-le par « aux arrondis près ».
\item \textbf{Une observation égale à une borne} : la valeur $44$ appartient à $[44;50[$, pas à $[38;44[$. C'est la convention semi-ouverte qui tranche.
\end{itemize}
\end{attention}

\begin{exercice}[À toi de jouer]
Un cybercafé de Douala relève le nombre de clients par jour pendant $20$ jours :
$12, 25, 8, 31, 19, 22, 15, 28, 10, 24, 17, 33, 21, 14, 26, 9, 30, 18, 23, 16$.
\begin{enumerate}
\item Déterminer le nombre de classes par la règle de Sturges et l'amplitude correspondante.
\item Construire le tableau complet (effectifs, fréquences, cumuls croissants et décroissants).
\item Quel pourcentage de journées compte moins de $26$ clients ?
\end{enumerate}
\end{exercice}

\begin{corrige}[Exercice 1]
\begin{enumerate}
\item $N = 20$, $m = 8$, $M = 33$. Sturges : $k \approx 1 + 3{,}3\log_{10}(20) \approx 1 + 3{,}3 \times 1{,}301 \approx 5{,}29$, donc $k = 5$ classes. Amplitude : $h = \frac{33-8}{5} = 5$. Classes : $[8;13[$, $[13;18[$, $[18;23[$, $[23;28[$, $[28;33]$.
\item Dépouillement : 
$[8;13[$ : $12, 8, 10, 9 \Rightarrow 4$ ; 
$[13;18[$ : $15, 17, 14, 16 \Rightarrow 4$ ; 
$[18;23[$ : $19, 22, 21, 18 \Rightarrow 4$ ; 
$[23;28[$ : $25, 24, 26, 23 \Rightarrow 4$ ; 
$[28;33]$ : $31, 28, 33, 30 \Rightarrow 4$. 
Cumuls croissants : $4, 8, 12, 16, 20$ ; décroissants : $20, 16, 12, 8, 4$ ; fréquences : $20\,\%$ chacune. Vérifications : $4\times 5 = 20$ ; dernière fréquence cumulée croissante $= 100\,\%$.
\item « Moins de $26$ clients » : la borne $26$ se trouve à l'intérieur de la classe $[23;28[$.
    \begin{itemize}
    \item \textbf{À partir du tableau seulement (données groupées)} : on utilise l'interpolation linéaire (méthode vue en section 3). L'effectif cumulé croissant au début de la classe est $12$. On suppose une répartition uniforme dans la classe : $12 + \frac{26-23}{5} \times 4 = 12 + 2,4 = 14,4$ journées, soit environ $72\,\%$.
    \item \textbf{À partir des données brutes (disponibles ici)} : on compte exactement les valeurs $< 26$ : $8, 9, 10, 12, 14, 15, 16, 17, 18, 19, 21, 22, 23, 24, 25$ soit $15$ journées, donc $75\,\%$.
    \end{itemize}
    Au Probatoire, si seules les classes sont données, on applique l'interpolation linéaire.
\end{enumerate}
\end{corrige}

\begin{aretenir}[L'essentiel de la section]
\begin{itemize}
\item Une étude statistique porte sur une \cle{population}, décrite par un \cle{caractère} dont les valeurs sont les \cle{modalités}.
\item Quand les modalités sont nombreuses, on regroupe en \cle{classes d'égales amplitudes} $[a ; b[$ (dernière fermée), de nombre $k \approx 1 + 3{,}3\log_{10} N$ (Sturges) et d'amplitude $h = \frac{M-m}{k}$ arrondie au supérieur.
\item Le tableau complet donne effectifs $n_i$, fréquences $f_i = n_i/N$, cumuls croissants (« moins de ») et décroissants (« au moins »).
\item Contrôles : $\sum n_i = N$, $\sum f_i = 1$, dernier cumul croissant $= N$, premier cumul décroissant $= N$.
\item Ce tableau est la \textbf{matière première} des sections suivantes : moyenne, médiane, variance, histogramme et polygones cumulés s'en déduiront directement.
\end{itemize}
\end{aretenir}

\coursec{Caractéristique de position : la moyenne}

Dans la section précédente, nous avons appris à \emph{organiser} les données : regrouper une série statistique en classes d'égales amplitudes, dresser le tableau des effectifs et des fréquences cumulés. C'est un premier pas, mais un tableau reste une masse d'informations difficile à lire d'un coup d'œil. Le directeur d'une coopérative de planteurs de la plaine de la Mvila ne veut pas contempler 40 lignes de chiffres : il veut savoir, en une phrase, « en moyenne, combien de régimes de plantains chaque producteur a-t-il vendus cette semaine ? ».

C'est exactement le rôle d'une \cle{caractéristique de position} : résumer toute la série par \emph{un seul nombre} qui indique autour de quelle valeur les données se situent. La plus célèbre — et la plus utilisée — est la \cle{moyenne}. Cette section lui est entièrement consacrée : nous la définirons rigoureusement, nous l'étendrons aux séries regroupées en classes, nous justifierons pourquoi le centre de classe joue un rôle central, et nous apprendrons à l'interpréter comme le « centre de gravité » de la distribution.

\courssub{La moyenne d'une série à modalités simples}

\textbf{Intuition.} Imagine que cinq élèves de ta classe ont obtenu les notes $8$, $10$, $12$, $12$ et $18$ sur $20$ au dernier devoir de mathématiques. Si l'on voulait « redistribuer équitablement » le total des points entre les cinq élèves, chacun recevrait
\[
\frac{8+10+12+12+18}{5}=\frac{60}{5}=12.
\]
La moyenne est exactement cela : la valeur que recevrait chaque individu si l'on partageait le total de manière égale. Remarque au passage que la note $12$ apparaît deux fois : plutôt que d'additionner deux fois $12$, on peut écrire $2\times 12$. C'est l'idée de la \emph{pondération par les effectifs}.

\begin{definition}[Moyenne arithmétique d'une série statistique]
Soit une série statistique à caractère discret, de modalités $x_1, x_2, \dots, x_p$, d'effectifs respectifs $n_1, n_2, \dots, n_p$, et d'effectif total $N = n_1+n_2+\dots+n_p$. La \cle{moyenne} de la série, notée $\bar{x}$, est le nombre
\[
\bar{x}=\frac{n_1x_1+n_2x_2+\dots+n_px_p}{n_1+n_2+\dots+n_p}
=\frac{\displaystyle\sum_{i=1}^{p} n_i\,x_i}{\displaystyle\sum_{i=1}^{p} n_i}
=\frac{1}{N}\sum_{i=1}^{p} n_i\,x_i.
\]
On dit que $\bar{x}$ est la \emph{moyenne arithmétique pondérée} des modalités par leurs effectifs.
\end{definition}

\begin{aretenir}[Moyenne et fréquences]
En divisant numérateur et dénominateur par $N$, on obtient une écriture équivalente avec les fréquences $f_i=\dfrac{n_i}{N}$ :
\[
\bar{x}=\sum_{i=1}^{p} f_i\,x_i.
\]
Cette formule est très pratique lorsque l'énoncé donne directement les fréquences (ou les pourcentages) plutôt que les effectifs.
\end{aretenir}

\begin{exemplebox}[Moyenne des notes d'un groupe]
Un groupe de $10$ élèves de Première A a obtenu les notes suivantes en français : $6$ ; $8$ ; $8$ ; $10$ ; $11$ ; $12$ ; $12$ ; $12$ ; $14$ ; $17$. Le tableau des effectifs est :
\begin{center}
\begin{tabular}{|c|ccccccc|}
\hline
\rowcolor{popblueL} Note $x_i$ & $6$ & $8$ & $10$ & $11$ & $12$ & $14$ & $17$ \\
\hline Effectif $n_i$ & $1$ & $2$ & $1$ & $1$ & $3$ & $1$ & $1$ \\
\hline
\end{tabular}
\end{center}
La moyenne vaut :
\[
\bar{x}=\frac{1\times 6+2\times 8+1\times 10+1\times 11+3\times 12+1\times 14+1\times 17}{10}
=\frac{110}{10}=11.
\]
En moyenne, les élèves du groupe ont obtenu $11$ sur $20$.
\end{exemplebox}

\courssub{Une propriété fondamentale : la somme des écarts est nulle}

Voici une propriété simple en apparence, mais qui éclaire profondément ce qu'est la moyenne : si l'on mesure, pour chaque donnée, son \emph{écart} à la moyenne (c'est-à-dire $x_i-\bar{x}$), et que l'on additionne tous ces écarts en tenant compte des effectifs, on obtient toujours zéro.

\begin{propriete}[Somme nulle des écarts à la moyenne]
Pour toute série statistique de moyenne $\bar{x}$, on a :
\[
\sum_{i=1}^{p} n_i\,(x_i-\bar{x})=0.
\]
Autrement dit, la somme des écarts positifs compense exactement la somme des écarts négatifs.
\end{propriete}

\begin{experience}[Démonstration guidée]
Développons la somme en séparant les deux termes :
\begin{align*}
\sum_{i=1}^{p} n_i(x_i-\bar{x})
&=\sum_{i=1}^{p} n_i x_i-\sum_{i=1}^{p} n_i\,\bar{x} \\
&=\sum_{i=1}^{p} n_i x_i-\bar{x}\sum_{i=1}^{p} n_i \\
&=N\bar{x}-\bar{x}\,N \qquad \text{car } \sum_{i=1}^{p}n_i x_i = N\bar{x} \text{ et } \sum_{i=1}^{p}n_i=N\\
&=0.
\end{align*}
La démonstration tient en trois lignes : elle repose uniquement sur la définition $\bar{x}=\dfrac{1}{N}\sum n_i x_i$. Retiens l'idée clé : $\bar{x}$ est une constante, donc elle peut sortir de la somme.
\end{experience}

Cette propriété justifie l'interprétation physique de la moyenne : c'est le \cle{centre de gravité} de la distribution. Si l'on place sur une tige des poids égaux aux positions $x_i$ (chaque position répétée $n_i$ fois), la tige est en équilibre exactement au point d'abscisse $\bar{x}$ : les « moments » à gauche et à droite se compensent. C'est pourquoi une seule valeur extrême peut « tirer » la moyenne vers elle, comme un poids lourd éloigné du point d'équilibre.

\courssub{Extension aux séries regroupées en classes}

\textbf{Le problème.} Revenons à notre enquête de la section 1 : la coopérative a relevé le nombre de régimes de plantains vendus par $40$ producteurs pendant une semaine, et nous avons regroupé les données en classes d'amplitude $10$ :
\begin{center}
\begin{tabular}{|c|ccccc|}
\hline
\rowcolor{popblueL} Classe (nombre de régimes) & $[0;10[$ & $[10;20[$ & $[20;30[$ & $[30;40[$ & $[40;50[$ \\
\hline Effectif $n_i$ & $4$ & $9$ & $14$ & $8$ & $5$ \\
\hline
\end{tabular}
\end{center}
Une fois le regroupement effectué, les valeurs individuelles sont « perdues » : nous savons que $9$ producteurs ont vendu entre $10$ et $19$ régimes, mais nous ne connaissons plus les ventes exactes de chacun. Comment alors calculer une moyenne ?

\textbf{L'idée.} On choisit, pour chaque classe, une \emph{valeur représentative} unique, et l'on fait comme si tous les individus de la classe avaient cette valeur. Le choix naturel est le \cle{centre de classe}.

\begin{definition}[Centre de classe et moyenne d'une série groupée]
Le \cle{centre} de la classe $[a_i\,;\,b_i[$ est le nombre
\[
c_i=\frac{a_i+b_i}{2}.
\]
La moyenne d'une série regroupée en classes est alors approchée par :
\[
\bar{x}\approx \frac{1}{N}\sum_{i=1}^{p} n_i\,c_i.
\]
\end{definition}

\begin{experience}[Pourquoi le centre de classe donne-t-il une bonne approximation ?]
Raisonnons sur une classe $[a\,;\,b[$ contenant $n$ individus, dont les valeurs exactes (inconnues) sont $v_1, v_2, \dots, v_n$. La contribution exacte de cette classe à la somme totale serait $v_1+v_2+\dots+v_n$.

Faisons l'\emph{hypothèse de répartition uniforme} : les valeurs sont régulièrement dispersées dans la classe, c'est-à-dire qu'il y en a « autant » dans la moitié basse $[a\,;\,c[$ que dans la moitié haute $[c\,;\,b]$, où $c=\dfrac{a+b}{2}$. Alors, pour chaque valeur située à une distance $d$ au-dessus du centre, il y a en moyenne une valeur située à peu près à la distance $d$ en dessous :
\[
(c+d)+(c-d)=2c.
\]
Les écarts se compensent deux à deux, et la somme des valeurs de la classe vaut approximativement $n\times c$ :
\[
v_1+v_2+\dots+v_n \approx n\,c.
\]
En sommant sur toutes les classes, la somme totale des données vaut approximativement $\sum n_i c_i$, d'où $\bar{x}\approx \dfrac{1}{N}\sum n_i c_i$.

Conclusion : remplacer chaque classe par son centre revient à supposer que, \emph{à l'intérieur de chaque classe, la moyenne des valeurs est le centre}. C'est une approximation d'autant meilleure que les classes sont étroites et que les données s'y répartissent de façon équilibrée.
\end{experience}

\begin{attention}[Une moyenne approchée, pas exacte]
Pour une série regroupée en classes, la valeur obtenue $\bar{x}=\dfrac{1}{N}\sum n_i c_i$ est une \emph{approximation} de la vraie moyenne : elle dépend du découpage en classes choisi. Au Probatoire, on écrira proprement « la moyenne est environ égale à … » et l'on utilisera le symbole $\approx$ dans la conclusion, même si les calculs intermédiaires se mènent avec le signe $=$.
\end{attention}

\courssub{Méthode de calcul à partir du tableau des effectifs}

\begin{methode}[Calculer la moyenne d'une série groupée en classes]
\begin{enumerate}
\item \textbf{Compléter le tableau} par une ligne (ou une colonne) « centre de classe $c_i$ » : pour chaque classe $[a_i\,;\,b_i[$, calculer $c_i=\dfrac{a_i+b_i}{2}$.
\item \textbf{Ajouter une ligne « produit $n_i c_i$ »} : multiplier chaque centre par l'effectif de sa classe.
\item \textbf{Additionner} les produits : $S=\sum n_i c_i$, et vérifier l'effectif total $N=\sum n_i$.
\item \textbf{Diviser} : $\bar{x}=\dfrac{S}{N}$, puis conclure par une phrase d'interprétation dans le contexte de l'énoncé.
\end{enumerate}
\end{methode}

\begin{exoresolu}[Moyenne des ventes de plantains]
\textbf{Énoncé.} La coopérative de la Mvila a relevé les ventes hebdomadaires (en régimes de plantains) de ses $40$ producteurs, regroupées dans le tableau ci-dessous. Calculer la moyenne de cette série et interpréter le résultat.
\begin{center}
\begin{tabular}{|c|ccccc|}
\hline
\rowcolor{popblueL} Classe & $[0;10[$ & $[10;20[$ & $[20;30[$ & $[30;40[$ & $[40;50[$ \\
\hline Effectif $n_i$ & $4$ & $9$ & $14$ & $8$ & $5$ \\
\hline
\end{tabular}
\end{center}
\tcblower
\textbf{Solution.} On complète le tableau avec les centres de classe et les produits $n_ic_i$ :
\begin{center}
\begin{tabular}{|c|ccccc|c|}
\hline
\rowcolor{popblueL} Classe & $[0;10[$ & $[10;20[$ & $[20;30[$ & $[30;40[$ & $[40;50[$ & Total \\
\hline Centre $c_i$ & $5$ & $15$ & $25$ & $35$ & $45$ & \\
\hline Effectif $n_i$ & $4$ & $9$ & $14$ & $8$ & $5$ & $40$ \\
\hline $n_i c_i$ & $20$ & $135$ & $350$ & $280$ & $225$ & $1010$ \\
\hline
\end{tabular}
\end{center}
Détail des produits : $4\times 5=20$ ; $9\times 15=135$ ; $14\times 25=350$ ; $8\times 35=280$ ; $5\times 45=225$. La somme vaut
\[
S=20+135+350+280+225=1010.
\]
On en déduit :
\[
\bar{x}=\frac{1010}{40}=25{,}25.
\]
\textbf{Interprétation.} En moyenne, un producteur de la coopérative a vendu environ $25$ régimes de plantains cette semaine-là. Si l'on redistribuait équitablement la totalité des $1010$ régimes vendus entre les $40$ producteurs, chacun en recevrait $25{,}25$.
\end{exoresolu}

La figure suivante visualise cette distribution : les barres représentent les effectifs de chaque classe, et la ligne verticale marque la position de la moyenne $\bar{x}=25{,}25$. On observe qu'elle se situe au cœur de la classe la plus peuplée, $[20;30[$, légèrement « tirée » vers la droite par les producteurs des classes élevées.

\begin{popfigure}
\begin{center}
\begin{tikzpicture}
\begin{axis}[
    width=0.85\linewidth, height=6.5cm,
    ybar, bar width=18,
    ymin=0, ymax=17,
    xmin=-2, xmax=52,
    xlabel={Nombre de régimes vendus},
    ylabel={Effectif},
    xtick={5,15,25,35,45},
    xticklabels={$[0;10[$,$[10;20[$,$[20;30[$,$[30;40[$,$[40;50[$},
    x tick label style={font=\small, rotate=20, anchor=east},
    axis lines=left,
    every axis plot/.append style={fill=popblue!55, draw=popblue},
]
\addplot coordinates {(5,4) (15,9) (25,14) (35,8) (45,5)};
\addplot[poporange, very thick, dashed] coordinates {(25.25,0) (25.25,16)};
\node[poporange, font=\small, anchor=west] at (axis cs:26,15.2) {$\bar{x}=25{,}25$};
\end{axis}
\end{tikzpicture}
\end{center}
\legende{Diagramme en barres des effectifs et position de la moyenne $\bar{x}=25{,}25$ : la moyenne est le « centre de gravité » de la distribution.}
\end{popfigure}

\begin{attention}[Les deux erreurs classiques au Probatoire]
\begin{itemize}
\item \textbf{Oublier la pondération.} Calculer $\dfrac{5+15+25+35+45}{5}=25$, c'est-à-dire la moyenne des centres \emph{sans tenir compte des effectifs}, est faux : la classe $[20;30[$ compte $14$ producteurs, elle doit « peser » plus lourd que la classe $[40;50[$ qui n'en compte que $5$. Ici, la moyenne correcte est $25{,}25$ et non $25$.
\item \textbf{Confondre série brute et série groupée.} Si l'énoncé donne les valeurs exactes (série brute), on calcule la moyenne exacte avec les modalités $x_i$ ; si l'énoncé donne des classes, on utilise les centres $c_i$ et l'on obtient une valeur \emph{approchée}. Ne mélange jamais les deux démarches dans un même calcul.
\end{itemize}
\end{attention}

\begin{exercice}[À toi de jouer]
Une boutique de recharge MTN et Orange a noté le montant (en FCFA) de $50$ transactions un samedi :
\begin{center}
\begin{tabular}{|c|cccc|}
\hline
\rowcolor{popblueL} Montant (FCFA) & $[0;1000[$ & $[1000;2000[$ & $[2000;3000[$ & $[3000;4000[$ \\
\hline Effectif & $12$ & $18$ & $14$ & $6$ \\
\hline
\end{tabular}
\end{center}
\begin{enumerate}
\item Compléter le tableau avec les centres de classe et les produits $n_ic_i$.
\item Calculer le montant moyen d'une transaction.
\item Le gérant affirme : « la moitié des clients a dépensé plus que la moyenne ». Cette affirmation est-elle nécessairement vraie ? Justifier.
\end{enumerate}
\end{exercice}

\begin{corrige}[Exercice : montant moyen des transactions]
\begin{enumerate}
\item Les centres sont $c_1=500$, $c_2=1500$, $c_3=2500$, $c_4=3500$. Les produits : $12\times 500=6000$ ; $18\times 1500=27000$ ; $14\times 2500=35000$ ; $6\times 3500=21000$.
\item Somme : $S=6000+27000+35000+21000=89000$. Moyenne :
\[
\bar{x}=\frac{89000}{50}=1780 \text{ FCFA}.
\]
\item Non. La moyenne ne partage pas la population en deux moitiés égales : c'est une caractéristique de \emph{répartition du total}, pas de \emph{partage de l'effectif}. Ici, $12+18=30$ clients sur $50$ ont dépensé moins de $2000$ FCFA, donc probablement plus de la moitié des clients a dépensé \emph{moins} que la moyenne de $1780$ FCFA. Le paramètre qui partage l'effectif en deux est la \emph{médiane}, étudiée dans la section suivante.
\end{enumerate}
\end{corrige}

\begin{savaistu}[D'autres moyennes existent]
La moyenne étudiée ici est la moyenne \emph{arithmétique}. En Terminale, tu rencontreras deux autres moyennes, adaptées aux données de nature \emph{multiplicative} :
\begin{itemize}
\item la \cle{moyenne géométrique} $\sqrt[n]{x_1 x_2 \dots x_n}$, idéale pour les taux d'évolution successifs (par exemple, si le prix du cacao augmente de $10\,\%$ puis de $21\,\%$, le taux moyen annuel n'est \emph{pas} $15{,}5\,\%$ mais se calcule avec la moyenne géométrique) ;
\item la \cle{moyenne harmonique} $\dfrac{n}{\frac{1}{x_1}+\dots+\frac{1}{x_n}}$, utilisée pour les moyennes de vitesses ou de débits.
\end{itemize}
Ces trois moyennes vérifient toujours : moyenne harmonique $\leq$ moyenne géométrique $\leq$ moyenne arithmétique. Un beau résultat que tu démontreras peut-être un jour !
\end{savaistu}

\begin{aretenir}[L'essentiel sur la moyenne]
\begin{itemize}
\item Série à modalités simples : $\bar{x}=\dfrac{1}{N}\sum n_i x_i=\sum f_i x_i$.
\item Série groupée en classes : on remplace chaque classe par son centre $c_i=\dfrac{a_i+b_i}{2}$ et $\bar{x}\approx\dfrac{1}{N}\sum n_i c_i$ (valeur \emph{approchée}, justifiée par l'hypothèse de répartition uniforme).
\item La somme des écarts à la moyenne est nulle : $\sum n_i(x_i-\bar{x})=0$ ; la moyenne est le centre de gravité de la distribution.
\item La moyenne est sensible aux valeurs extrêmes : elle ne résume pas toujours bien une série très étalée ou asymétrique — d'où l'intérêt de la médiane (section 3) et des caractéristiques de dispersion (section 4).
\end{itemize}
\end{aretenir}

\coursec{Caractéristique de position : la médiane}

Dans la section précédente, nous avons appris à calculer la \cle{moyenne} d'une série statistique : elle résume la série par une valeur « centrale » obtenue en répartissant équitablement le total des valeurs. Mais la moyenne possède un défaut bien connu : elle est très sensible aux valeurs extrêmes. Imaginons une classe de Première A où neuf élèves ont entre 8 et 12 sur 20 en mathématiques, et où un dixième élève, excellent, a 19. La moyenne de la classe sera tirée vers le haut par cette seule note, et ne reflétera plus vraiment le niveau « typique » du groupe. Pour décrire le centre d'une série autrement, les statisticiens utilisent un deuxième paramètre de position : la \cle{médiane}. C'est l'objet de cette section.

\courssub{Intuition et définition de la médiane}

\textbf{Une situation concrète.}\par

Le marché central de Douala accueille chaque samedi des dizaines de vendeuses de plantains. Une enquête menée par des élèves a relevé le chiffre d'affaires (en milliers de francs CFA) de 60 vendeuses. Si l'on range ces 60 valeurs de la plus petite à la plus grande, on peut chercher la valeur qui \emph{coupe la population en deux} : la moitié des vendeuses gagne moins que cette valeur, l'autre moitié gagne plus. Cette valeur « du milieu » est la médiane. Contrairement à la moyenne, une vendeuse qui aurait réalisé des ventes exceptionnelles ne la déplace pas : seule compte la \emph{position} des valeurs, pas leur grandeur.

\begin{definition}[Médiane d'une série statistique]
Soit une série statistique d'effectif total $N$, dont les valeurs sont rangées dans l'ordre croissant. On appelle \cle{médiane}, notée $Me$, toute valeur qui partage la population en deux parties d'effectifs égaux :
\begin{itemize}
\item au moins $50\,\%$ des valeurs sont inférieures ou égales à $Me$ ;
\item au moins $50\,\%$ des valeurs sont supérieures ou égales à $Me$.
\end{itemize}
Pour une série discrète non groupée, si $N$ est impair, la médiane est la valeur de rang $\dfrac{N+1}{2}$ ; si $N$ est pair, on prend par convention le centre de l'intervalle formé par les valeurs de rangs $\dfrac{N}{2}$ et $\dfrac{N}{2}+1$ (soit la moyenne arithmétique de ces deux valeurs).
\end{definition}

\begin{exemplebox}[Médiane d'une petite série discrète]
Les notes de 7 élèves à un devoir sont, rangées dans l'ordre croissant :
\[ 5 \;;\; 7 \;;\; 9 \;;\; 10 \;;\; 11 \;;\; 13 \;;\; 18. \]
Ici $N = 7$ est impair : la médiane est la valeur de rang $\dfrac{7+1}{2} = 4$, soit $Me = 10$. Trois notes sont en dessous, trois au-dessus. Remarquons que la moyenne vaut $\dfrac{73}{7} \approx 10,4$ : la note 18 l'a légèrement tirée vers le haut, alors que la médiane reste à 10.
\end{exemplebox}

\courssub{Classe médiane et effectifs cumulés croissants}

Lorsque les données sont regroupées en classes d'égales amplitudes, on ne connaît plus les valeurs individuelles : on ne peut donc pas « prendre la valeur du milieu » directement. On procède en deux temps : on repère d'abord la \cle{classe médiane}, puis on estime la médiane à l'intérieur de cette classe.

\begin{definition}[Classe médiane, effectif cumulé antérieur]
Soit une série regroupée en classes d'effectif total $N$. On calcule les \cle{effectifs cumulés croissants} (ECC) : pour chaque classe, l'ECC est la somme des effectifs de cette classe et de toutes les classes précédentes.
\begin{itemize}
\item La \cle{classe médiane} est la première classe dont l'effectif cumulé croissant est supérieur ou égal à $\dfrac{N}{2}$.
\item L'\cle{effectif cumulé antérieur}, noté $F_{\text{prec}}$, est l'effectif cumulé croissant de la classe qui précède la classe médiane (il vaut $0$ si la classe médiane est la première).
\end{itemize}
\end{definition}

\begin{methode}[Déterminer la médiane à partir du tableau]
\begin{enumerate}
\item Calculer l'effectif total $N$ puis le \emph{rang médian} $\dfrac{N}{2}$.
\item Construire la colonne des effectifs cumulés croissants.
\item Repérer la \cle{classe médiane} : la première classe dont l'ECC est supérieur ou égal à $\dfrac{N}{2}$.
\item Noter $L$ la borne inférieure de la classe médiane, $h$ son amplitude, $f_{\text{med}}$ son effectif et $F_{\text{prec}}$ l'effectif cumulé antérieur.
\item Appliquer la formule d'interpolation linéaire :
\[ Me = L + \frac{\dfrac{N}{2} - F_{\text{prec}}}{f_{\text{med}}} \times h. \]
\end{enumerate}
\end{methode}

\courssub{La formule d'interpolation linéaire : d'où vient-elle ?}

La formule précédente n'est pas une recette magique : elle repose sur une hypothèse simple et naturelle, la \emph{répartition uniforme} des valeurs à l'intérieur de la classe médiane. La démonstration guidée ci-dessous est formatrice : elle t'aide à retenir la formule en la comprenant.

\begin{experience}[Démonstration guidée de la formule d'interpolation]
On considère la classe médiane $[L \,;\, L+h[$, d'effectif $f_{\text{med}}$, et on note $F_{\text{prec}}$ l'effectif cumulé antérieur.
\begin{enumerate}
\item \textbf{Cumul avant la classe.} Avant d'entrer dans la classe médiane, $F_{\text{prec}}$ individus ont déjà été comptés. Pour atteindre le rang médian $\dfrac{N}{2}$, il manque donc
\[ \frac{N}{2} - F_{\text{prec}} \quad \text{individus}. \]
\item \textbf{Hypothèse de répartition uniforme.} On suppose que les $f_{\text{med}}$ individus de la classe médiane sont répartis régulièrement sur l'intervalle $[L \,;\, L+h[$. Ainsi, la proportion d'individus de la classe situés avant une valeur $x$ de la classe est proportionnelle à la longueur $x - L$ :
\[ \frac{\text{nombre d'individus avant } x}{f_{\text{med}}} = \frac{x - L}{h}. \]
\item \textbf{Application à la médiane.} La médiane est la valeur $x = Me$ pour laquelle le nombre d'individus avant $x$, dans la classe, est exactement celui qui manque, à savoir $\dfrac{N}{2} - F_{\text{prec}}$. On obtient :
\[ \frac{\dfrac{N}{2} - F_{\text{prec}}}{f_{\text{med}}} = \frac{Me - L}{h}. \]
\item \textbf{Conclusion.} En multipliant les deux membres par $h$ puis en ajoutant $L$ :
\[ Me = L + \frac{\dfrac{N}{2} - F_{\text{prec}}}{f_{\text{med}}} \times h. \qquad \square \]
\end{enumerate}
Autrement dit : la médiane avance dans la classe médiane d'une \emph{fraction de l'amplitude} égale à la \emph{fraction de l'effectif} qu'il faut parcourir pour atteindre le rang médian.
\end{experience}

\begin{attention}[Quand la médiane tombe sur une borne de classe]
Si $\dfrac{N}{2}$ est exactement égal à un effectif cumulé croissant (cas fréquent quand $N$ est pair), la formule donne $Me = L + h$, c'est-à-dire la \emph{borne supérieure} de la classe médiane, qui est aussi la borne inférieure de la classe suivante. Ce n'est pas une erreur de calcul ! La médiane est alors la borne commune des deux classes : interprète le résultat en disant que la moitié de la population se situe en dessous de cette borne. Vérifie toujours la cohérence de ta réponse avec le tableau des ECC.
\end{attention}

\courssub{Exemple chiffré complet : les ventes de plantains}

Reprenons l'enquête auprès des 60 vendeuses de plantains du marché central. Les chiffres d'affaires hebdomadaires (en milliers de FCFA) ont été regroupés en classes d'égales amplitudes :

\begin{center}
\begin{tabularx}{\linewidth}{|l|X|X|X|X|X|}
\hline
\rowcolor{popblueL} Classe (milliers FCFA) & $[10;20[$ & $[20;30[$ & $[30;40[$ & $[40;50[$ & $[50;60[$ \\
\hline
Effectif $n_i$ & 8 & 15 & 20 & 12 & 5 \\
\hline
ECC & 8 & 23 & 43 & 55 & 60 \\
\hline
\end{tabularx}
\end{center}

\begin{exoresolu}[Médiane des ventes de plantains]
\textbf{Énoncé.} Déterminer la médiane de cette série par le calcul, puis interpréter le résultat.
\tcblower
\textbf{Solution détaillée.}
\begin{enumerate}
\item Effectif total : $N = 8+15+20+12+5 = 60$. Rang médian : $\dfrac{N}{2} = 30$.
\item Dans la colonne des ECC, on cherche le premier cumul supérieur ou égal à $30$ : c'est $43$, correspondant à la classe $[30\,;\,40[$. C'est la \cle{classe médiane}.
\item On relève : $L = 30$, $h = 10$, $f_{\text{med}} = 20$, et l'effectif cumulé antérieur $F_{\text{prec}} = 23$.
\item Interpolation linéaire :
\begin{align*}
Me &= L + \frac{\dfrac{N}{2} - F_{\text{prec}}}{f_{\text{med}}} \times h \\
&= 30 + \frac{30 - 23}{20} \times 10 \\
&= 30 + \frac{7}{20} \times 10 = 30 + 3{,}5 = 33{,}5.
\end{align*}
\end{enumerate}
\textbf{Interprétation.} $Me = 33{,}5$ milliers de FCFA, soit \num{33500} FCFA : au moins la moitié des vendeuses réalisent un chiffre d'affaires hebdomadaire inférieur à \num{33500} FCFA, et au moins la moitié réalisent plus. À titre de comparaison, la moyenne de cette série (calculée avec les centres de classes) vaut $33{,}5$ milliers de FCFA : les deux paramètres sont égaux, signe d'une distribution assez équilibrée.
\end{exoresolu}

\courssub{Lecture graphique de la médiane sur le polygone des effectifs cumulés}

Le \cle{polygone des effectifs cumulés croissants} s'obtient en plaçant, pour chaque classe, le point dont l'abscisse est la \emph{borne supérieure} de la classe et l'ordonnée l'effectif cumulé croissant, puis en reliant ces points par des segments (on part de la borne inférieure de la première classe, d'ordonnée $0$). Puisque l'ECC en $x$ représente le nombre d'individus dont la valeur est inférieure à $x$, la médiane est l'\emph{abscisse du point du polygone dont l'ordonnée est} $\dfrac{N}{2}$.

\begin{methode}[Lecture graphique de la médiane]
\begin{enumerate}
\item Tracer le polygone des ECC.
\item Tracer la droite horizontale d'équation $y = \dfrac{N}{2}$.
\item Repérer son point d'intersection avec le polygone.
\item Projeter ce point sur l'axe des abscisses : la valeur lue est la médiane.
\end{enumerate}
Graphiquement, l'interpolation linéaire revient exactement à lire l'abscisse sur le segment qui traverse la ligne $y = \dfrac{N}{2}$ : calcul et graphique sont deux visages de la même méthode.
\end{methode}

\begin{popfigure}
\begin{center}
\begin{tikzpicture}
\begin{axis}[
    width=0.95\linewidth, height=7.5cm,
    xlabel={Chiffre d'affaires (milliers de FCFA)},
    ylabel={Effectifs cumul\'es croissants},
    xmin=5, xmax=65, ymin=0, ymax=70,
    xtick={10,20,30,40,50,60},
    ytick={0,10,20,30,40,50,60},
    grid=both, grid style={popline!40},
    axis lines=left,
    clip=false
]
% Polygone des ECC
\addplot[popcurve, mark=*, mark options={fill=popblue}, mark size=1.8pt]
coordinates {(10,0) (20,8) (30,23) (40,43) (50,55) (60,60)};
% Ligne horizontale y = N/2 = 30
\addplot[poporange, dashed, thick, domain=5:65] {30};
% Projection verticale en x = 33.5
\draw[poporange, dashed, thick] (axis cs:33.5,0) -- (axis cs:33.5,30);
% Point d'intersection
\addplot[only marks, mark=*, mark options={fill=poporange}, mark size=2.5pt] coordinates {(33.5,30)};
\node[poporange, above] at (axis cs:33.5,31) {\footnotesize $\left(33{,}5 \,;\, 30\right)$};
\node[poporange, below] at (axis cs:33.5,-2) {\footnotesize $Me = 33{,}5$};
\node[popmuted, right] at (axis cs:60,30) {\footnotesize $y = \dfrac{N}{2} = 30$};
\end{axis}
\end{tikzpicture}
\end{center}
\legende{Polygone des effectifs cumulés croissants des ventes de plantains. La médiane est l'abscisse du point d'intersection du polygone avec la droite horizontale d'ordonnée $\frac{N}{2} = 30$ : on lit $Me = 33{,}5$ milliers de FCFA.}
\end{popfigure}

Sur la figure, le segment qui joint les points $(30\,;\,23)$ et $(40\,;\,43)$ traverse la ligne $y = 30$ exactement en $x = 33{,}5$ : la lecture graphique confirme le calcul par interpolation. C'est un excellent réflexe au Probatoire : \emph{vérifier ton calcul sur ton graphique}.

\courssub{Moyenne ou médiane : lequel choisir ?}

Les deux paramètres décrivent le « centre » d'une série, mais ils ne réagissent pas de la même manière aux valeurs extrêmes.

\begin{propriete}[Propriété de minimisation de la médiane (admise)]
La médiane $Me$ est une valeur qui rend minimale la somme des écarts absolus :
\[ \sum_{i} n_i \, |x_i - m|. \]
Autrement dit, parmi toutes les valeurs $m$ candidates, la médiane est celle qui est, en moyenne, la « moins éloignée » des valeurs de la série au sens de la distance absolue. (On admet ce résultat ; rappelons de manière analogue que la moyenne minimise la somme des carrés des écarts.)
\end{propriete}

Cette propriété explique la \emph{robustesse} de la médiane : déplacer une valeur extrême très loin ne change pas son rang, donc ne change pas la médiane, alors que la moyenne, elle, est entraînée.

\begin{exemplebox}[Sensibilité aux valeurs extrêmes]
Les revenus mensuels (en milliers de FCFA) de 5 employés d'une petite entreprise de Bafoussam sont : $80 \;;\; 90 \;;\; 100 \;;\; 110 \;;\; 120$.
\begin{itemize}
\item Moyenne : $\bar{x} = 100$ ; médiane : $Me = 100$. Les deux coïncident.
\item Le directeur s'attribue un revenu de $900$ au lieu de $120$ : la série devient $80 \;;\; 90 \;;\; 100 \;;\; 110 \;;\; 900$. La moyenne grimpe à $\bar{x} = 256$, qui ne représente plus personne, tandis que la médiane reste $Me = 100$, fidèle au revenu « typique ».
\end{itemize}
C'est pourquoi les instituts de statistique publient souvent le \emph{salaire médian} plutôt que le salaire moyen.
\end{exemplebox}

\begin{aretenir}[L'essentiel sur la médiane]
\begin{itemize}
\item La médiane partage la population en deux effectifs égaux : $Me$ est la valeur de rang $\dfrac{N}{2}$.
\item Série groupée : la classe médiane est la première classe dont l'ECC atteint ou dépasse $\dfrac{N}{2}$, puis
\[ Me = L + \frac{\dfrac{N}{2} - F_{\text{prec}}}{f_{\text{med}}} \times h. \]
\item Graphiquement, $Me$ est l'abscisse du point du polygone des ECC d'ordonnée $\dfrac{N}{2}$.
\item La médiane est peu sensible aux valeurs extrêmes ; la moyenne l'est beaucoup.
\end{itemize}
\end{aretenir}

\begin{attention}[Pièges fréquents au Probatoire]
\begin{itemize}
\item Ne confonds pas l'\emph{effectif} de la classe médiane $f_{\text{med}}$ avec l'\emph{effectif cumulé} : la formule utilise les deux, mais à des places différentes.
\item Le rang médian est $\dfrac{N}{2}$, et non $\dfrac{N+1}{2}$, pour une série groupée en classes.
\item Pour le polygone des ECC, on place les points aux \emph{bornes supérieures} des classes, jamais aux centres.
\item Donne toujours la médiane avec son unité et une phrase d'interprétation : un résultat nu sans commentaire coûte des points.
\end{itemize}
\end{attention}

\begin{savaistu}[Quartiles et déciles : vers la Terminale]
La médiane coupe la population en deux. Rien n'empêche de la couper en quatre : les \cle{quartiles} $Q_1$, $Q_2 = Me$ et $Q_3$ sont les valeurs correspondant aux rangs $\dfrac{N}{4}$, $\dfrac{N}{2}$ et $\dfrac{3N}{4}$. On peut même couper en dix avec les \cle{déciles}. En Terminale, tu utilisera ces paramètres pour construire des \emph{diagrammes en boîte} et comparer finement des distributions — par exemple les notes de deux classes de Première au même devoir, ou les rendements de deux variétés de maïs dans la plaine des Mbo. La méthode reste la même : effectifs cumulés et interpolation linéaire.
\end{savaistu}

\begin{exercice}[À toi de jouer]
Une coopérative de la région de l'Ouest a relevé les masses (en kg) de sacs de haricots livrés par 50 producteurs :
\begin{center}
\begin{tabularx}{\linewidth}{|l|X|X|X|X|X|}
\hline
\rowcolor{popblueL} Masse (kg) & $[40;50[$ & $[50;60[$ & $[60;70[$ & $[70;80[$ & $[80;90[$ \\
\hline
Effectif & 6 & 12 & 18 & 9 & 5 \\
\hline
\end{tabularx}
\end{center}
\begin{enumerate}
\item Construire la colonne des effectifs cumulés croissants.
\item Déterminer la classe médiane puis calculer la médiane par interpolation linéaire.
\item Tracer le polygone des ECC et vérifier graphiquement le résultat.
\item La moyenne de la série vaut $64$ kg. Comparer avec la médiane et commenter.
\end{enumerate}
\end{exercice}

\begin{corrige}[Exercice : masses des sacs de haricots]
\begin{enumerate}
\item Les ECC sont : $6 \;;\; 18 \;;\; 36 \;;\; 45 \;;\; 50$.
\item $N = 50$, donc $\dfrac{N}{2} = 25$. Le premier ECC supérieur ou égal à $25$ est $36$ : la classe médiane est $[60\,;\,70[$. On a $L = 60$, $h = 10$, $f_{\text{med}} = 18$, $F_{\text{prec}} = 18$. D'où :
\[ Me = 60 + \frac{25 - 18}{18} \times 10 = 60 + \frac{70}{18} \approx 63{,}9 \text{ kg}. \]
\item On place les points $(40\,;\,0)$, $(50\,;\,6)$, $(60\,;\,18)$, $(70\,;\,36)$, $(80\,;\,45)$, $(90\,;\,50)$, on les relie, puis on lit l'abscisse du point d'ordonnée $25$ : on retrouve environ $64$ kg.
\item $Me \approx 63{,}9$ kg et $\bar{x} = 64$ kg sont très proches : la distribution est presque symétrique, les quelques sacs lourds de la classe $[80\,;\,90[$ tirent très légèrement la moyenne vers le haut.
\end{enumerate}
\end{corrige}

\coursec{Caractéristiques de dispersion : variance et écart-type}

Dans la section précédente, nous avons appris à situer le « centre » d'une série statistique grâce à la moyenne et à la médiane. Mais connaître le centre ne suffit pas. Imagine deux vendeuses du marché Mokolo à Yaoundé : chacune vend en moyenne 30 régimes de plantains par semaine. La première vend entre 28 et 32 régimes chaque semaine ; la seconde vend tantôt 5 régimes, tantôt 55. Leurs moyennes sont identiques, pourtant leurs situations sont totalement différentes : la première peut prévoir ses achats sereinement, la seconde vit dans l'incertitude. Il nous manque un outil pour mesurer cette différence : un indicateur de \cle{dispersion}, qui quantifie à quel point les valeurs s'éloignent de la moyenne. C'est précisément le rôle de la \cle{variance} et de l'\cle{écart-type}.

\courssub{Pourquoi mesurer la dispersion ? Intuition}

Pour mesurer à quel point une valeur $x_i$ s'éloigne de la moyenne $\bar{x}$, l'idée naturelle est de regarder l'\cle{écart} $x_i - \bar{x}$. Mais si l'on additionne tous ces écarts, on obtient toujours $0$ : les écarts positifs compensent exactement les écarts négatifs. Il faut donc rendre tous les écarts positifs. Deux choix sont possibles : prendre la valeur absolue $|x_i - \bar{x}|$, ou prendre le carré $(x_i - \bar{x})^2$. Les mathématiciens ont retenu le carré, pour deux raisons : le carré est une fonction dérivable (ce qui simplifie tous les calculs théoriques), et il accentue les grands écarts, ce qui rend l'indicateur sensible aux valeurs extrêmes — exactement ce que l'on souhaite.

La variance est donc la \emph{moyenne des carrés des écarts à la moyenne}.

\courssub{La variance : définition}

\begin{definition}[Variance d'une série statistique]
Soit une série statistique d'effectif total $N$, de valeurs $x_1, x_2, \ldots, x_p$ d'effectifs respectifs $n_1, n_2, \ldots, n_p$, et de moyenne $\bar{x}$. La \cle{variance} de la série, notée $\sigma^2$ (ou $V$), est le nombre positif défini par :
\[ \sigma^2 = \frac{\displaystyle\sum_{i=1}^{p} n_i \left(x_i - \bar{x}\right)^2}{\displaystyle\sum_{i=1}^{p} n_i} = \frac{1}{N}\sum_{i=1}^{p} n_i \left(x_i - \bar{x}\right)^2. \]
Pour une série \textbf{regroupée en classes}, on remplace chaque valeur $x_i$ par le \cle{centre} $c_i$ de la classe :
\[ \sigma^2 = \frac{1}{N}\sum_{i=1}^{p} n_i \left(c_i - \bar{x}\right)^2. \]
\end{definition}

\begin{aretenir}[Lecture de la formule]
La variance se lit ainsi : « pour chaque valeur, je calcule son écart à la moyenne, j'élève cet écart au carré, je multiplie par l'effectif, j'additionne le tout, et je divise par l'effectif total ». C'est bien une \emph{moyenne} (pondérée) des carrés des écarts.
\end{aretenir}

\begin{attention}[La variance n'est pas dans l'unité de la variable]
Parce que l'on élève les écarts au carré, la variance s'exprime dans l'unité \emph{au carré}. Si les ventes sont en régimes de plantains, la variance est en « régimes$^2$ », ce qui n'a pas de signification concrète. C'est pourquoi on introduit l'écart-type, qui « défait » le carré.
\end{attention}

\courssub{L'écart-type : définition et interprétation}

\begin{definition}[Écart-type]
L'\cle{écart-type} d'une série statistique, noté $\sigma$, est la racine carrée de la variance :
\[ \sigma = \sqrt{\sigma^2} = \sqrt{\frac{1}{N}\sum_{i=1}^{p} n_i \left(x_i - \bar{x}\right)^2}. \]
\end{definition}

L'écart-type s'interprète comme un \cle{écart moyen quadratique} : il donne l'ordre de grandeur de la distance « typique » entre les valeurs de la série et la moyenne.

\begin{propriete}[Interprétation de l'écart-type]
\begin{itemize}
\item Plus $\sigma$ est \textbf{petit}, plus les valeurs sont \textbf{concentrées} autour de la moyenne : la série est \emph{homogène}, la moyenne est représentative.
\item Plus $\sigma$ est \textbf{grand}, plus les valeurs sont \textbf{dispersées} : la série est \emph{hétérogène}, la moyenne résume mal la situation.
\item On a toujours $\sigma \geq 0$, et $\sigma = 0$ si et seulement si toutes les valeurs sont égales.
\end{itemize}
\end{propriete}

\begin{savaistu}[L'écart-type parle votre langue]
Contrairement à la variance, l'écart-type s'exprime \textbf{dans la même unité que la variable} : en régimes de plantains, en francs CFA, en notes sur 20, en minutes. Dire « l'écart-type des notes de la classe est de 3 points » se comprend immédiatement : les notes s'écartent typiquement d'environ 3 points de la moyenne. C'est ce qui fait de l'écart-type l'indicateur de dispersion le plus utilisé au monde, des laboratoires pharmaceutiques aux sondages électoraux.
\end{savaistu}

\courssub{La formule de König–Huygens : un calcul allégé}

La formule de définition exige de calculer d'abord $\bar{x}$, puis chaque écart $x_i - \bar{x}$, puis son carré : c'est long, et les arrondis de $\bar{x}$ faussent le résultat. Il existe une formule équivalente, bien plus pratique.

\begin{propriete}[Formule de König–Huygens]
La variance peut se calculer par la formule dite de \cle{König–Huygens} :
\[ \sigma^2 = \frac{\displaystyle\sum_{i=1}^{p} n_i x_i^2}{\displaystyle\sum_{i=1}^{p} n_i} - \bar{x}^2 = \frac{1}{N}\sum_{i=1}^{p} n_i x_i^2 - \bar{x}^2. \]
Autrement dit : \emph{la variance est la moyenne des carrés moins le carré de la moyenne}. Pour une série en classes, on remplace $x_i$ par le centre $c_i$.
\end{propriete}

\begin{experience}[Démonstration guidée de la formule de König–Huygens]
Cette démonstration est formatrice et sa démarche est exigible. Suivons-la pas à pas. On part de la définition et on développe le carré :
\begin{align*}
\sum_{i=1}^{p} n_i (x_i - \bar{x})^2 &= \sum_{i=1}^{p} n_i \left(x_i^2 - 2\bar{x}\,x_i + \bar{x}^2\right) & &\text{(identité remarquable)}\\
&= \sum_{i=1}^{p} n_i x_i^2 - 2\bar{x}\sum_{i=1}^{p} n_i x_i + \bar{x}^2 \sum_{i=1}^{p} n_i & &\text{(on distribue et on sort les constantes)}
\end{align*}
Or, par définition de la moyenne, $\displaystyle\sum_{i=1}^{p} n_i x_i = N\bar{x}$, et $\displaystyle\sum_{i=1}^{p} n_i = N$. En remplaçant :
\begin{align*}
\sum_{i=1}^{p} n_i (x_i - \bar{x})^2 &= \sum_{i=1}^{p} n_i x_i^2 - 2\bar{x} \times N\bar{x} + \bar{x}^2 \times N\\
&= \sum_{i=1}^{p} n_i x_i^2 - 2N\bar{x}^2 + N\bar{x}^2\\
&= \sum_{i=1}^{p} n_i x_i^2 - N\bar{x}^2.
\end{align*}
En divisant les deux membres par $N$, on obtient :
\[ \sigma^2 = \frac{1}{N}\sum_{i=1}^{p} n_i x_i^2 - \bar{x}^2. \]
La formule est démontrée. Retiens l'idée clé : développer le carré, puis utiliser les deux identités $\sum n_i x_i = N\bar{x}$ et $\sum n_i = N$.
\end{experience}

\begin{methode}[Calculer variance et écart-type avec König–Huygens]
\begin{enumerate}
\item Compléter le tableau avec deux colonnes : $n_i x_i$ (pour la moyenne) et $n_i x_i^2$ (pour la variance). Pour une série en classes, utiliser les centres $c_i$ : colonnes $n_i c_i$ et $n_i c_i^2$.
\item Calculer la moyenne : $\bar{x} = \dfrac{\sum n_i x_i}{N}$.
\item Calculer la moyenne des carrés : $\dfrac{\sum n_i x_i^2}{N}$.
\item Appliquer König–Huygens : $\sigma^2 = \dfrac{\sum n_i x_i^2}{N} - \bar{x}^2$.
\item Prendre la racine carrée : $\sigma = \sqrt{\sigma^2}$, en arrondissant raisonnablement (souvent au centième).
\end{enumerate}
Cette méthode demande moins d'opérations que la définition et est plus stable numériquement : on évite de soustraire un $\bar{x}$ arrondi à chaque valeur.
\end{methode}

\begin{attention}[Pièges fréquents au Probatoire]
\begin{itemize}
\item \textbf{Diviser par $N$, pas par $N-1$} : pour décrire une population étudiée (le cadre de la Première), on divise par l'effectif total $N$. L'estimateur « non biaisé » avec $N-1$ n'intervient que lorsqu'on veut estimer la variance d'une population à partir d'un échantillon : ce sera vu en Terminale.
\item Ne pas confondre $\dfrac{\sum n_i x_i^2}{N}$ et $\bar{x}^2$ : ce sont deux quantités différentes, et c'est leur \emph{différence} qui donne la variance.
\item Une variance est toujours \textbf{positive ou nulle}. Si ton calcul donne un nombre négatif, il y a forcément une erreur (souvent un arrondi de $\bar{x}$ trop brutal).
\item Ne jamais écrire $\sigma^2 = \sum n_i(x_i - \bar{x})^2$ en oubliant la division par $N$.
\end{itemize}
\end{attention}

\courssub{Exemple chiffré : les ventes de plantains}

\begin{exoresolu}[Variance et écart-type des ventes de plantains]
Mama Françoise, vendeuse au marché de Bafoussam, a relevé ses ventes quotidiennes de régimes de plantains pendant 20 jours :
\begin{center}
\begin{tabular}{|l|ccccc|}
\hline
\rowcolor{popblueL} Nombre de régimes $x_i$ & 10 & 20 & 30 & 40 & 50 \\
\hline
Nombre de jours $n_i$ & 2 & 5 & 7 & 4 & 2 \\
\hline
\end{tabular}
\end{center}
Calculer la moyenne, la variance et l'écart-type de cette série, puis interpréter.
\tcblower
\textbf{Solution détaillée.} Complétons le tableau avec les colonnes $n_i x_i$ et $n_i x_i^2$ :
\begin{center}
\begin{tabular}{|l|cccccc|}
\hline
\rowcolor{popblueL} $x_i$ & 10 & 20 & 30 & 40 & 50 & \textbf{Totaux} \\
\hline
$n_i$ & 2 & 5 & 7 & 4 & 2 & $N = 20$ \\
\hline
$n_i x_i$ & 20 & 100 & 210 & 160 & 100 & $590$ \\
\hline
$n_i x_i^2$ & 200 & 2000 & 6300 & 6400 & 5000 & $19\,900$ \\
\hline
\end{tabular}
\end{center}
\textbf{Étape 1 — Moyenne :}
\[ \bar{x} = \frac{590}{20} = 29{,}5 \text{ régimes.} \]
\textbf{Étape 2 — Moyenne des carrés :}
\[ \frac{\sum n_i x_i^2}{N} = \frac{19\,900}{20} = 995. \]
\textbf{Étape 3 — Variance par König–Huygens :}
\[ \sigma^2 = 995 - (29{,}5)^2 = 995 - 870{,}25 = 124{,}75. \]
\textbf{Étape 4 — Écart-type :}
\[ \sigma = \sqrt{124{,}75} \approx 11{,}17 \text{ régimes.} \]
\textbf{Interprétation :} Mama Françoise vend en moyenne 29,5 régimes par jour, avec un écart-type d'environ 11,2 régimes. Ses ventes quotidiennes s'écartent donc typiquement d'une dizaine de régimes de la moyenne : l'activité est assez variable. Pour sécuriser ses achats auprès des producteurs, elle devra prévoir une marge.
\end{exoresolu}

\courssub{Comparer deux séries de même moyenne}

Revenons à nos deux vendeuses du marché Mokolo. Les deux séries ci-dessous ont la même moyenne $\bar{x} = 30$ régimes, mais des dispersions très différentes.

\begin{popfigure}
\begin{center}
\begin{tikzpicture}
\begin{axis}[
  width=0.95\linewidth, height=6.2cm,
  ybar, bar width=7pt,
  ymin=0, ymax=9,
  xmin=0, xmax=60,
  xlabel={Ventes hebdomadaires (régimes)},
  ylabel={Effectif (semaines)},
  xtick={10,20,30,40,50},
  legend style={at={(0.5,1.02)}, anchor=south, legend columns=2, draw=none},
  axis lines=left,
]
\addplot+[ybar, fill=popblue, draw=popdark] coordinates {(10,0) (20,1) (25,3) (30,4) (35,3) (40,1) (50,0)};
\addlegendentry{Vendeuse A : $\sigma \approx 5$}
\addplot+[ybar, fill=poporange, draw=popdark] coordinates {(10,2) (20,2) (25,1) (30,1) (35,1) (40,2) (50,2)};
\addlegendentry{Vendeuse B : $\sigma \approx 14$}
\end{axis}
\end{tikzpicture}
\end{center}
\legende{Deux séries de même moyenne $\bar{x} = 30$ : la série A (bleue) est concentrée autour de la moyenne (petit écart-type), la série B (orange) est très dispersée (grand écart-type).}
\end{popfigure}

\begin{propriete}[Rôle de l'écart-type dans la comparaison]
Pour comparer deux séries de \textbf{même nature} (même unité), la moyenne indique le niveau général, et l'écart-type indique la régularité :
\begin{itemize}
\item à moyennes égales, la série au plus petit écart-type est la plus \emph{régulière}, la plus prévisible ;
\item à moyennes différentes, on compare d'abord les moyennes, puis on nuance avec les écarts-types.
\end{itemize}
\end{propriete}

\begin{exemplebox}[Notes de deux classes de Première A]
Deux classes ont toutes deux une moyenne de $11/20$ au même devoir de mathématiques. La classe A a un écart-type de $2$ points, la classe B un écart-type de $5$ points. En classe A, les notes sont homogènes : la plupart des élèves sont autour de 11. En classe B, les notes sont très dispersées : d'excellents élèves côtoient des élèves en grande difficulté. Le professeur de la classe B devra différencier son accompagnement, alors que celui de la classe A peut miser sur un enseignement uniforme. Même moyenne, réalités pédagogiques opposées : voilà toute la puissance de l'écart-type.
\end{exemplebox}

\begin{exercice}[À toi de jouer]
On a relevé les durées de communication (en minutes) d'un échantillon d'abonnés MTN : les classes sont $[0;10[$, $[10;20[$, $[20;30[$, $[30;40[$, d'effectifs respectifs $8$, $15$, $12$, $5$.
\begin{enumerate}
\item Construire le tableau avec les centres $c_i$, les colonnes $n_i c_i$ et $n_i c_i^2$.
\item Calculer la moyenne $\bar{x}$.
\item Calculer la variance par la formule de König–Huygens, puis l'écart-type. Interpréter.
\end{enumerate}
\end{exercice}

\begin{corrige}[Exercice : durées de communication]
\begin{enumerate}
\item Les centres sont $c_1 = 5$, $c_2 = 15$, $c_3 = 25$, $c_4 = 35$ et $N = 8+15+12+5 = 40$.
\begin{center}
\begin{tabular}{|l|ccccc|}
\hline
\rowcolor{popblueL} $c_i$ & 5 & 15 & 25 & 35 & \textbf{Totaux} \\
\hline
$n_i$ & 8 & 15 & 12 & 5 & $40$ \\
\hline
$n_i c_i$ & 40 & 225 & 300 & 175 & $740$ \\
\hline
$n_i c_i^2$ & 200 & 3375 & 7500 & 6125 & $17\,200$ \\
\hline
\end{tabular}
\end{center}
\item $\bar{x} = \dfrac{740}{40} = 18{,}5$ minutes.
\item Variance : $\sigma^2 = \dfrac{17\,200}{40} - (18{,}5)^2 = 430 - 342{,}25 = 87{,}75$. Écart-type : $\sigma = \sqrt{87{,}75} \approx 9{,}37$ minutes. Les durées de communication s'écartent typiquement d'environ 9,4 minutes de la durée moyenne de 18,5 minutes : les comportements des abonnés sont assez dispersés.
\end{enumerate}
\end{corrige}

\begin{aretenir}[L'essentiel de la section]
\begin{itemize}
\item La \cle{variance} mesure la dispersion : $\sigma^2 = \dfrac{1}{N}\sum n_i (x_i - \bar{x})^2$ (avec les centres $c_i$ pour les classes).
\item La formule de \cle{König–Huygens} allège les calculs : $\sigma^2 = \dfrac{1}{N}\sum n_i x_i^2 - \bar{x}^2$ — « moyenne des carrés moins carré de la moyenne ».
\item L'\cle{écart-type} $\sigma = \sqrt{\sigma^2}$ s'exprime dans l'unité de la variable : petit $\sigma$ = série homogène ; grand $\sigma$ = série dispersée.
\item On divise toujours par l'effectif total $N$ en Première.
\end{itemize}
\end{aretenir}

\coursec{Représentations graphiques : histogramme et polygones des effectifs cumulés}

Dans la section précédente, nous avons résumé une série statistique par des nombres : la moyenne, la médiane, la variance et l'écart-type. Ces indicateurs sont précieux, mais ils restent abstraits. Un gestionnaire du marché de Mokolo, un chef d'établissement ou un agent de CAMTEL préfère souvent \emph{voir} les données. C'est toute la force des représentations graphiques : un bon graphique révèle d'un coup d'œil la forme de la distribution (concentrée ou étalée, symétrique ou non), là où un tableau de chiffres demande un effort de lecture. Dans cette section, tu vas apprendre à construire rigoureusement les deux représentations exigibles au programme de Première : l'\cle{histogramme} et les \cle{polygones des effectifs cumulés}, et à lire graphiquement la médiane.

\courssub{L'histogramme : représenter une série regroupée en classes}

\textbf{Intuition.} Reprenons notre exemple fil rouge : un vendeur de plantains au marché a relevé, pendant $50$ jours, la masse de plantains vendue chaque jour (en kg). Les données, regroupées en classes à la section 1, sont les suivantes :

\begin{center}
\begin{tabularx}{\linewidth}{|l|X|X|X|X|}\hline
\rowcolor{popblueL} Classe (kg) & $[0\,;10[$ & $[10\,;20[$ & $[20\,;30[$ & $[30\,;50[$ \\ \hline
Effectif $n_i$ & $5$ & $15$ & $20$ & $10$ \\ \hline
Fréquence $f_i$ & $0,10$ & $0,30$ & $0,40$ & $0,20$ \\ \hline
\end{tabularx}
\end{center}

Comment dessiner cela ? L'idée naturelle est de représenter chaque classe par un rectangle dont la \emph{base} est l'amplitude de la classe. Mais que mettre en hauteur ? Si toutes les classes ont la même amplitude, on peut mettre l'effectif directement. Mais ici, la dernière classe $[30\,;50[$ est \emph{deux fois plus large} que les autres. Si l'on mettait simplement l'effectif $10$ en hauteur, cette classe paraîtrait « petite » alors qu'elle couvre une large plage de valeurs : le dessin mentirait. La solution consiste à raisonner en \emph{aires}.

\begin{definition}[Histogramme et densité de fréquence]
L'\cle{histogramme} d'une série regroupée en classes est un graphique constitué de rectangles juxtaposés :
\begin{itemize}
\item la \textbf{base} de chaque rectangle, portée sur l'axe des abscisses, est l'intervalle de la classe (sa longueur est l'\cle{amplitude} $a_i$) ;
\item la \textbf{hauteur} du rectangle est la \cle{densité} de la classe : densité d'effectif $d_i = \dfrac{n_i}{a_i}$ ou densité de fréquence $d_i = \dfrac{f_i}{a_i}$.
\end{itemize}
Ainsi, l'\textbf{aire} de chaque rectangle est proportionnelle à l'effectif (ou égale à la fréquence) de la classe.
\end{definition}

\begin{attention}[Le piège des classes d'amplitudes inégales]
Si les classes n'ont \textbf{pas} toutes la même amplitude, il est \textbf{interdit} de porter les effectifs (ou les fréquences) directement en hauteur. Un rectangle deux fois plus large aurait, à hauteur égale, une aire deux fois plus grande : le lecteur croirait que la classe contient deux fois plus d'individus. C'est l'\emph{aire}, et non la hauteur, qui doit représenter l'effectif. En revanche, si toutes les classes ont la même amplitude, diviser par $a_i$ ne change que l'échelle verticale : on peut alors porter directement les effectifs en hauteur.
\end{attention}

\begin{propriete}[Aire totale de l'histogramme]
Dans un histogramme construit avec les densités de fréquences $d_i = \dfrac{f_i}{a_i}$, l'aire totale des rectangles vaut $1$ (soit $100\,\%$). En effet, l'aire du rectangle de la classe $i$ est $a_i \times d_i = a_i \times \dfrac{f_i}{a_i} = f_i$, donc :
\[ \text{Aire totale} = \sum_{i} f_i = 1. \]
De même, avec les densités d'effectifs $d_i = \dfrac{n_i}{a_i}$, l'aire totale vaut l'effectif total $N$. Cette propriété permet de \emph{vérifier} qu'un histogramme est correctement construit.
\end{propriete}

\begin{methode}[Tracer correctement un histogramme]
\begin{enumerate}
\item Repérer les classes et vérifier si les amplitudes sont égales ou non.
\item Calculer l'amplitude $a_i$ de chaque classe, puis la densité $d_i = \dfrac{n_i}{a_i}$ (ou $\dfrac{f_i}{a_i}$).
\item Choisir les échelles : sur l'axe des abscisses, une unité graphique pour un nombre fixe d'unités du caractère ; sur l'axe des ordonnées, une échelle adaptée à la plus grande densité. \textbf{Indiquer les unités sur chaque axe.}
\item Tracer, pour chaque classe, le rectangle de base la classe et de hauteur sa densité, les rectangles étant juxtaposés sans espace.
\item Vérifier : la somme des aires doit valoir $N$ (densités d'effectifs) ou $1$ (densités de fréquences).
\end{enumerate}
\end{methode}

\begin{exemplebox}[Histogramme des ventes de plantains]
Pour notre vendeur, calculons les densités d'effectifs :
\begin{center}
\begin{tabularx}{\linewidth}{|l|X|X|X|X|}\hline
\rowcolor{popblueL} Classe & $[0\,;10[$ & $[10\,;20[$ & $[20\,;30[$ & $[30\,;50[$ \\ \hline
Amplitude $a_i$ & $10$ & $10$ & $10$ & $20$ \\ \hline
Effectif $n_i$ & $5$ & $15$ & $20$ & $10$ \\ \hline
Densité $d_i = n_i/a_i$ & $0,5$ & $1,5$ & $2$ & $0,5$ \\ \hline
\end{tabularx}
\end{center}
Remarque bien la dernière colonne : l'effectif $10$ divisé par l'amplitude $20$ donne une densité de $0,5$, deux fois moins que ce qu'aurait suggéré une lecture naïve. L'aire totale vaut $10 \times 0,5 + 10 \times 1,5 + 10 \times 2 + 20 \times 0,5 = 5 + 15 + 20 + 10 = 50 = N$ : l'histogramme est cohérent.
\end{exemplebox}

\begin{popfigure}
\begin{center}
\begin{tikzpicture}
\begin{axis}[
  width=0.95\linewidth, height=7cm,
  xlabel={Masse vendue (kg)}, ylabel={Densité d'effectif},
  xmin=0, xmax=52, ymin=0, ymax=2.4,
  xtick={0,10,20,30,40,50},
  ytick={0,0.5,1,1.5,2},
  grid=major, grid style={popline!50},
  axis lines=left, tick label style={font=\small}
]
\draw[fill=popblue!45, draw=popblue, thick] (axis cs:0,0) rectangle (axis cs:10,0.5);
\draw[fill=popblue!45, draw=popblue, thick] (axis cs:10,0) rectangle (axis cs:20,1.5);
\draw[fill=popblue!45, draw=popblue, thick] (axis cs:20,0) rectangle (axis cs:30,2);
\draw[fill=poporange!45, draw=poporange, thick] (axis cs:30,0) rectangle (axis cs:50,0.5);
\node[font=\small] at (axis cs:5,0.75) {$5$ jours};
\node[font=\small] at (axis cs:15,1.75) {$15$ jours};
\node[font=\small] at (axis cs:25,2.2) {$20$ jours};
\node[font=\small] at (axis cs:40,0.75) {$10$ jours};
\end{axis}
\end{tikzpicture}
\end{center}
\legende{Histogramme en densités d'effectifs des ventes journalières de plantains. La classe $[30\,;50[$, deux fois plus large, a une hauteur réduite de moitié : c'est l'\emph{aire} de chaque rectangle qui est égale à l'effectif.}
\end{popfigure}

\courssub{Les polygones des effectifs cumulés}

\textbf{Intuition.} Supposons que le vendeur se pose la question : « Combien de jours ai-je vendu \emph{moins de} $20$ kg ? » La réponse, $5 + 15 = 20$ jours, est un effectif cumulé croissant. S'il se demande « combien de jours ai-je vendu \emph{au moins} $20$ kg ? », la réponse, $20 + 10 = 30$ jours, est un effectif cumulé décroissant. Plutôt que de refaire ces additions à chaque question, on dessine une courbe qui y répond pour \emph{toutes} les bornes possibles : c'est le polygone des effectifs cumulés.

\begin{definition}[Polygones des effectifs cumulés]
Pour une série regroupée en classes de bornes $x_0 < x_1 < \dots < x_k$ :
\begin{itemize}
\item le \cle{polygone des effectifs cumulés croissants} est la ligne brisée joignant les points $\bigl(x_i\,;\,ECC_i\bigr)$, où $ECC_i$ est le nombre d'individus dont la valeur est \textbf{strictement inférieure} à la borne supérieure $x_i$ ; on le fait partir du point $(x_0\,;\,0)$ ;
\item le \cle{polygone des effectifs cumulés décroissants} est la ligne brisée joignant les points $\bigl(x_i\,;\,ECD_i\bigr)$, où $ECD_i$ est le nombre d'individus dont la valeur est \textbf{supérieure ou égale} à la borne inférieure $x_i$ ; il se termine au point $(x_k\,;\,0)$.
\end{itemize}
On définit de même les polygones de \emph{fréquences} cumulées, en remplaçant les effectifs par les fréquences cumulées (comprises entre $0$ et $1$, ou $0$ et $100\,\%$).
\end{definition}

\begin{exemplebox}[Tableau des cumuls pour les ventes de plantains]
\begin{center}
\begin{tabularx}{\linewidth}{|l|X|X|X|X|X|}\hline
\rowcolor{popblueL} Borne $x_i$ (kg) & $0$ & $10$ & $20$ & $30$ & $50$ \\ \hline
$ECC$ (moins de $x_i$ kg) & $0$ & $5$ & $20$ & $40$ & $50$ \\ \hline
$ECD$ (au moins $x_i$ kg) & $50$ & $45$ & $30$ & $10$ & $0$ \\ \hline
\end{tabularx}
\end{center}
Le polygone croissant joint donc les points $(0\,;0)$, $(10\,;5)$, $(20\,;20)$, $(30\,;40)$, $(50\,;50)$ ; le polygone décroissant joint $(0\,;50)$, $(10\,;45)$, $(20\,;30)$, $(30\,;10)$, $(50\,;0)$.
\end{exemplebox}

\begin{popfigure}
\begin{center}
\begin{tikzpicture}
\begin{axis}[
  width=0.95\linewidth, height=7.5cm,
  xlabel={Masse vendue (kg)}, ylabel={Effectif cumulé (jours)},
  xmin=0, xmax=52, ymin=0, ymax=55,
  xtick={0,10,20,30,40,50},
  ytick={0,5,10,20,25,30,40,50},
  grid=major, grid style={popline!50},
  axis lines=left, tick label style={font=\small},
  legend style={font=\small, at={(0.03,0.97)}, anchor=north west}
]
\addplot[popblue, very thick, mark=*, mark size=1.5pt] coordinates {(0,0) (10,5) (20,20) (30,40) (50,50)};
\addlegendentry{Cumulés croissants}
\addplot[poporange, very thick, mark=square*, mark size=1.5pt] coordinates {(0,50) (10,45) (20,30) (30,10) (50,0)};
\addlegendentry{Cumulés décroissants}
\draw[dashed, popgreen, thick] (axis cs:0,25) -- (axis cs:22.5,25) -- (axis cs:22.5,0);
\node[popgreen, font=\small, anchor=west] at (axis cs:23,4) {Médiane $\approx 22,5$ kg};
\node[popgreen, font=\small, anchor=south west] at (axis cs:1,26) {$N/2 = 25$};
\end{axis}
\end{tikzpicture}
\end{center}
\legende{Polygones des effectifs cumulés croissants (bleu) et décroissants (orange). Les deux courbes se coupent à l'abscisse de la médiane ; on peut aussi lire la médiane en cherchant l'abscisse du point du polygone croissant d'ordonnée $N/2 = 25$.}
\end{popfigure}

\begin{propriete}[Lecture graphique de la médiane et des quartiles]
\begin{itemize}
\item La \cle{médiane} est l'abscisse du point d'intersection des polygones des effectifs cumulés croissants et décroissants. De façon équivalente, c'est l'abscisse du point du polygone croissant d'ordonnée $\dfrac{N}{2}$ (ou $50\,\%$ sur un polygone de fréquences cumulées).
\item Plus généralement (préparation à la Terminale), le \cle{premier quartile} $Q_1$ se lit à l'ordonnée $\dfrac{N}{4}$ (soit $25\,\%$) et le \cle{troisième quartile} $Q_3$ à l'ordonnée $\dfrac{3N}{4}$ (soit $75\,\%$) sur le polygone croissant.
\end{itemize}
Ces lectures reposent sur l'hypothèse d'\emph{équirépartition} des valeurs à l'intérieur de chaque classe : c'est exactement l'hypothèse qui justifie l'interpolation linéaire vue à la section 3.
\end{propriete}

\begin{exoresolu}[Lecture graphique complète]
Énoncé. À l'aide des polygones cumulés des ventes de plantains, déterminer graphiquement puis par interpolation : (a) la médiane ; (b) les quartiles $Q_1$ et $Q_3$ ; (c) le nombre de jours où le vendeur a vendu moins de $25$ kg.
\tcblower
Solution détaillée.
(a) On a $N = 50$, donc $\dfrac{N}{2} = 25$. Sur le polygone croissant, l'ordonnée $25$ est atteinte dans la classe $[20\,;30[$, entre les points $(20\,;20)$ et $(30\,;40)$. Par interpolation linéaire :
\[ Me = 20 + \dfrac{25 - 20}{40 - 20}\times 10 = 20 + 2{,}5 = 22{,}5 \text{ kg}. \]
Graphiquement, c'est aussi l'abscisse du point de croisement des deux polygones (pointillés verts sur la figure).
(b) $\dfrac{N}{4} = 12{,}5$ : cette ordonnée est atteinte dans la classe $[10\,;20[$ entre $(10\,;5)$ et $(20\,;20)$ :
\[ Q_1 = 10 + \dfrac{12{,}5 - 5}{20 - 5}\times 10 = 10 + 5 = 15 \text{ kg}. \]
$\dfrac{3N}{4} = 37{,}5$ : atteinte dans $[20\,;30[$ entre $(20\,;20)$ et $(30\,;40)$ :
\[ Q_3 = 20 + \dfrac{37{,}5 - 20}{40 - 20}\times 10 = 20 + 8{,}75 = 28{,}75 \text{ kg}. \]
Ainsi, la moitié centrale des journées correspond à des ventes comprises entre $15$ kg et $28{,}75$ kg.
(c) Pour $x = 25$, milieu de la classe $[20\,;30[$, l'interpolation sur le polygone croissant donne :
\[ ECC(25) = 20 + \dfrac{25 - 20}{30 - 20}\times(40 - 20) = 20 + 10 = 30. \]
Le vendeur a vendu moins de $25$ kg pendant environ $30$ jours sur $50$, soit $60\,\%$ du temps.
\end{exoresolu}

\begin{attention}[Erreurs fréquentes sur les polygones cumulés]
\begin{itemize}
\item \textbf{Mauvaise abscisse} : pour le polygone \emph{croissant}, on place les effectifs cumulés aux \emph{bornes supérieures} des classes ; pour le \emph{décroissant}, aux \emph{bornes inférieures}. Inverser les deux décale toute la courbe.
\item \textbf{Oublier les points d'attache} : le polygone croissant doit partir de $(x_0\,;\,0)$ et le décroissant de $(x_0\,;\,N)$. Sans ces points, la lecture graphique de la médiane par intersection est faussée.
\item \textbf{Relier par une courbe arrondie} : ce sont des \emph{segments de droite} (d'où le nom « polygone »), car on suppose la répartition uniforme dans chaque classe.
\end{itemize}
\end{attention}

\begin{savaistu}[Des polygones cumulés à la mesure des inégalités]
En Terminale, tu rencontreras la \cle{courbe de Lorenz} : on y reporte les fréquences cumulées de la \emph{population} (par exemple les ménages camerounais rangés du plus pauvre au plus riche) contre les fréquences cumulées du \emph{caractère} (le revenu). Plus la courbe s'écarte de la diagonale, plus les inégalités sont fortes. On mesure cet écart par l'\cle{indice de Gini}, un nombre compris entre $0$ (égalité parfaite) et $1$ (inégalité maximale), utilisé par l'Institut National de la Statistique du Cameroun et la Banque mondiale pour comparer les pays. Tout cela n'est qu'une généralisation des polygones de fréquences cumulées que tu viens d'apprendre à tracer !
\end{savaistu}

\begin{aretenir}[L'essentiel de la section]
\begin{itemize}
\item Dans un histogramme, c'est l'\textbf{aire} du rectangle qui représente l'effectif : la hauteur est la densité $d_i = \dfrac{n_i}{a_i}$. Avec les densités de fréquences, l'aire totale vaut $1$.
\item Le polygone des effectifs cumulés \textbf{croissants} joint les points (borne supérieure ; $ECC$) ; le \textbf{décroissant} joint les points (borne inférieure ; $ECD$).
\item La médiane se lit à l'ordonnée $\dfrac{N}{2}$ sur le polygone croissant, ou à l'intersection des deux polygones ; $Q_1$ et $Q_3$ se lisent aux ordonnées $\dfrac{N}{4}$ et $\dfrac{3N}{4}$.
\end{itemize}
\end{aretenir}

\coursec{Interprétation et communication des résultats d'une enquête}

Dans la section précédente, tu as appris à construire un histogramme et les polygones des effectifs cumulés. Mais un tableau ou un graphique, aussi beau soit-il, ne parle pas tout seul : il faut encore le \emph{lire}, le \emph{commenter} et le \emph{communiquer}. C'est toute la différence entre un élève qui « calcule » et un statisticien qui \emph{conclut}. Dans cette dernière grande étape de la leçon, nous allons apprendre à rédiger un compte-rendu d'enquête complet, à croiser intelligemment la moyenne, la médiane et l'écart-type, à choisir la bonne représentation graphique selon le public visé, et surtout à connaître les \cle{limites} de nos conclusions. C'est exactement le travail que font les instituts de sondage, les journalistes économiques ou les services de planification du MINESEC lorsqu'ils analysent les résultats du Probatoire.

\courssub{Rédiger un compte-rendu statistique}

\textbf{Pourquoi rédiger ?}\par
Une enquête statistique n'a de valeur que si elle peut être comprise et vérifiée par d'autres. Imagine que le proviseur de ton lycée te demande d'étudier le temps de trajet domicile--lycée des élèves pour décider de l'heure d'ouverture des portes. Si tu lui remets simplement un tableau de nombres, il ne pourra rien en faire. Il te faut un \cle{rapport} structuré, qui raconte l'enquête du début à la fin.

\begin{methode}[Structure d'un rapport d'enquête statistique]
Un bon rapport d'enquête se rédige en cinq parties, toujours dans le même ordre :
\begin{enumerate}
\item \textbf{Introduction} : on présente le problème posé et l'objectif de l'enquête (« Pourquoi cette étude ? »).
\item \textbf{Méthode} : on décrit la \cle{population} étudiée (qui ? combien ?), le \cle{caractère} observé (quoi ? avec quelle unité ?), et le \cle{mode de collecte} des données (questionnaire, mesure directe, registre, date et lieu de l'enquête).
\item \textbf{Résultats} : on présente les tableaux (effectifs, fréquences, cumuls), les paramètres calculés (moyenne, médiane, écart-type) et les graphiques (histogramme, polygones).
\item \textbf{Interprétation} : on commente les résultats en langage courant, en croisant les indicateurs (voir 6.2).
\item \textbf{Conclusion} : on répond à la question de départ, on propose éventuellement une décision, et on signale honnêtement les limites de l'étude.
\end{enumerate}
\end{methode}

\begin{attention}[Présente toujours la population et le caractère avant tout calcul]
Une moyenne de « 25 » ne veut rien dire si l'on ne sait pas \emph{qui} a été interrogé et \emph{quoi} on a mesuré. Dans toute copie de Probatoire et dans tout rapport, écris systématiquement une phrase du type : « La population étudiée est l'ensemble des 40 élèves de la classe de Première A du lycée de \ldots ; le caractère étudié est le temps de trajet domicile--lycée, exprimé en minutes ; les données ont été recueillies par questionnaire le 12 mars. » Oublier cette présentation est une faute de méthode, même si tous les calculs sont justes.
\end{attention}

\courssub{Analyse conjointe : moyenne, médiane et écart-type}

\textbf{L'intuition.}\par
Un seul indicateur ne suffit jamais à décrire une série. Deux classes peuvent avoir la même moyenne de notes et des réalités totalement différentes : dans l'une, tous les élèves sont autour de 10 ; dans l'autre, la moitié a 4 et l'autre moitié 16. C'est en \emph{croisant} les trois paramètres que le portrait devient fidèle.

\begin{propriete}[Lecture croisée des trois paramètres]
Soit une série statistique de moyenne $\bar{x}$, de médiane $Me$ et d'écart-type $\sigma$.
\begin{itemize}
\item \textbf{Comparaison moyenne--médiane (forme de la distribution).} Si $\bar{x} \approx Me$, la distribution est à peu près \cle{symétrique}. Si $\bar{x} > Me$, la série est \cle{étalée vers les grandes valeurs} (quelques valeurs élevées « tirent » la moyenne vers le haut) : on parle d'asymétrie à droite. Si $\bar{x} < Me$, la série est étalée vers les petites valeurs (asymétrie à gauche).
\item \textbf{Écart-type (concentration).} Plus $\sigma$ est petit devant $\bar{x}$, plus les données sont \cle{concentrées} autour de la moyenne (population homogène). Plus $\sigma$ est grand, plus les données sont \cle{dispersées} (population hétérogène). En pratique, pour beaucoup de séries « raisonnablement régulières » (unimodales, symétriques, type « cloche »), environ les deux tiers des données se trouvent dans l'intervalle $[\bar{x}-\sigma\,;\,\bar{x}+\sigma]$.
\end{itemize}
\end{propriete}

\begin{attention}[La moyenne est sensible aux valeurs extrêmes, pas la médiane]
Remplace une seule donnée par une valeur très grande : la moyenne augmente fortement, la médiane ne bouge presque pas. Ainsi, quand une série contient des valeurs exceptionnelles (un très gros salaire, un trajet de 2 heures), la médiane est souvent un meilleur résumé de la « valeur typique » que la moyenne. Dire le contraire, ou comparer deux séries uniquement par leurs moyennes alors que leurs dispersions sont très différentes, est l'erreur d'interprétation la plus fréquente.
\end{attention}

\begin{exemplebox}[Interprétation complète : les ventes de plantains]
Une coopérative du village de Obala a relevé, pendant 20 semaines, la masse hebdomadaire de plantains vendue au marché (en kg). Les calculs des sections précédentes donnent : moyenne $\bar{x} = 45$ kg, médiane $Me = 42$ kg, écart-type $\sigma = 12$ kg. Commentons :
\begin{itemize}
\item \textbf{Position.} En moyenne, la coopérative vend 45 kg par semaine ; une semaine sur deux, elle vend moins de 42 kg.
\item \textbf{Forme.} On a $\bar{x} = 45 > Me = 42$ : la moyenne dépasse la médiane, ce qui signale quelques semaines de \emph{très fortes} ventes (par exemple pendant les fêtes) qui tirent la moyenne vers le haut. La distribution est légèrement asymétrique vers les grandes valeurs.
\item \textbf{Dispersion.} $\sigma = 12$ kg, soit environ $27\,\%$ de la moyenne : les ventes sont assez dispersées. L'intervalle $[\bar{x}-\sigma\,;\,\bar{x}+\sigma] = [33\,;\,57]$ contient la majorité des semaines : on peut annoncer prudemment que « une semaine ordinaire, la coopérative vend entre 33 et 57 kg ».
\item \textbf{Décision.} Pour fixer un stock hebdomadaire de sécurité, viser 45 kg (la moyenne) est risqué la moitié du temps ; viser plutôt $\bar{x}+\sigma = 57$ kg couvre la grande majorité des semaines.
\end{itemize}
Voilà ce qu'est une \emph{interprétation} : chaque nombre est traduit en une phrase utile pour agir.
\end{exemplebox}

\begin{savaistu}[Pourquoi les journalistes préfèrent la médiane pour les salaires]
Quand un journal annonce « le salaire médian au Cameroun est de \ldots FCFA » plutôt que le salaire moyen, ce n'est pas un hasard. Quelques très hauts revenus (grands dirigeants, footballeurs professionnels) gonflent énormément la moyenne, qui donnerait une image trop optimiste du niveau de vie. La médiane, elle, n'est pas influencée par ces valeurs extrêmes : elle dit simplement que la moitié des salariés gagne moins, et l'autre moitié plus. Retiens ce réflexe de citoyen : quand on te donne une « moyenne », demande toujours si la médiane ne raconterait pas une autre histoire !
\end{savaistu}

\courssub{Choisir la représentation graphique adaptée au public}

\textbf{L'intuition.}\par
Un graphique est un \emph{message}. La question n'est pas « quel graphique est le plus beau ? » mais « quel graphique fait passer le mieux \emph{mon} message à \emph{ce} public ? ».

\begin{aretenir}[Quel graphique pour quel message ?]
\begin{itemize}
\item \textbf{Histogramme} : il montre la \cle{forme de la distribution} (où se concentrent les données, y a-t-il des classes dominantes, la répartition est-elle symétrique ?). C'est le graphique à présenter quand le message est : « voici comment la population se répartit ».
\item \textbf{Polygone des effectifs (ou fréquences) cumulés croissants} : il répond aux questions de \cle{position} et de \cle{seuil} : « combien d'individus en dessous de telle valeur ? », « quelle est la médiane ? ». C'est le graphique à présenter quand le message est : « voici où se situe la moitié (ou tel pourcentage) de la population ».
\item \textbf{Polygone cumulé décroissant} : il répond à la question inverse : « combien d'individus \emph{au-dessus} de telle valeur ? » (utile pour les seuils de réussite, les dépassements de norme).
\end{itemize}
Devant un public non mathématicien (conseil d'administration, parents d'élèves), on privilégie l'histogramme, plus intuitif ; devant un jury ou un correcteur, on ajoute le polygone cumulé qui prouve la lecture graphique de la médiane.
\end{aretenir}

\begin{exemplebox}[Mini-rapport visuel : l'enquête « temps de trajet »]
La figure ci-dessous résume l'enquête menée auprès des 40 élèves d'une classe de Première A sur le temps de trajet domicile--lycée (en minutes), regroupés en classes d'amplitude 10. On y voit côte à côte le tableau récapitulatif, l'histogramme (forme de la distribution) et le polygone des effectifs cumulés croissants (lecture de la médiane, voisine de 25 minutes).
\end{exemplebox}

\begin{popfigure}
\begin{center}
\begin{tabular}{|c|c|c|}
\hline
\rowcolor{popblueL} Classe (min) & Effectif $n_i$ & Effectif cumulé croissant \\
\hline
$[0\,;\,10[$ & 4 & 4 \\
\hline
$[10\,;\,20[$ & 10 & 14 \\
\hline
$[20\,;\,30[$ & 14 & 28 \\
\hline
$[30\,;\,40[$ & 8 & 36 \\
\hline
$[40\,;\,50[$ & 4 & 40 \\
\hline
\end{tabular}
\end{center}
\vspace{0.4cm}
\begin{center}
\begin{tikzpicture}
\begin{axis}[
  width=0.48\linewidth, height=5.2cm,
  ybar, bar width=10pt,
  ymin=0, ymax=16,
  xmin=-2, xmax=52,
  xtick={0,10,20,30,40,50},
  xlabel={temps (min)}, ylabel={effectif},
  xlabel style={font=\small}, ylabel style={font=\small},
  tick label style={font=\small},
  axis lines=left,
  title={\small Histogramme}, title style={font=\small}
]
\addplot+[ybar, fill=popblue!70, draw=popblue] coordinates {(5,4) (15,10) (25,14) (35,8) (45,4)};
\end{axis}
\end{tikzpicture}
\hfill
\begin{tikzpicture}
\begin{axis}[
  width=0.48\linewidth, height=5.2cm,
  ymin=0, ymax=44,
  xmin=0, xmax=52,
  xtick={0,10,20,30,40,50},
  xlabel={temps (min)}, ylabel={effectif cumulé},
  xlabel style={font=\small}, ylabel style={font=\small},
  tick label style={font=\small},
  axis lines=left,
  title={\small Polygone cumulé croissant}, title style={font=\small}
]
\addplot[popcurve, mark=*, mark size=1.5pt] coordinates {(0,0) (10,4) (20,14) (30,28) (40,36) (50,40)};
\addplot[dashed, poporange, thick] coordinates {(0,20) (25,20)};
\addplot[dashed, poporange, thick] coordinates {(25,0) (25,20)};
\node[font=\small, poporange, anchor=south west] at (axis cs:1,20.5) {$N/2 = 20$};
\node[font=\small, poporange, anchor=north] at (axis cs:25,1) {$Me \approx 25$};
\end{axis}
\end{tikzpicture}
\end{center}
\legende{Mini-rapport visuel de l'enquête « temps de trajet » : tableau récapitulatif, histogramme (distribution concentrée sur $[20\,;\,30[$) et polygone des effectifs cumulés croissants (lecture graphique de la médiane $Me \approx 25$ min).}
\end{popfigure}

\textbf{Lecture du mini-rapport.} L'histogramme montre une distribution presque symétrique, dominée par la classe $[20\,;\,30[$ : la plupart des élèves mettent entre 20 et 30 minutes. Le polygone cumulé croissant coupe la droite d'ordonnée $N/2 = 20$ en une abscisse voisine de 25 : la moitié des élèves met moins de 25 minutes pour venir au lycée. La moyenne calculée à partir des centres de classes vaut $\bar{x} = 24,5$ min ; la médiane lue graphiquement est voisine de 24,3 min (environ 25 min). Ici $\bar{x} \approx Me$, ce qui confirme la quasi-symétrie observée sur l'histogramme. Tout est cohérent : c'est cela, un rapport bien construit.

\courssub{Exemples d'enquêtes élèves et critique des limites}

\textbf{Trois enquêtes à portée de classe.}\par
Le programme recommande de conduire de \emph{petites enquêtes} réelles. En voici trois, faciles à organiser :
\begin{itemize}
\item \textbf{Temps de trajet domicile--école} (caractère continu, regroupé en classes de 10 min) : c'est l'exemple traité ci-dessus. Collecte par questionnaire anonyme, en une heure de cours.
\item \textbf{Nombre de frères et sœurs} (caractère \emph{discret}, modalités $0, 1, 2, \ldots$) : inutile de faire des classes ; on dresse directement le tableau des effectifs par modalité, et l'on compare moyenne et médiane. Dans beaucoup de classes camerounaises, on observe $\bar{x} > Me$ à cause de quelques familles très nombreuses : belle illustration de l'asymétrie à droite.
\item \textbf{Hauteur des plants du jardin scolaire} (caractère continu, mesuré à la règle en cm, regroupé en classes de 5 cm) : enquête idéale en lien avec les travaux de SVT ; l'écart-type mesure alors l'homogénéité de la croissance, information utile pour juger de la qualité du semis ou de l'arrosage.
\end{itemize}

\begin{attention}[Ne conclus jamais au-delà de ce que tes données permettent]
C'est la règle d'or de l'honnêteté statistique. Si tu as interrogé \emph{ta} classe de 40 élèves, tes conclusions valent pour \emph{ta} classe, éventuellement pour ton lycée avec prudence, mais certainement pas pour « tous les élèves de Douala » ou « les jeunes du Cameroun ». Trois limites doivent toujours être mentionnées dans la conclusion d'un rapport :
\begin{enumerate}
\item \textbf{Effectif réduit} : avec 30 ou 40 données, un seul questionnaire mal rempli peut déformer la moyenne. Plus l'effectif est grand, plus les résultats sont stables.
\item \textbf{Classes arbitraires} : le découpage en classes est un choix de l'enquêteur. Avec des classes de 5 min au lieu de 10 min, l'histogramme aurait une autre allure et la médiane lue graphiquement changerait légèrement.
\item \textbf{Hypothèse d'uniformité} : pour calculer la moyenne d'une série en classes, on remplace chaque classe par son centre, ce qui revient à supposer que les valeurs s'y répartissent uniformément. Ce n'est qu'une \emph{approximation} : la moyenne obtenue est une valeur approchée de la vraie moyenne des données brutes.
\end{enumerate}
Écrire ces limites dans ta conclusion n'affaiblit pas ton travail : au contraire, c'est la marque d'un esprit scientifique rigoureux, et c'est apprécié par les correcteurs.
\end{attention}

\begin{exoresolu}[Rédiger la conclusion d'une enquête]
\textbf{Énoncé.} Les 36 élèves d'une classe de Première A ont mesuré la hauteur (en cm) de 36 plants de maïs du jardin scolaire, trois semaines après le semis. Les résultats, regroupés en classes d'amplitude 5 cm, donnent : moyenne $\bar{x} = 22$ cm, médiane lue sur le polygone cumulé $Me \approx 21$ cm, écart-type $\sigma = 6$ cm. Rédige la partie « interprétation et conclusion » du rapport d'enquête, en mentionnant les limites de l'étude.
\tcblower
\textbf{Solution détaillée.}
\emph{Interprétation.} Trois semaines après le semis, les plants mesurent en moyenne 22 cm, et la moitié d'entre eux mesure moins de 21 cm. La moyenne (22 cm) est très proche de la médiane (21 cm) : la distribution des hauteurs est à peu près symétrique, il n'y a pas de groupe de plants anormalement grands ou petits. L'écart-type de 6 cm, soit environ $27\,\%$ de la moyenne, indique une dispersion modérée : la majorité des plants a une hauteur comprise entre $\bar{x}-\sigma = 16$ cm et $\bar{x}+\sigma = 28$ cm. La croissance du semis est donc globalement régulière, avec toutefois quelques plants en retard (classe des plus petits) qu'un apport d'eau ou d'engrais ciblé pourrait aider.

\emph{Conclusion et limites.} Le semis du jardin scolaire présente une croissance homogène et satisfaisante. Cependant, cette conclusion porte sur 36 plants seulement, issus d'une seule parcelle : on ne peut pas la généraliser à toutes les parcelles de la région. De plus, le regroupement en classes de 5 cm et l'utilisation des centres de classes font de la moyenne une valeur approchée ; un autre découpage aurait pu donner des résultats légèrement différents. Une nouvelle mesure dans deux semaines permettra de suivre l'évolution.
\end{exoresolu}

\begin{exercice}[À toi de jouer : enquête sur les forfaits téléphoniques]
Les 32 élèves d'une classe ont relevé leur dépense mensuelle en forfaits internet (MTN ou Orange), en centaines de FCFA. Les résultats regroupés donnent : $\bar{x} = 25$ (soit \num{2500} FCFA), $Me = 20$ (soit \num{2000} FCFA), $\sigma = 14$.
\begin{enumerate}
\item Compare la moyenne et la médiane : que révèle leur écart sur la forme de la distribution ?
\item Interprète l'écart-type en langage courant.
\item Rédige en cinq ou six lignes la conclusion du rapport, en précisant la population, le caractère et deux limites de l'étude.
\end{enumerate}
\end{exercice}

\begin{corrige}[Exercice : enquête sur les forfaits téléphoniques]
\begin{enumerate}
\item On a $\bar{x} = 25 > Me = 20$ : la moyenne est nettement supérieure à la médiane. Quelques élèves dépensent beaucoup (gros forfaits) et tirent la moyenne vers le haut : la distribution est asymétrique, étalée vers les grandes valeurs. La moitié des élèves dépense en fait moins de \num{2000} FCFA par mois.
\item $\sigma = 14$ centaines de FCFA, soit \num{1400} FCFA, ce qui représente $56\,\%$ de la moyenne : la dispersion est forte, les dépenses sont très inégales d'un élève à l'autre. La majorité des dépenses se situe approximativement entre $25-14 = 11$ et $25+14 = 39$ centaines, soit entre \num{1100} et \num{3900} FCFA.
\item \emph{Conclusion possible.} « La population étudiée est l'ensemble des 32 élèves de la classe ; le caractère est la dépense mensuelle en forfaits internet, en FCFA, recueillie par questionnaire. La dépense moyenne est de \num{2500} FCFA, mais la médiane de \num{2000} FCFA montre que cette moyenne est gonflée par quelques gros consommateurs ; la forte dispersion ($\sigma = \num{1400}$ FCFA) traduit de grandes inégalités de dépense. Ces résultats ne concernent que cette classe (effectif réduit de 32 élèves) et ne peuvent être généralisés à tous les lycéens de la ville ; de plus, le regroupement en classes rend la moyenne approximative. »
\end{enumerate}
\end{corrige}

\begin{aretenir}[L'essentiel de la section]
\begin{itemize}
\item Un rapport d'enquête suit cinq étapes : introduction, méthode (population, caractère, collecte), résultats, interprétation, conclusion.
\item On croise toujours les trois paramètres : $\bar{x}$ et $Me$ pour la forme (symétrie si $\bar{x} \approx Me$ ; asymétrie sinon), $\sigma$ pour la concentration.
\item L'histogramme décrit la \emph{distribution} ; le polygone cumulé sert à lire la \emph{médiane} et les seuils.
\item Toute conclusion doit citer ses limites : effectif réduit, choix arbitraire des classes, approximation par les centres de classes, et interdiction de généraliser au-delà de la population réellement étudiée.
\end{itemize}
\end{aretenir}

\coursec{Compléments utiles pour la Terminale (optionnel)}

Dans la section précédente, tu as appris à interpréter et à communiquer les résultats d'une enquête : moyenne, médiane, écart-type, histogramme. Ces outils suffisent largement pour le programme de Première A. Mais en Terminale, tu rencontreras des situations plus fines : comparer deux séries qui n'ont pas la même unité, résumer une série entière en un seul dessin, ou encore étudier le lien entre \emph{deux} caractères à la fois (par exemple le prix et la quantité vendue). Cette section optionnelle te donne une longueur d'avance : les notions qui suivent prolongent naturellement ce que tu sais déjà.

\begin{attention}[Hors programme de Première A]
Les notions de cette section (variance corrigée, quartiles, boîte à moustaches, coefficient de variation, covariance, corrélation) \textbf{ne sont pas exigibles au Probatoire A}. Elles sont présentées ici pour préparer sereinement la transition vers la Terminale et pour enrichir ta culture statistique. Aucune n'est à restituer en devoir de Première, sauf demande explicite de ton professeur.
\end{attention}

\courssub{La variance corrigée : pourquoi diviser par $N-1$ ?}

\textbf{Intuition.} Jusqu'ici, tu as calculé la variance d'une \emph{population entière} : les 40 élèves d'une classe, les 50 clients d'une boutique. Mais dans la vraie vie, on ne peut presque jamais interroger toute la population : on ne va pas peser tous les sacs de maïs du Cameroun ! On prélève alors un \emph{échantillon} (par exemple 30 sacs) et on espère que la variance de l'échantillon donne une bonne idée de la variance de la population totale. Or un échantillon, par hasard, tend à être un peu moins dispersé que la population : les valeurs extrêmes y sont souvent absentes. La variance « classique » de l'échantillon \emph{sous-estime} donc, en moyenne, la vraie variance. Pour corriger ce défaut, les statisticiens divisent par $N-1$ au lieu de $N$.

\begin{definition}[Variance corrigée d'un échantillon]
Soit un échantillon de $N$ valeurs $x_1, x_2, \ldots, x_N$, de moyenne $\bar{x}$. La \cle{variance corrigée} (ou variance d'échantillon, dite \emph{non biaisée}) est :
\[
s^2 = \frac{1}{N-1}\sum_{i=1}^{N} (x_i - \bar{x})^2.
\]
L'\cle{écart-type corrigé} est $s = \sqrt{s^2}$. On dit que $s^2$ est un \cle{estimateur non biaisé} de la variance de la population : en moyenne, sur de nombreux échantillons, il ne surestime ni ne sous-estime la vraie variance.
\end{definition}

\begin{exemplebox}[Variance classique ou corrigée ?]
On relève les masses (en kg) de $N = 5$ sacs d'oignons au marché de Mokolo : $48$ ; $50$ ; $51$ ; $49$ ; $52$. La moyenne est $\bar{x} = 50$ kg. Les écarts à la moyenne sont $-2$ ; $0$ ; $1$ ; $-1$ ; $2$, de carrés $4, 0, 1, 1, 4$, de somme $10$.
\begin{itemize}
\item Variance « population » : $V = \dfrac{10}{5} = 2$, soit $\sigma = \sqrt{2} \approx 1,41$ kg.
\item Variance corrigée : $s^2 = \dfrac{10}{5-1} = 2,5$, soit $s \approx 1,58$ kg.
\end{itemize}
Si ces 5 sacs sont un \emph{échantillon} destiné à estimer la dispersion de toute la récolte, on retient $s^2 = 2,5$.
\end{exemplebox}

\begin{attention}[Quelle variance au Probatoire ?]
En Première A, on divise \textbf{toujours par $N$} (variance de la série étudiée, considérée comme complète). La division par $N-1$ ne concerne que l'estimation à partir d'un échantillon, notion de Terminale. Ne mélange pas les deux dans un devoir !
\end{attention}

\courssub{Quartiles, déciles, percentiles et boîte à moustaches}

\textbf{Intuition.} La médiane partage une série ordonnée en deux moitiés. Pourquoi s'arrêter là ? En partageant chaque moitié à son tour, on obtient quatre quarts : c'est l'idée des \emph{quartiles}. Trois nombres ($Q_1$, $Q_2$, $Q_3$) résument alors toute la dispersion, et un dessin très compact — la \emph{boîte à moustaches} — les visualise d'un coup d'œil.

\begin{definition}[Quartiles et écart interquartile]
Soit une série statistique ordonnée par valeurs croissantes, d'effectif total $N$.
\begin{itemize}
\item Le \cle{premier quartile} $Q_1$ est la plus petite valeur de la série telle qu'au moins $25\,\%$ des valeurs lui soient inférieures ou égales.
\item Le \cle{deuxième quartile} $Q_2$ n'est autre que la \cle{médiane} $Me$ (au moins $50\,\%$ des valeurs).
\item Le \cle{troisième quartile} $Q_3$ est la plus petite valeur telle qu'au moins $75\,\%$ des valeurs lui soient inférieures ou égales.
\end{itemize}
L'\cle{écart interquartile} est $Q_3 - Q_1$ : il mesure l'étalement de la « moitié centrale » des données. De même, les \cle{déciles} $D_1, D_2, \ldots, D_9$ partagent la série en dix dixièmes, et les \cle{percentiles} en cent centièmes (le $25$\textsuperscript{e} percentile est $Q_1$, le $50$\textsuperscript{e} est la médiane).
\end{definition}

\begin{definition}[Boîte à moustaches]
La \cle{boîte à moustaches} (ou \emph{box-plot}) d'une série est un diagramme formé de :
\begin{itemize}
\item un \textbf{rectangle} (la « boîte ») allant de $Q_1$ à $Q_3$, traversé par un trait à la médiane ;
\item deux segments (les « moustaches ») allant de la boîte jusqu'au minimum et au maximum de la série (ou jusqu'à des bornes excluant les valeurs aberrantes, selon les conventions).
\end{itemize}
\end{definition}

\begin{exemplebox}[Calcul des quartiles]
Les ventes quotidiennes de régimes de plantains d'un vendeur du marché Central de Yaoundé, sur 11 jours (série ordonnée) :
\[
8 \;;\; 10 \;;\; 12 \;;\; 13 \;;\; 15 \;;\; 16 \;;\; 17 \;;\; 19 \;;\; 20 \;;\; 22 \;;\; 25.
\]
\begin{itemize}
\item Médiane : $6$\textsuperscript{e} valeur, $Q_2 = Me = 16$.
\item $Q_1$ : médiane de la première moitié $\{8, 10, 12, 13, 15\}$, soit $Q_1 = 12$.
\item $Q_3$ : médiane de la seconde moitié $\{17, 19, 20, 22, 25\}$, soit $Q_3 = 20$.
\item Écart interquartile : $Q_3 - Q_1 = 20 - 12 = 8$ régimes.
\end{itemize}
La moitié centrale des journées voit des ventes comprises entre 12 et 20 régimes.
\end{exemplebox}

\begin{popfigure}
\begin{tikzpicture}[x=0.32cm,y=0.5cm]
% Axe
\draw[->,popdark] (5,0) -- (28,0) node[below left,font=\small]{ventes (régimes)};
\foreach \x in {6,8,...,26} \draw[popdark] (\x,0.12) -- (\x,-0.12) node[below,font=\scriptsize]{\x};
% Histogramme (barres)
\fill[popblueL,draw=popblue] (7,0) rectangle (9,1);
\fill[popblueL,draw=popblue] (9,0) rectangle (11,1);
\fill[popblueL,draw=popblue] (11,0) rectangle (13,2);
\fill[popblueL,draw=popblue] (13,0) rectangle (15,2);
\fill[popblueL,draw=popblue] (15,0) rectangle (17,2);
\fill[popblueL,draw=popblue] (17,0) rectangle (19,1);
\fill[popblueL,draw=popblue] (19,0) rectangle (21,2);
\fill[popblueL,draw=popblue] (21,0) rectangle (23,1);
\fill[popblueL,draw=popblue] (23,0) rectangle (25,0);
\fill[popblueL,draw=popblue] (25,0) rectangle (27,1);
% Box-plot au-dessus
\draw[poporange,very thick] (8,4.2) -- (8,5.8);   % min
\draw[poporange,very thick] (25,4.2) -- (25,5.8);  % max
\draw[poporange,very thick] (8,5) -- (12,5);       % moustache gauche
\draw[poporange,very thick] (20,5) -- (25,5);      % moustache droite
\fill[poporangeL,draw=poporange,very thick] (12,4.2) rectangle (20,5.8); % boîte
\draw[popink,very thick] (16,4.2) -- (16,5.8);     % médiane
% Annotations
\node[above,font=\scriptsize,poporange] at (8,5.8){min $= 8$};
\node[above,font=\scriptsize,poporange] at (12,5.8){$Q_1 = 12$};
\node[above,font=\scriptsize,popink] at (16,5.8){$Me = 16$};
\node[above,font=\scriptsize,poporange] at (20,5.8){$Q_3 = 20$};
\node[above,font=\scriptsize,poporange] at (25,5.8){max $= 25$};
\end{tikzpicture}
\legende{Boîte à moustaches superposée à l'histogramme des ventes quotidiennes de plantains : la boîte concentre 50 \% des journées, entre $Q_1 = 12$ et $Q_3 = 20$ régimes.}
\end{popfigure}

\begin{savaistu}[John Tukey et l'analyse exploratoire]
La boîte à moustaches a été popularisée en 1970 par le statisticien américain \emph{John Tukey}, pionnier de « l'analyse exploratoire des données ». Son idée : avant tout calcul savant, un bon dessin doit permettre de \emph{voir} la médiane, la dispersion et les valeurs extrêmes en une seconde. Tukey est aussi l'inventeur du mot « bit » et un père de la transformée de Fourier rapide, utilisée aujourd'hui dans tous les téléphones portables... y compris ceux des réseaux MTN et Orange !
\end{savaistu}

\begin{attention}[Conventions variables]
Selon les manuels et les logiciels, $Q_1$ et $Q_3$ peuvent être définis par des méthodes légèrement différentes (rang $\frac{N}{4}$, médiane des moitiés, interpolation). Les valeurs obtenues peuvent différer d'un cran pour les petites séries : ce n'est pas une erreur, mais une différence de convention. En Terminale, applique la convention de ton professeur et \textbf{annonce-la clairement}.
\end{attention}

\courssub{Le coefficient de variation : comparer des dispersions d'échelles différentes}

\textbf{Intuition.} Un grossiste de Douala vend du riz (prix moyen $\bar{x}_1 = 350$ FCFA le kg, écart-type $\sigma_1 = 20$ FCFA) et des sacs de ciment (prix moyen $\bar{x}_2 = 5\,500$ FCFA, écart-type $\sigma_2 = 200$ FCFA). Quel produit a les prix les plus dispersés ? Comparer $20$ et $200$ FCFA n'a pas de sens : les ordres de grandeur sont différents. Il faut \emph{relativiser} l'écart-type par la moyenne.

\begin{definition}[Coefficient de variation]
Le \cle{coefficient de variation} d'une série de moyenne $\bar{x} \neq 0$ et d'écart-type $\sigma$ est :
\[
CV = \frac{\sigma}{\bar{x}}, \qquad \text{souvent exprimé en pourcentage : } CV = \frac{\sigma}{\bar{x}} \times 100\,\%.
\]
\end{definition}

\begin{propriete}[Dispersion relative]
Le coefficient de variation est un nombre \textbf{sans unité} : il ne dépend ni de l'unité de mesure, ni de l'ordre de grandeur de la série. Il permet donc de comparer la \emph{dispersion relative} de séries d'échelles ou de natures différentes : plus $CV$ est grand, plus la série est dispersée \emph{relativement à sa moyenne}. (Non exigible au Probatoire A.)
\end{propriete}

\begin{exemplebox}[Riz ou ciment : lequel fluctue le plus ?]
Avec les données du grossiste :
\[
CV_{\text{riz}} = \frac{20}{350} \approx 0,057 = 5,7\,\% \qquad ; \qquad CV_{\text{ciment}} = \frac{200}{5\,500} \approx 0,036 = 3,6\,\%.
\]
Conclusion : bien que l'écart-type du ciment soit dix fois plus grand, c'est le prix du \textbf{riz} qui est relativement le plus instable ($5,7\,\% > 3,6\,\%$).
\end{exemplebox}

\begin{attention}[Moyenne nulle ou négative]
Le coefficient de variation n'a de sens que si $\bar{x} > 0$. Pour une série de moyenne nulle (écarts de température, gains et pertes) ou négative, le quotient $\sigma/\bar{x}$ est infini ou absurde : on ne l'utilise pas.
\end{attention}

\courssub{Covariance et corrélation linéaire : étudier deux caractères à la fois}

\textbf{Intuition.} Jusqu'ici, chaque individu était décrit par \emph{un seul} caractère. Mais souvent, on en observe \emph{deux} simultanément : les heures de révision $x$ et la note au Probatoire blanc $y$ de chaque élève. Question naturelle : quand $x$ augmente, $y$ a-t-il tendance à augmenter aussi ? La \emph{covariance} mesure cette « tendance à varier ensemble ».

\begin{definition}[Covariance]
Pour une série double $(x_i, y_i)$, $1 \le i \le N$, de moyennes $\bar{x}$ et $\bar{y}$, la \cle{covariance} est :
\[
\mathrm{Cov}(x, y) = \frac{1}{N}\sum_{i=1}^{N} (x_i - \bar{x})(y_i - \bar{y}) = \frac{1}{N}\sum_{i=1}^{N} x_i y_i - \bar{x}\,\bar{y}.
\]
Si $\mathrm{Cov}(x,y) > 0$, les deux caractères ont tendance à varier dans le même sens ; si $\mathrm{Cov}(x,y) < 0$, en sens contraires.
\end{definition}

La covariance dépend des unités : pour obtenir une mesure universelle, on la normalise par les écarts-types.

\begin{definition}[Coefficient de corrélation linéaire]
Le \cle{coefficient de corrélation linéaire} de $x$ et $y$ est :
\[
r = \frac{\mathrm{Cov}(x, y)}{\sigma_x \, \sigma_y}, \qquad \text{et l'on a toujours } -1 \le r \le 1.
\]
\begin{itemize}
\item $r$ proche de $1$ : forte dépendance croissante (les points du nuage s'alignent sur une droite montante) ;
\item $r$ proche de $-1$ : forte dépendance décroissante ;
\item $r$ proche de $0$ : absence de relation \emph{linéaire} entre $x$ et $y$.
\end{itemize}
\end{definition}

\begin{exoresolu}[Révisions et notes au Probatoire blanc]
Pour 5 élèves d'une classe de Première A à Bafoussam, on note $x$ les heures de révision hebdomadaires et $y$ la note sur 20 :
\begin{center}
\begin{tabular}{|c|ccccc|}\hline
\rowcolor{popblueL} Élève & 1 & 2 & 3 & 4 & 5 \\ \hline
$x_i$ (heures) & 2 & 4 & 6 & 8 & 10 \\ \hline
$y_i$ (note) & 8 & 10 & 11 & 14 & 17 \\ \hline
\end{tabular}
\end{center}
Calculer le coefficient de corrélation linéaire et commenter.
\tcblower
\textbf{Étape 1 : moyennes.} $\bar{x} = \dfrac{2+4+6+8+10}{5} = 6$ et $\bar{y} = \dfrac{8+10+11+14+17}{5} = 12$.

\textbf{Étape 2 : covariance.} On calcule les produits $x_i y_i$ : $16$ ; $40$ ; $66$ ; $112$ ; $170$, de somme $404$. Donc :
\[
\mathrm{Cov}(x,y) = \frac{404}{5} - 6 \times 12 = 80,8 - 72 = 8,8.
\]

\textbf{Étape 3 : écarts-types.}
\begin{align*}
\sigma_x^2 &= \frac{4^2 + 2^2 + 0^2 + 2^2 + 4^2}{5} = \frac{16+4+0+4+16}{5} = 8, \quad \sigma_x = \sqrt{8} \approx 2,83 ;\\
\sigma_y^2 &= \frac{(-4)^2 + (-2)^2 + (-1)^2 + 2^2 + 5^2}{5} = \frac{16+4+1+4+25}{5} = 10, \quad \sigma_y = \sqrt{10} \approx 3,16.
\end{align*}

\textbf{Étape 4 : corrélation.}
\[
r = \frac{8,8}{\sqrt{8}\times\sqrt{10}} = \frac{8,8}{\sqrt{80}} \approx \frac{8,8}{8,94} \approx 0,98.
\]
\textbf{Interprétation :} $r \approx 0,98$ est très proche de $1$ : il existe une forte relation croissante entre le temps de révision et la note. Les points $(x_i, y_i)$ sont presque alignés sur une droite montante. En Terminale, tu apprendras à tracer cette droite, appelée \emph{droite d'ajustement} ou \emph{droite de régression}.
\end{exoresolu}

\begin{attention}[Corrélation n'est pas causalité]
Un coefficient $r$ proche de $1$ ne prouve pas que $x$ \emph{provoque} $y$ ! Si les ventes de moustiquaires et celles de ventilateurs sont corrélées, c'est sans doute la saison sèche qui agit sur les deux. De même, $r \approx 0$ n'exclut pas une relation \emph{non linéaire} (par exemple en arc de parabole). Retiens : la corrélation mesure uniquement la force d'une relation \textbf{linéaire}.
\end{attention}

\begin{aretenir}[Ce qu'il faut retenir de cette ouverture]
\begin{itemize}
\item Variance corrigée d'un échantillon : $s^2 = \dfrac{1}{N-1}\sum (x_i - \bar{x})^2$ (division par $N-1$).
\item Quartiles $Q_1$, $Q_2 = Me$, $Q_3$ ; écart interquartile $Q_3 - Q_1$ ; la boîte à moustaches résume min, $Q_1$, $Me$, $Q_3$, max en un seul dessin.
\item Coefficient de variation $CV = \dfrac{\sigma}{\bar{x}}$ : sans unité, il compare les dispersions relatives.
\item Covariance $\mathrm{Cov}(x,y) = \dfrac{1}{N}\sum x_i y_i - \bar{x}\bar{y}$ et corrélation $r = \dfrac{\mathrm{Cov}(x,y)}{\sigma_x \sigma_y} \in [-1\,;\,1]$.
\end{itemize}
Ces quatre outils seront tes compagnons de Terminale : les avoir rencontrés dès maintenant, c'est déjà les avoir à moitié compris.
\end{aretenir}

\begin{exercice}[Pour aller plus loin]
Une coopérative de café de l'Ouest prélève 6 sacs : masses (kg) $60$ ; $62$ ; $58$ ; $61$ ; $59$ ; $60$.
\begin{enumerate}
\item Calculer la moyenne, puis la variance corrigée $s^2$ et l'écart-type corrigé $s$.
\item Déterminer $Q_1$, $Me$, $Q_3$ et l'écart interquartile.
\item Une autre coopérative a $\bar{x}' = 60$ kg et $\sigma' = 3$ kg. Comparer les dispersions relatives à l'aide du coefficient de variation.
\end{enumerate}
\end{exercice}

\begin{corrige}[Exercice : coopérative de café]
\textbf{1.} $\bar{x} = \dfrac{60+62+58+61+59+60}{6} = \dfrac{360}{6} = 60$ kg. Les écarts sont $0, 2, -2, 1, -1, 0$, de carrés $0, 4, 4, 1, 1, 0$, de somme $10$. Variance corrigée : $s^2 = \dfrac{10}{6-1} = 2$, d'où $s = \sqrt{2} \approx 1,41$ kg.

\textbf{2.} Série ordonnée : $58$ ; $59$ ; $60$ ; $60$ ; $61$ ; $62$. Médiane : $Me = \dfrac{60+60}{2} = 60$ kg. Première moitié $\{58, 59, 60\}$ : $Q_1 = 59$ kg. Seconde moitié $\{60, 61, 62\}$ : $Q_3 = 61$ kg. Écart interquartile : $Q_3 - Q_1 = 2$ kg.

\textbf{3.} Première coopérative : $CV = \dfrac{1,41}{60} \approx 2,4\,\%$. Seconde : $CV' = \dfrac{3}{60} = 5\,\%$. La seconde coopérative a des masses relativement plus dispersées : le remplissage de la première est plus régulier.
\end{corrige}

\newpage

\coursec{Activité d'intégration}

\coursec{Activité d'intégration : Enquête statistique sur les temps de trajet domicile--lycée}

\courssub{Objectifs de l'activité}

À l'issue de cette activité d'intégration, tu seras capable de mobiliser \emph{simultanément} l'ensemble des savoir-faire de la leçon dans une situation authentique d'enquête statistique. Concrètement, tu devras :

\begin{itemize}
\item organiser une collecte de données réelles et regrouper une série brute en \cle{classes d'égales amplitudes} ;
\item construire un tableau complet des effectifs, des fréquences et des \cle{effectifs cumulés croissants et décroissants} ;
\item calculer la \cle{moyenne}, la \cle{médiane}, la \cle{variance} et l'\cle{écart-type} d'une série regroupée en classes ;
\item construire un \cle{histogramme} et les deux \cle{polygones des effectifs cumulés} ;
\item rédiger un rapport court interprétant les paramètres obtenus et le présenter oralement.
\end{itemize}

Cette activité reproduit exactement la démarche d'un statisticien : collecter, organiser, calculer, représenter, interpréter, communiquer.

\courssub{Situation}

Le proviseur du lycée bilingue de Nkongsamba s'inquiète des retards répétés constatés lors du premier cours de la journée. Le conseil d'administration envisage d'ajuster l'heure de début des cours, actuellement fixée à 7 h 30. Avant toute décision, le proviseur mandate la classe de Première A pour mener une enquête statistique sur le \textbf{temps de trajet domicile--lycée} des élèves (marche à pied, moto-taxi « benskin », taxi collectif ou car).

Chaque élève de la classe a noté anonymement, sur une fiche, la durée en minutes de son trajet le matin même. Voici les $30$ données brutes recueillies (en minutes) :

\begin{center}
\begin{tabular}{|c|c|c|c|c|c|c|c|c|c|}
\hline
\rowcolor{popblueL}
5 & 8 & 12 & 15 & 18 & 20 & 22 & 25 & 25 & 27 \\
\hline
28 & 30 & 32 & 35 & 38 & 40 & 42 & 45 & 10 & 14 \\
\hline
16 & 22 & 24 & 26 & 28 & 33 & 36 & 44 & 48 & 20 \\
\hline
\end{tabular}
\end{center}

Le proviseur pose trois questions à la classe :
\begin{enumerate}
\item Quel est le temps de trajet \emph{moyen} d'un élève du lycée ?
\item À partir de quelle durée peut-on dire qu'un élève sur deux met \emph{plus de temps} que les autres ?
\item Les temps de trajet sont-ils \emph{homogènes} (regroupés autour de la moyenne) ou très dispersés ? Cela justifie-t-il de décaler l'heure des cours ?
\end{enumerate}

\courssub{Consignes}

\textbf{Partie A : Organisation des données}

\begin{enumerate}
\item Justifier que le caractère étudié est quantitatif continu et proposer un regroupement en $5$ classes d'amplitude $10$, en commençant par $[0\,;\,10[$.
\item Construire le tableau statistique complet : classes, centres $c_i$, effectifs $n_i$, fréquences $f_i$, effectifs cumulés croissants (ECC) et effectifs cumulés décroissants (ECD).
\end{enumerate}

\textbf{Partie B : Caractéristiques de position et de dispersion}

\begin{enumerate}[resume]
\item Calculer la moyenne $\bar{x}$ de la série à l'aide des centres de classes.
\item Déterminer la classe médiane, puis calculer une valeur approchée de la médiane $M_e$ par interpolation linéaire.
\item Calculer la variance $V$ puis l'écart-type $\sigma$ de la série.
\item Comparer la moyenne et la médiane. Que peut-on en déduire sur la forme de la distribution ?
\end{enumerate}

\textbf{Partie C : Représentations graphiques}

\begin{enumerate}[resume]
\item Construire l'histogramme des effectifs (sur papier millimétré : 1 cm pour 10 min en abscisse, 1 cm pour 2 élèves en ordonnée).
\item Construire, sur un même graphique, le polygone des ECC et le polygone des ECD. Retrouver graphiquement une valeur approchée de la médiane.
\end{enumerate}

\textbf{Partie D : Interprétation et communication}

\begin{enumerate}[resume]
\item En groupes de quatre, rédiger un rapport d'une page répondant aux trois questions du proviseur, en citant tous les paramètres calculés.
\item Chaque groupe présente ses conclusions en 3 minutes, puis la classe vote une recommandation argumentée (maintenir ou décaler l'heure de 7 h 30).
\end{enumerate}

\courssub{Ressources à mobiliser}

\begin{aretenir}[Boîte à outils de la leçon]
\begin{itemize}
\item \textbf{Centre d'une classe} $[a\,;\,b[$ : $c_i=\dfrac{a+b}{2}$.
\item \textbf{Moyenne} (série en classes) : $\bar{x}=\dfrac{\sum n_i c_i}{N}$.
\item \textbf{Variance} : $V=\dfrac{\sum n_i(c_i-\bar{x})^2}{N}=\dfrac{\sum n_i c_i^2}{N}-\bar{x}^2$ (formule de König--Huygens) ; \textbf{écart-type} : $\sigma=\sqrt{V}$.
\item \textbf{Médiane} : valeur qui partage la série en deux effectifs égaux ; pour une série en classes, on repère la classe médiane avec les ECC puis on interpole linéairement.
\item \textbf{Polygones cumulés} : le polygone des ECC joint les points $(b_i\,;\,\mathrm{ECC}_i)$ où $b_i$ est la borne supérieure ; le polygone des ECD joint les points $(a_i\,;\,\mathrm{ECD}_i)$ où $a_i$ est la borne inférieure ; l'abscisse du point d'intersection des deux polygones approche la médiane.
\end{itemize}
\end{aretenir}

\begin{methode}[Calcul de la médiane par interpolation linéaire (série en classes)]
\begin{enumerate}
\item Calculer $N/2$.
\item Repérer dans la colonne des ECC la première valeur $\ge N/2$ ; la classe correspondante est la \cle{classe médiane} $[a\,;\,b[$.
\item Noter $F_{\text{av}}$ l'ECC de la classe précédente (effectif cumulé avant la classe médiane) et $n_{\text{med}}$ l'effectif de la classe médiane.
\item Appliquer la formule : $M_e = a + \dfrac{b-a}{n_{\text{med}}}\left(\dfrac{N}{2} - F_{\text{av}}\right)$.
\end{enumerate}
\end{methode}

\courssub{Démarche et corrigé commenté}

\begin{exoresolu}[Enquête « Temps de trajet domicile--lycée » : corrigé complet]
Répondre à l'ensemble des consignes des parties A, B, C et D.

\tcblower

\textbf{Partie A. Organisation des données.}

\textbf{1.} Le temps de trajet peut prendre n'importe quelle valeur réelle dans un intervalle (on peut mesurer 12,3 min, 12,34 min, etc.) : c'est un caractère \cle{quantitatif continu}. L'étendue vaut $48-5=43$ minutes ; un regroupement en $5$ classes d'amplitude $10$ est donc adapté : $[0\,;\,10[$, $[10\,;\,20[$, $[20\,;\,30[$, $[30\,;\,40[$, $[40\,;\,50]$.

\textbf{2.} En dépouillant les $30$ fiches (chaque valeur est rangée dans sa classe, la borne inférieure étant incluse), on obtient :

\begin{center}
\begin{tabularx}{\linewidth}{|l|X|X|X|X|X|}
\hline
\rowcolor{popblueL}
Classe & Centre $c_i$ & $n_i$ & $f_i$ & ECC & ECD \\
\hline
$[0\,;\,10[$ & 5 & 2 & 0,067 & 2 & 30 \\
\hline
$[10\,;\,20[$ & 15 & 6 & 0,200 & 8 & 28 \\
\hline
$[20\,;\,30[$ & 25 & 11 & 0,367 & 19 & 22 \\
\hline
$[30\,;\,40[$ & 35 & 6 & 0,200 & 25 & 11 \\
\hline
$[40\,;\,50]$ & 45 & 5 & 0,167 & 30 & 5 \\
\hline
\rowcolor{popblueL}
Total & & 30 & 1 & & \\
\hline
\end{tabularx}
\end{center}

Vérification indispensable : la somme des effectifs vaut bien $N=30$ et la somme des fréquences vaut $1$ (aux arrondis près).

\textbf{Partie B. Position et dispersion.}

\textbf{3.} Moyenne :
\[
\bar{x}=\frac{2\times 5+6\times 15+11\times 25+6\times 35+5\times 45}{30}
=\frac{10+90+275+210+225}{30}=\frac{810}{30}=27.
\]
Le temps de trajet moyen est donc de $\bar{x}=27$ minutes.

\textbf{4.} On a $\dfrac{N}{2}=15$. Dans la colonne des ECC, $15$ se situe entre $8$ et $19$ : la \cle{classe médiane} est donc $[20\,;\,30[$. Par interpolation linéaire (formule de la méthode) :
\[
M_e = 20+10\times\frac{15-8}{19-8}=20+\frac{70}{11}\approx 26,4 \text{ min}.
\]
Environ la moitié des élèves met moins de $26,4$ minutes, l'autre moitié davantage.

\textbf{5.} Variance, avec les écarts au carré (on peut aussi utiliser la formule de König--Huygens) :
\begin{align*}
V&=\frac{2(5-27)^2+6(15-27)^2+11(25-27)^2+6(35-27)^2+5(45-27)^2}{30}\\
&=\frac{2\times 484+6\times 144+11\times 4+6\times 64+5\times 324}{30}\\
&=\frac{968+864+44+384+1620}{30}=\frac{3880}{30}\approx 129,3.
\end{align*}
D'où $\sigma=\sqrt{129,3}\approx 11,4$ minutes.

\textbf{6.} On observe $M_e\approx 26,4 < \bar{x}=27$ : la moyenne, sensible aux valeurs élevées (élèves venant des quartiers éloignés), est légèrement supérieure à la médiane. La distribution est donc \emph{légèrement étirée vers les longs trajets} (asymétrie positive). L'écart-type de $11,4$ min, soit environ $42\,\%$ de la moyenne, traduit une dispersion notable.

\textbf{Partie C. Représentations graphiques.}

\textbf{7.} Histogramme des effectifs :

\begin{popfigure}
\begin{center}
\begin{tikzpicture}[x=0.22cm,y=0.5cm]
\draw[->,thick] (0,0) -- (52,0) node[below left]{\scriptsize temps (min)};
\draw[->,thick] (0,0) -- (0,12.5) node[above left]{\scriptsize effectif};
\foreach \x in {0,10,20,30,40,50}{\draw (\x,0.15)--(\x,-0.15) node[below]{\scriptsize \x};}
\foreach \y in {2,4,6,8,10}{\draw (0.4,\y)--(-0.4,\y) node[left]{\scriptsize \y};}
\fill[popblue!60,draw=popdark] (0,0) rectangle (10,2);
\fill[popblue!60,draw=popdark] (10,0) rectangle (20,6);
\fill[popblue!60,draw=popdark] (20,0) rectangle (30,11);
\fill[popblue!60,draw=popdark] (30,0) rectangle (40,6);
\fill[popblue!60,draw=popdark] (40,0) rectangle (50,5);
\end{tikzpicture}
\end{center}
\legende{Histogramme des effectifs : la classe modale $[20\,;\,30[$ concentre $11$ élèves.}
\end{popfigure}

\textbf{8.} Polygones cumulés : le polygone des ECC joint les points $(0\,;\,0)$, $(10\,;\,2)$, $(20\,;\,8)$, $(30\,;\,19)$, $(40\,;\,25)$, $(50\,;\,30)$ ; le polygone des ECD joint $(0\,;\,30)$, $(10\,;\,28)$, $(20\,;\,22)$, $(30\,;\,11)$, $(40\,;\,5)$, $(50\,;\,0)$.

\begin{popfigure}
\begin{center}
\begin{tikzpicture}[x=0.22cm,y=0.16cm]
\draw[->,thick] (0,0) -- (54,0) node[below left]{\scriptsize temps (min)};
\draw[->,thick] (0,0) -- (0,33) node[above left]{\scriptsize effectif cumulé};
\foreach \x in {0,10,20,30,40,50}{\draw (\x,0.6)--(\x,-0.6) node[below]{\scriptsize \x};}
\foreach \y in {5,10,15,20,25,30}{\draw (0.8,\y)--(-0.8,\y) node[left]{\scriptsize \y};}
\draw[very thick,popblue] (0,0)--(10,2)--(20,8)--(30,19)--(40,25)--(50,30);
\draw[very thick,poporange] (0,30)--(10,28)--(20,22)--(30,11)--(40,5)--(50,0);
\foreach \p in {(10,2),(20,8),(30,19),(40,25)}{\fill[popblue] \p circle (0.35);}
\foreach \p in {(10,28),(20,22),(30,11),(40,5)}{\fill[poporange] \p circle (0.35);}
\draw[dashed,popgreen,thick] (26.4,0) -- (26.4,15) -- (0,15);
\fill[popgreen] (26.4,15) circle (0.45);
\node[popgreen,below] at (26.4,0) {\scriptsize $M_e\approx 26{,}4$};
\node[popblue] at (44,29) {\scriptsize ECC};
\node[poporange] at (44,3.2) {\scriptsize ECD};
\end{tikzpicture}
\end{center}
\legende{Les deux polygones se coupent au point d'abscisse $M_e\approx 26{,}4$ min : on retrouve graphiquement la médiane calculée.}
\end{popfigure}

Le point d'intersection des deux polygones a pour ordonnée $15=\dfrac{N}{2}$ et pour abscisse environ $26,4$ : la lecture graphique confirme le calcul par interpolation.

\textbf{Partie D. Interprétation et communication (éléments attendus pour le rapport).}

Le rapport doit contenir : le tableau statistique, les paramètres ($\bar{x}=27$ min ; $M_e\approx 26,4$ min ; $\sigma\approx 11,4$ min), les graphiques, et une conclusion argumentée. Exemple de réponse aux questions du proviseur :
\begin{itemize}
\item \textbf{Temps moyen :} Un élève met en moyenne $27$ minutes pour venir au lycée.
\item \textbf{Médiane :} Un élève sur deux met plus de $26,4$ minutes (la médiane).
\item \textbf{Dispersion :} L'écart-type $\sigma\approx 11,4$ min montre une forte hétérogénéité. L'intervalle $[\bar{x}-\sigma\,;\,\bar{x}+\sigma]=[15,6\,;\,38,4]$ contient, d'après les données brutes, $19$ élèves sur $30$ (soit $63\,\%$), confirmant une dispersion importante. Les extrêmes vont de $5$ min (marche à pied, quartier proche) à $48$ min (benskin depuis la périphérie).
\item \textbf{Recommandation argumentée :} Décaler le début des cours à 7 h 45 ne résoudrait pas le problème des élèves les plus éloignés (ils arriveraient encore en retard). Mieux vaut négocier avec le transporteur local (CAMTEL/SOTUC) un circuit de ramassage scolaire desservant les quartiers situés à plus de $40$ minutes, ou tolérer un accueil progressif entre 7 h 30 et 8 h.
\end{itemize}
\end{exoresolu}

\begin{attention}[Pièges fréquents dans ce type d'enquête]
\begin{itemize}
\item Ne confonds pas \emph{borne de classe} et \emph{centre de classe} : tous les calculs (moyenne, variance) utilisent les centres $c_i$.
\item Pour la médiane, n'oublie pas l'interpolation linéaire : répondre « la médiane est dans $[20\,;\,30[$ » est insuffisant au Probatoire.
\item Dans le polygone des ECC, on place les effectifs cumulés aux \emph{bornes supérieures} des classes ; aux bornes inférieures pour les ECD.
\item Vérifie toujours que $\sum n_i = N$ avant tout calcul : une erreur de dépouillement fausse toute la suite.
\item Pour l'intervalle $[\bar{x}-\sigma\,;\,\bar{x}+\sigma]$, ne compte pas les effectifs de classes entières si les bornes coupent des classes ; utilise les données brutes ou une interpolation.
\end{itemize}
\end{attention}

\begin{savaistu}[La règle empirique (loi 68--95--99,7)]
Pour une distribution \emph{proche de la loi normale} (en forme de cloche symétrique), on sait que :
\begin{itemize}
\item environ $68\,\%$ des données sont dans $[\bar{x}-\sigma\,;\,\bar{x}+\sigma]$,
\item environ $95\,\%$ dans $[\bar{x}-2\sigma\,;\,\bar{x}+2\sigma]$,
\item environ $99,7\,\%$ dans $[\bar{x}-3\sigma\,;\,\bar{x}+3\sigma]$.
\end{itemize}
Ici, la distribution est asymétrique (queue à droite), donc la règle n'est pas exacte : on trouve $63\,\%$ au lieu de $68\,\%$. Cela rappelle qu'il faut toujours \emph{regarder la forme de l'histogramme} avant d'appliquer une règle théorique.
\end{savaistu}

\courssub{Bilan de l'activité}

Cette enquête a permis de mettre en évidence plusieurs idées essentielles de la leçon :

\begin{itemize}
\item \textbf{La statistique est une démarche complète} : aucun paramètre n'a de sens sans une collecte soigneuse et un regroupement judicieux des données en classes d'égales amplitudes.
\item \textbf{Moyenne et médiane se complètent} : ici, $\bar{x}=27$ et $M_e\approx 26,4$ sont proches, mais leur comparaison révèle l'étirement de la distribution vers les longs trajets, information invisible si l'on ne calcule qu'un seul paramètre.
\item \textbf{L'écart-type mesure l'homogénéité} : $\sigma\approx 11,4$ min montre que le « trajet moyen » ne décrit fidèlement personne ; les écarts individuels sont importants.
\item \textbf{Les graphiques confirment les calculs} : l'histogramme visualise la classe modale, et l'intersection des polygones cumulés redonne la médiane sans nouveau calcul.
\item \textbf{Les mathématiques éclairent la décision} : la recommandation finale au proviseur s'appuie sur des nombres, pas sur des impressions — c'est le cœur du savoir-être statistique : argumenter avec rigueur, honnêteté et esprit critique.
\end{itemize}

Tu es désormais prêt à mener ta propre mini-enquête (notes à un devoir, dépenses hebdomadaires en crédit téléphonique MTN/Orange, distances parcourues) et à en communiquer les résultats comme un véritable statisticien.

\newpage

\coursec{Méthodes et exercices résolus}

\coursec{Statistiques : Organisation et gestion des données}

\courssub{Regrouper une population en classes d'égales amplitudes}

\begin{methode}[Regrouper des données en classes d'égales amplitudes]
Pour regrouper une série de $N$ données continues (ou discrètes nombreuses) en classes d'égales amplitudes :
\begin{enumerate}
\item Repérer la valeur minimale $x_{\min}$ et la valeur maximale $x_{\max}$, puis calculer l'\cle{étendue} $e = x_{\max}-x_{\min}$.
\item Choisir le nombre de classes $k$ (souvent imposé par l'énoncé ; sinon entre 5 et 8).
\item Calculer l'\cle{amplitude} commune $a \approx \dfrac{e}{k}$, en l'arrondissant à une valeur simple (multiple de 5, 10, etc.).
\item Écrire les classes sous la forme $[b_i\,;\,b_{i+1}[$ : borne inférieure incluse, borne supérieure exclue (sauf éventuellement la dernière classe, fermée des deux côtés $[b_k\,;\,b_{k+1}]$).
\item Compter l'effectif $n_i$ de chaque classe en dépouillant les données une par une, puis \textbf{vérifier} que $\sum n_i = N$.
\item Compléter le tableau avec le centre de classe $c_i=\dfrac{b_i+b_{i+1}}{2}$, les fréquences $f_i=\dfrac{n_i}{N}$, les effectifs cumulés croissants (E.C.C.) et décroissants (E.C.D.), ainsi que les fréquences cumulées croissantes (F.C.C.) et décroissantes (F.C.D.).
\end{enumerate}
\end{methode}

\begin{exoresolu}[Regroupement en classes et tableaux cumulés]
Un grossiste de Mokolo a relevé les masses (en kg) de 40 sacs d'arachide :
\[ 42\,;\,45\,;\,47\,;\,50\,;\,52\,;\,53\,;\,55\,;\,55\,;\,56\,;\,57\,;\,58\,;\,58\,;\,59\,;\,60\,;\,60\,;\,61\,;\,61\,;\,62\,;\,62\,;\,63 \]
\[ 63\,;\,64\,;\,64\,;\,65\,;\,65\,;\,66\,;\,66\,;\,67\,;\,68\,;\,68\,;\,69\,;\,70\,;\,70\,;\,71\,;\,72\,;\,73\,;\,74\,;\,75\,;\,76\,;\,78 \]
Regrouper ces données en classes d'amplitude $8$ en partant de $40$, et dresser le tableau complet des effectifs, fréquences, effectifs et fréquences cumulés croissants et décroissants.
\tcblower
\textbf{Étape 1 : étendue.} $x_{\min}=42$, $x_{\max}=78$, donc $e = 78-42 = 36$.

\textbf{Étape 2 : classes d'amplitude 8 à partir de 40.}
\[ [40;48[\,;\quad [48;56[\,;\quad [56;64[\,;\quad [64;72[\,;\quad [72;80] \]

\textbf{Étape 3 : dépouillement.} On compte les données dans chaque classe :
\begin{itemize}
\item $[40;48[$ : 42, 45, 47 $\Rightarrow n_1=3$ ;
\item $[48;56[$ : 50, 52, 53, 55, 55 $\Rightarrow n_2=5$ ;
\item $[56;64[$ : 56, 57, 58, 58, 59, 60, 60, 61, 61, 62, 62, 63, 63 $\Rightarrow n_3=13$ ;
\item $[64;72[$ : 64, 64, 65, 65, 66, 66, 67, 68, 68, 69, 70, 70, 71 $\Rightarrow n_4=13$ ;
\item $[72;80]$ : 72, 73, 74, 75, 76, 78 $\Rightarrow n_5=6$.
\end{itemize}
Vérification : $3+5+13+13+6 = 40 = N$. \checkmark

\textbf{Étape 4 : tableau complet.}
\begin{center}
\begin{tabular}{|l|c|c|c|c|c|}
\hline
\rowcolor{popblueL} Classe & $[40;48[$ & $[48;56[$ & $[56;64[$ & $[64;72[$ & $[72;80]$ \\ \hline
Centre $c_i$ & 44 & 52 & 60 & 68 & 76 \\ \hline
Effectif $n_i$ & 3 & 5 & 13 & 13 & 6 \\ \hline
Fréquence $f_i$ & 0,075 & 0,125 & 0,325 & 0,325 & 0,15 \\ \hline
E.C.C. & 3 & 8 & 21 & 34 & 40 \\ \hline
E.C.D. & 40 & 37 & 32 & 19 & 6 \\ \hline
F.C.C. & 0,075 & 0,20 & 0,525 & 0,85 & 1 \\ \hline
F.C.D. & 1 & 0,925 & 0,80 & 0,475 & 0,15 \\ \hline
\end{tabular}
\end{center}
Les fréquences sont calculées par $f_i=\dfrac{n_i}{40}$ : par exemple $f_3=\dfrac{13}{40}=0,325$. On vérifie que la somme des fréquences vaut $1$, que le dernier E.C.C. vaut $N=40$, que le premier E.C.D. vaut $N=40$, et que la dernière F.C.C. vaut $1$ (ou $100\,\%$).
\end{exoresolu}

\courssub{Caractéristique de position : la moyenne}

\begin{methode}[Calculer une moyenne]
\begin{enumerate}
\item \textbf{Série à modalités simples} (valeurs $x_i$ d'effectifs $n_i$) :
\[ \bar{x} = \dfrac{\sum_{i} n_i x_i}{N} = \dfrac{n_1x_1+n_2x_2+\cdots+n_px_p}{N}. \]
\item \textbf{Série regroupée en classes} : on remplace chaque classe par son \cle{centre} $c_i$ (valeur approchée) :
\[ \bar{x} \approx \dfrac{\sum_{i} n_i c_i}{N}. \]
\item Construire une colonne $n_ix_i$ (ou $n_ic_i$) dans le tableau, faire la somme, puis diviser par $N$.
\item Conclure par une phrase avec l'\textbf{unité} et une valeur arrondie cohérente avec les données (généralement au dixième ou centième près).
\end{enumerate}
\textbf{Réflexe contrôle} : la moyenne doit se situer entre la plus petite et la plus grande valeur de la série (ou entre les bornes extrêmes des classes).
\end{methode}

\begin{exoresolu}[Moyenne d'une série regroupée en classes]
On reprend la répartition des 40 sacs d'arachide de l'exercice précédent. Calculer la masse moyenne d'un sac.
\tcblower
On utilise les centres de classes et on ajoute une colonne $n_ic_i$ :
\begin{center}
\begin{tabular}{|l|c|c|c|}
\hline
\rowcolor{popblueL} Classe & Centre $c_i$ & $n_i$ & $n_i c_i$ \\ \hline
$[40;48[$ & 44 & 3 & 132 \\ \hline
$[48;56[$ & 52 & 5 & 260 \\ \hline
$[56;64[$ & 60 & 13 & 780 \\ \hline
$[64;72[$ & 68 & 13 & 884 \\ \hline
$[72;80]$ & 76 & 6 & 456 \\ \hline
\rowcolor{poporangeL} Total & & 40 & 2512 \\ \hline
\end{tabular}
\end{center}
La masse moyenne est donc :
\[ \bar{x} = \dfrac{2512}{40} = 62,8. \]
\textbf{Contrôle} : $62,8 \in [42;78]$, le résultat est plausible. \checkmark

\textbf{Conclusion} : la masse moyenne d'un sac d'arachide est d'environ $62,8$ kg.
\end{exoresolu}

\courssub{Caractéristique de position : la médiane}

\begin{methode}[Déterminer la médiane]
\begin{enumerate}
\item \textbf{Série discrète} : ranger les $N$ valeurs dans l'ordre croissant.
\begin{itemize}
\item Si $N$ est impair, $Me$ est la valeur de rang $\dfrac{N+1}{2}$.
\item Si $N$ est pair, $Me$ est la demi-somme des valeurs de rangs $\dfrac{N}{2}$ et $\dfrac{N}{2}+1$.
\end{itemize}
\item \textbf{Série en classes} : repérer la \cle{classe médiane}, première classe dont l'effectif cumulé croissant atteint ou dépasse $\dfrac{N}{2}$, puis appliquer l'\cle{interpolation linéaire} (formule exigible au Probatoire) :
\[ Me = b_i + a \times \dfrac{\frac{N}{2}-N_{i-1}}{n_i}, \]
où $b_i$ est la borne inférieure de la classe médiane, $a$ son amplitude, $N_{i-1}$ l'effectif cumulé croissant de la classe précédente et $n_i$ l'effectif de la classe médiane.
\item \textbf{Graphiquement} : la médiane est l'abscisse du point du polygone des effectifs cumulés croissants d'ordonnée $\dfrac{N}{2}$. C'est aussi l'abscisse du point d'intersection des polygones des effectifs cumulés croissants et décroissants.
\end{enumerate}
\end{methode}

\begin{exoresolu}[Médiane par interpolation linéaire]
On reprend la série des 40 sacs d'arachide. Déterminer la médiane par interpolation linéaire, puis interpréter le résultat.
\tcblower
\textbf{Étape 1 : position de la médiane.} $N=40$, donc $\dfrac{N}{2}=20$. On cherche la première classe dont l'E.C.C. atteint 20.

\textbf{Étape 2 : classe médiane.} D'après le tableau des E.C.C. : 3, 8, \textbf{21}, 34, 40. Le premier E.C.C. supérieur ou égal à 20 est 21 : la classe médiane est $[56;64[$.

\textbf{Étape 3 : interpolation linéaire.} Ici $b_i=56$, $a=8$, $N_{i-1}=8$, $n_i=13$ :
\[ Me = 56 + 8\times\dfrac{20-8}{13} = 56 + 8\times\dfrac{12}{13} = 56 + \dfrac{96}{13} \approx 56 + 7,3846 \approx 63,4. \]

\textbf{Conclusion} : la médiane vaut environ $63,4$ kg. Cela signifie qu'environ la moitié des sacs pèse moins de $63,4$ kg et l'autre moitié plus de $63,4$ kg.
\end{exoresolu}

\courssub{Caractéristiques de dispersion : variance et écart-type}

\begin{methode}[Calculer variance et écart-type]
\begin{enumerate}
\item Calculer d'abord la moyenne $\bar{x}$.
\item \textbf{Formule de définition} (utile pour comprendre le sens) :
\[ V = \dfrac{1}{N}\sum_i n_i (x_i-\bar{x})^2 \qquad\text{puis}\qquad \sigma = \sqrt{V}. \]
\item \textbf{Formule de König-Huygens} (souvent plus rapide, moins de risques d'erreurs) :
\[ V = \dfrac{1}{N}\sum_i n_i x_i^{\,2} - \bar{x}^2. \]
\item Pour une série en classes, remplacer $x_i$ par le centre $c_i$ (valeurs approchées).
\item Ajouter les colonnes $n_ix_i^2$ (ou $n_i(x_i-\bar{x})^2$) au tableau pour organiser les calculs.
\item Conclure : $\sigma$ s'exprime dans la même unité que les données ; plus $\sigma$ est grand, plus la série est dispersée autour de sa moyenne.
\end{enumerate}
\end{methode}

\begin{exoresolu}[Variance et écart-type]
Les notes de mathématiques de 5 élèves de Première A au premier devoir sont : $8\,;\,10\,;\,11\,;\,14\,;\,17$. Calculer la moyenne, la variance et l'écart-type de cette série. Interpréter.
\tcblower
\textbf{Étape 1 : moyenne.}
\[ \bar{x} = \dfrac{8+10+11+14+17}{5} = \dfrac{60}{5} = 12. \]

\textbf{Étape 2 : variance par la définition.} On calcule les écarts à la moyenne :
\begin{center}
\begin{tabular}{|c|c|c|c|}
\hline
\rowcolor{popblueL} $x_i$ & $x_i-\bar{x}$ & $(x_i-\bar{x})^2$ \\ \hline
8 & $-4$ & 16 \\ \hline
10 & $-2$ & 4 \\ \hline
11 & $-1$ & 1 \\ \hline
14 & $2$ & 4 \\ \hline
17 & $5$ & 25 \\ \hline
\rowcolor{poporangeL} Total & 0 & 50 \\ \hline
\end{tabular}
\end{center}
\[ V = \dfrac{50}{5} = 10. \]
\textbf{Vérification par König-Huygens} : $\sum x_i^2 = 64+100+121+196+289 = 770$, donc
\[ V = \dfrac{770}{5} - 12^2 = 154 - 144 = 10. \checkmark \]

\textbf{Étape 3 : écart-type.}
\[ \sigma = \sqrt{10} \approx 3,16. \]

\textbf{Conclusion} : la moyenne de la classe est $12$ et l'écart-type vaut environ $3,2$ points : les notes sont assez dispersées autour de la moyenne (de 8 à 17).
\end{exoresolu}

\courssub{Représentations graphiques : histogramme et polygones des effectifs cumulés}

\begin{methode}[Histogramme et polygones cumulés]
\textbf{Histogramme} (classes d'égales amplitudes) :
\begin{enumerate}
\item En abscisse : les bornes des classes, à échelle régulière ; en ordonnée : les effectifs (ou fréquences).
\item Tracer des rectangles \textbf{juxtaposés} (sans espace) dont la base est l'amplitude de la classe et la hauteur l'effectif $n_i$ (ou la fréquence $f_i$).
\item \textbf{Attention} : si les classes n'ont pas la même amplitude, la hauteur du rectangle est la \cle{densité} $d_i = \dfrac{n_i}{a_i}$ (effectif divisé par amplitude) ; l'aire du rectangle est alors proportionnelle à l'effectif.
\end{enumerate}
\textbf{Polygone des effectifs cumulés croissants} :
\begin{enumerate}
\item Placer les points de coordonnées $(\text{borne supérieure de la classe}\,;\,\text{E.C.C.})$, en ajoutant le point de départ $(b_0\,;\,0)$ où $b_0$ est la borne inférieure de la première classe.
\item Relier les points par des segments de droite.
\item Pour la médiane graphique : tracer la droite horizontale d'ordonnée $\dfrac{N}{2}$, lire l'abscisse du point d'intersection avec le polygone.
\end{enumerate}
\textbf{Polygone des effectifs cumulés décroissants} :
\begin{enumerate}
\item Placer les points de coordonnées $(\text{borne inférieure de la classe}\,;\,\text{E.C.D.})$, en ajoutant le point final $(b_{k+1}\,;\,0)$ où $b_{k+1}$ est la borne supérieure de la dernière classe.
\item Relier les points par des segments de droite.
\end{enumerate}
\textbf{Propriété} : les deux polygones (croissant et décroissant) se coupent en un point d'abscisse égale à la médiane.
\end{methode}

\begin{exoresolu}[Histogramme et lecture graphique de la médiane]
On reprend la série des 40 sacs d'arachide. Construire l'histogramme des effectifs et le polygone des effectifs cumulés croissants, puis retrouver graphiquement la médiane.
\tcblower
\textbf{Histogramme} : on trace un rectangle par classe, de hauteur égale à l'effectif.
\begin{popfigure}
\begin{center}
\begin{tikzpicture}
\begin{axis}[width=0.85\linewidth,height=5.2cm, ybar, bar width=14pt,
xlabel={Masse (kg)}, ylabel={Effectif},
symbolic x coords={[40;48[,[48;56[,[56;64[,[64;72[,[72;80]},
xtick=data, x tick label style={font=\small, rotate=20, anchor=east},
ymin=0, ymax=15, ytick={0,3,5,13},
nodes near coords, every node near coord/.append style={font=\small},
axis lines=left, enlarge x limits=0.12]
\addplot[fill=popblue!60, draw=popblue] coordinates {([40;48[,3) ([48;56[,5) ([56;64[,13) ([64;72[,13) ([72;80],6)};
\end{axis}
\end{tikzpicture}
\end{center}
\legende{Histogramme des effectifs de la répartition des masses des 40 sacs.}
\end{popfigure}

\textbf{Polygone des E.C.C.} : on place les points $(40;0)$, $(48;3)$, $(56;8)$, $(64;21)$, $(72;34)$, $(80;40)$ et on les relie.
\begin{popfigure}
\begin{center}
\begin{tikzpicture}
\begin{axis}[width=0.85\linewidth,height=5.6cm,
xlabel={Masse (kg)}, ylabel={E.C.C.},
xmin=38, xmax=82, ymin=0, ymax=42,
xtick={40,48,56,64,72,80}, ytick={0,10,20,30,40},
grid=major, grid style={popline!40}, axis lines=left]
\addplot[popcurve, mark=*, mark options={fill=poporange}] coordinates {(40,0) (48,3) (56,8) (64,21) (72,34) (80,40)};
\addplot[dashed, poppink, thick] coordinates {(38,20) (63.4,20)};
\addplot[dashed, poppink, thick] coordinates {(63.4,20) (63.4,0)};
\node[popdark, font=\small, anchor=west] at (axis cs:41,22) {$\dfrac{N}{2}=20$};
\node[popdark, font=\small, anchor=north] at (axis cs:63.4,-1) {$Me \approx 63,4$};
\end{axis}
\end{tikzpicture}
\end{center}
\legende{Polygone des effectifs cumulés croissants et détermination graphique de la médiane.}
\end{popfigure}

\textbf{Lecture graphique} : la droite d'ordonnée $\dfrac{N}{2}=20$ coupe le polygone en un point d'abscisse comprise entre 56 et 64, très proche de $63,4$ : on retrouve la médiane calculée par interpolation linéaire. \checkmark

\textbf{Conclusion} : l'histogramme montre que les masses se concentrent dans les classes $[56;64[$ et $[64;72[$, et la médiane graphique confirme $Me \approx 63,4$ kg.
\end{exoresolu}

\courssub{Interprétation et communication des résultats d'une enquête}

\begin{methode}[Comparer et interpréter deux séries statistiques]
\begin{enumerate}
\item Comparer les \textbf{moyennes} (ou médianes) pour juger du \emph{niveau} de chaque série.
\item Comparer les \textbf{écarts-types} pour juger de l'\emph{homogénéité} (ou régularité) : le plus petit écart-type correspond à la série la plus régulière, la plus homogène.
\item Rédiger la conclusion en langage courant, adaptée au contexte de l'énoncé, sans jamais se contenter de donner les nombres.
\end{enumerate}
\textbf{Réflexe} : une moyenne élevée avec un grand écart-type peut cacher de fortes inégalités ; toujours croiser les deux paramètres.
\end{methode}

\begin{exoresolu}[Comparaison de deux classes]
Deux classes de Première A d'un lycée de Bafoussam ont passé le même devoir. Les résultats sont résumés ci-dessous :
\begin{center}
\begin{tabular}{|l|c|c|}
\hline
\rowcolor{popblueL} & Moyenne & Écart-type \\ \hline
Première A1 & 11,5 & 4,2 \\ \hline
Première A2 & 11,2 & 1,8 \\ \hline
\end{tabular}
\end{center}
Le proviseur doit désigner la classe qui représentera le lycée à un concours inter-établissements où la régularité des copies est primordiale. Quelle classe conseiller ? Justifier.
\tcblower
\textbf{Comparaison des moyennes} : $11,5 > 11,2$, donc la Première A1 a un niveau moyen très légèrement supérieur (écart de seulement $0,3$ point, faible).

\textbf{Comparaison des écarts-types} : $1,8 < 4,2$, donc les notes de la Première A2 sont beaucoup moins dispersées autour de la moyenne : cette classe est nettement plus \textbf{homogène}. Dans la Première A1, l'écart-type de $4,2$ traduit de fortes inégalités : d'excellentes copies côtoient des copies très faibles.

\textbf{Conclusion} : les moyennes étant quasi identiques, c'est le critère de régularité qui doit décider. On conseille au proviseur de désigner la \textbf{Première A2}, dont l'écart-type plus faible ($1,8$ contre $4,2$) garantit des copies plus régulières.
\end{exoresolu}

\courssub{Compléments utiles pour la Terminale}

\begin{savaistu}[Vers la Terminale : coefficient de variation et boîte à moustaches]
En Terminale, on approfondit l'analyse de la dispersion avec :
\begin{itemize}
\item Le \textbf{coefficient de variation} $CV = \dfrac{\sigma}{\bar{x}} \times 100$ (en \%), qui permet de comparer la dispersion de deux séries de moyennes très différentes (ex: masses d'éléphants vs masses de souris).
\item Les \textbf{quartiles} $Q_1$ (premier quartile, 25\% des données en dessous) et $Q_3$ (troisième quartile, 75\% des données en dessous), obtenus par interpolation linéaire aux rangs $N/4$ et $3N/4$.
\item L'\cle{écart interquartile} $EI = Q_3 - Q_1$, mesure de dispersion robuste (insensible aux valeurs extrêmes).
\item La \textbf{boîte à moustaches} (box-plot) : représentation graphique synthétique utilisant $Min$, $Q_1$, $Me$, $Q_3$, $Max$, idéale pour comparer visuellement plusieurs séries.
\end{itemize}
Ces outils permettent une analyse bien plus fine que la seule variance, surtout en présence de valeurs aberrantes.
\end{savaistu}

\begin{attention}[Les 5 erreurs qui coûtent des points au Probatoire]
\begin{enumerate}
\item \textbf{Oublier les effectifs dans la moyenne.} La moyenne d'une série avec effectifs n'est pas $\dfrac{\sum x_i}{p}$ mais $\dfrac{\sum n_i x_i}{N}$. De même, pour des classes, on utilise les \textbf{centres} $c_i$, jamais les bornes.
\item \textbf{Se tromper de rang pour la médiane.} Pour $N$ pair, la médiane est la demi-somme des valeurs de rangs $\dfrac{N}{2}$ et $\dfrac{N}{2}+1$, et non la valeur de rang $\dfrac{N}{2}$. Et il faut avoir \textbf{rangé la série} dans l'ordre croissant avant toute recherche.
\item \textbf{Confondre variance et écart-type.} L'écart-type est $\sigma=\sqrt{V}$ : oublier la racine carrée fait perdre le point de la question. La variance est toujours \textbf{positive} ; un résultat négatif signale une erreur de calcul (souvent dans König-Huygens : c'est $\dfrac{\sum n_i x_i^2}{N}-\bar{x}^2$, pas l'inverse).
\item \textbf{Mal construire le tableau cumulé.} Le dernier effectif cumulé croissant doit être égal à $N$ et la somme des fréquences doit valoir $1$ (ou $100\,\%$) : ces deux vérifications doivent apparaître sur la copie.
\item \textbf{Tracer un histogramme avec des rectangles espacés ou des hauteurs fausses.} Les rectangles sont \textbf{juxtaposés} (le caractère est continu), les bornes sont placées en abscisse à échelle régulière, et pour des classes d'amplitudes inégales ce sont les \textbf{aires} qui sont proportionnelles aux effectifs (hauteur = densité). Ne pas oublier non plus la légende des axes, les unités et le titre du graphique.
\end{enumerate}
\end{attention}

\newpage

\coursec{Exercices}

\courssub{Je vérifie mes connaissances}

\begin{exercice}[QCM : les bons réflexes]
Pour chaque question, une seule réponse est correcte. Entoure la bonne lettre.
\begin{enumerate}
\item Une série statistique de 80 élèves est regroupée en classes $[0\,; 5[$, $[5\,; 10[$, $[10\,; 15[$, $[15\,; 20]$. Le centre de la classe $[10\,; 15[$ est :
\begin{enumerate}
\item[a)] $10$ \quad b) $12$ \quad c) $12,5$ \quad d) $15$
\end{enumerate}
\item L'effectif cumulé croissant de la classe $[5\,; 10[$ représente :
\begin{enumerate}
\item[a)] le nombre d'individus de la classe $[5\,; 10[$ ;
\item[b)] le nombre d'individus dont la valeur du caractère est strictement inférieure à $10$ ;
\item[c)] le nombre d'individus dont la valeur du caractère est supérieure à $5$ ;
\item[d)] la fréquence de la classe $[5\,; 10[$.
\end{enumerate}
\item La médiane d'une série d'effectif total $N = 101$ est :
\begin{enumerate}
\item[a)] la valeur du $50^{\text{e}}$ individu ;
\item[b)] la valeur du $51^{\text{e}}$ individu ;
\item[c)] la moyenne des valeurs du $50^{\text{e}}$ et du $51^{\text{e}}$ individu ;
\item[d)] la valeur du $50,5^{\text{e}}$ individu.
\end{enumerate}
\item Si tous les écarts à la moyenne sont nuls, alors :
\begin{enumerate}
\item[a)] la variance est égale à la moyenne ;
\item[b)] l'écart-type vaut $1$ ;
\item[c)] la variance est nulle ;
\item[d)] la série n'existe pas.
\end{enumerate}
\end{enumerate}
\end{exercice}

\begin{exercice}[Vrai ou faux, en justifiant]
Réponds par VRAI ou FAUX en justifiant chaque réponse par une phrase ou un contre-exemple.
\begin{enumerate}
\item Deux séries ayant la même moyenne ont nécessairement le même écart-type.
\item La somme des fréquences de toutes les classes d'une série est toujours égale à $1$.
\item Dans un histogramme à classes d'égales amplitudes, les aires des rectangles sont proportionnelles aux effectifs.
\item La médiane d'une série peut être calculée par la formule $\overline{x} = \dfrac{\sum n_i x_i}{N}$.
\item Si on ajoute $3$ à chaque valeur d'une série, la moyenne augmente de $3$ mais l'écart-type ne change pas.
\end{enumerate}
\end{exercice}

\begin{exercice}[Définitions et vocabulaire]
\begin{enumerate}
\item Définis les termes suivants : \cle{population}, \cle{caractère}, \cle{modalité}, \cle{effectif cumulé croissant}.
\item Quelle différence fais-tu entre un caractère statistique discret et un caractère continu ? Donne un exemple de chaque, tiré de la vie scolaire.
\item Écris la formule de la moyenne d'une série regroupée en classes, en précisant la signification de chaque symbole utilisé.
\item Écris les deux formules de la variance (formule de définition et formule de König-Huygens) et précise celle qui est la plus pratique en calcul.
\end{enumerate}
\end{exercice}

\begin{exercice}[Calculs directs]
On a relevé le nombre de buts marqués par match lors d'un tournoi inter-établissements de Yaoundé :
\begin{center}
\begin{tabular}{|l|ccccc|}
\hline
\rowcolor{popblueL} Nombre de buts $x_i$ & $0$ & $1$ & $2$ & $3$ & $4$ \\
\hline
Nombre de matchs $n_i$ & $2$ & $5$ & $6$ & $4$ & $3$ \\
\hline
\end{tabular}
\end{center}
\begin{enumerate}
\item Détermine l'effectif total $N$.
\item Calcule la moyenne $\overline{x}$ de cette série.
\item Dresse le tableau des effectifs cumulés croissants et déduis-en la médiane.
\item Calcule la variance puis l'écart-type (arrondi au centième).
\end{enumerate}
\end{exercice}

\courssub{Je m'entraîne}

\begin{exercice}[Regroupement en classes]
Voici les durées (en minutes) de $30$ appels passés un soir depuis une cabine téléphonique de Douala :
\begin{center}
$2$ ; $5$ ; $8$ ; $1$ ; $12$ ; $7$ ; $3$ ; $9$ ; $15$ ; $4$ ; $6$ ; $11$ ; $2$ ; $8$ ; $10$ ; $5$ ; $13$ ; $1$ ; $7$ ; $9$ ; $3$ ; $6$ ; $14$ ; $4$ ; $8$ ; $11$ ; $2$ ; $10$ ; $5$ ; $7$.
\end{center}
\begin{enumerate}
\item Regroupe ces données en classes d'amplitude $4$ en partant de $0$ : $[0\,; 4[$, $[4\,; 8[$, etc.
\item Dresse le tableau complet : effectifs, fréquences (en \%), effectifs cumulés croissants et décroissants.
\item Quel est le pourcentage d'appels ayant duré moins de $12$ minutes ?
\end{enumerate}
\end{exercice}

\begin{exercice}[Moyenne, médiane et dispersion d'une série en classes]
Une coopérative agricole de l'Ouest a pesé $50$ sacs de maïs. Les masses (en kg) sont regroupées ainsi :
\begin{center}
\begin{tabular}{|l|ccccc|}
\hline
\rowcolor{popblueL} Masse (kg) & $[45\,; 50[$ & $[50\,; 55[$ & $[55\,; 60[$ & $[60\,; 65[$ & $[65\,; 70[$ \\
\hline
Effectif $n_i$ & $6$ & $12$ & $18$ & $10$ & $4$ \\
\hline
\end{tabular}
\end{center}
\begin{enumerate}
\item Calcule les centres $c_i$ des classes.
\item Calcule la masse moyenne d'un sac.
\item Détermine la classe médiane, puis calcule la médiane par interpolation linéaire.
\item Calcule la variance et l'écart-type (arrondis au centième). Interprète l'écart-type par une phrase.
\end{enumerate}
\end{exercice}

\begin{exercice}[Comparaison de deux classes]
Les notes de mathématiques de deux classes de Première A de $40$ élèves chacune ont donné les résultats suivants :
\begin{center}
\begin{tabular}{|l|c|c|}
\hline
\rowcolor{popblueL} Paramètre & Première A1 & Première A2 \\
\hline
Moyenne $\overline{x}$ & $10,4$ & $10,4$ \\
\hline
Écart-type $\sigma$ & $1,8$ & $4,6$ \\
\hline
\end{tabular}
\end{center}
\begin{enumerate}
\item Peut-on dire que les deux classes ont le même niveau ? Justifie.
\item Laquelle des deux classes est la plus homogène ? Explique le rôle de l'écart-type.
\item Le proviseur souhaite constituer une classe de renforcement avec les élèves les plus en difficulté. Quelle classe fournira vraisemblablement le plus d'élèves de très faible niveau ? Pourquoi ?
\end{enumerate}
\end{exercice}

\begin{exercice}[Chasse aux erreurs]
Un élève a construit le tableau suivant à partir d'une enquête sur les distances (en km) parcourues par $60$ moto-taxis de Bafoussam :
\begin{center}
\begin{tabular}{|l|c|c|}
\hline
\rowcolor{popblueL} Distance (km) & Effectif & Effectif cumulé croissant \\
\hline
$[0\,; 20]$ & $10$ & $10$ \\
\hline
$[20\,; 40]$ & $18$ & $30$ \\
\hline
$[40\,; 60]$ & $20$ & $48$ \\
\hline
$[60\,; 80]$ & $12$ & $62$ \\
\hline
\end{tabular}
\end{center}
Il affirme ensuite : « la variance vaut $\sigma^2 = \dfrac{1}{60}\sum n_i c_i^2 - \overline{x}$ ».
\begin{enumerate}
\item Repère et corrige toutes les erreurs de construction du tableau (bornes des classes, effectifs cumulés).
\item Corrige la formule de la variance écrite par l'élève.
\item Avec le tableau corrigé, calcule la moyenne et la médiane par interpolation linéaire.
\end{enumerate}
\end{exercice}

\courssub{Je me prépare au Probatoire}

\begin{exercice}[Enquête CAMTEL — exercice type Probatoire]
Une agence CAMTEL de Yaoundé a relevé le temps d'attente (en minutes) de $100$ clients avant d'être servis. Les résultats sont consignés dans le tableau ci-dessous :
\begin{center}
\begin{tabular}{|l|ccccc|}
\hline
\rowcolor{popblueL} Temps d'attente (min) & $[0\,; 10[$ & $[10\,; 20[$ & $[20\,; 30[$ & $[30\,; 40[$ & $[40\,; 50[$ \\
\hline
Effectif $n_i$ & $12$ & $28$ & $35$ & $17$ & $8$ \\
\hline
\end{tabular}
\end{center}
\begin{enumerate}
\item Reproduis et complète le tableau avec : les centres des classes, les fréquences, les effectifs cumulés croissants.
\item Calcule le temps d'attente moyen d'un client (arrondi au dixième de minute).
\item \begin{enumerate}
\item Détermine la classe médiane.
\item Calcule la médiane par interpolation linéaire (arrondie au dixième).
\end{enumerate}
\item Calcule la variance puis l'écart-type de la série (arrondis au centième).
\item Construis l'histogramme des effectifs (unités : $1$ cm pour $10$ min en abscisse, $1$ cm pour $5$ clients en ordonnée).
\item \begin{enumerate}
\item Construis le polygone des effectifs cumulés croissants.
\item Détermine graphiquement la médiane et compare avec le résultat de la question 3.b).
\end{enumerate}
\item Le chef d'agence affirme : « plus de la moitié des clients attendent moins de $25$ minutes ». À l'aide du polygone des effectifs cumulés croissants, dis si cette affirmation est exacte.
\end{enumerate}
\end{exercice}

\begin{exercice}[Récoltes de cacao dans le Centre]
Un acheteur de cacao a mesuré la production annuelle (en tonnes) de $80$ planteurs d'une localité de la région du Centre :
\begin{center}
\begin{tabular}{|l|ccccc|}
\hline
\rowcolor{popblueL} Production (t) & $[0\,; 2[$ & $[2\,; 4[$ & $[4\,; 6[$ & $[6\,; 8[$ & $[8\,; 10[$ \\
\hline
Effectif $n_i$ & $8$ & $20$ & $a$ & $18$ & $6$ \\
\hline
\end{tabular}
\end{center}
où $a$ désigne un effectif inconnu.
\begin{enumerate}
\item Détermine la valeur de $a$.
\item Construis le tableau des fréquences et des fréquences cumulées croissantes (en \%).
\item Calcule la production moyenne d'un planteur.
\item \begin{enumerate}
\item Détermine la médiane par interpolation linéaire.
\item Interprète cette médiane par une phrase compréhensible par un non-spécialiste.
\end{enumerate}
\item Calcule la variance et l'écart-type (arrondis au centième).
\item Un acheteur hésite entre cette localité et une autre où la production moyenne est de $4,8$ tonnes avec un écart-type de $0,9$ tonne. Dans quelle localité les productions sont-elles les plus homogènes ? Justifie et conclus sur le choix le plus sûr pour l'acheteur.
\end{enumerate}
\end{exercice}

\begin{exercice}[Lecture graphique et rédaction d'un rapport]
Dans le cadre de la semaine du livre, le club lecture d'un lycée de Garoua a mené une enquête sur le nombre de livres lus par trimestre par les élèves. Le polygone des effectifs cumulés croissants est donné ci-dessous : les points de coordonnées $(0\,; 0)$, $(2\,; 15)$, $(4\,; 45)$, $(6\,; 90)$, $(8\,; 110)$ et $(10\,; 120)$ sont reliés par des segments de droite.
\begin{enumerate}
\item Reproduis le graphique (unités : $1$ cm pour $2$ livres en abscisse, $1$ cm pour $10$ élèves en ordonnée).
\item \begin{enumerate}
\item Détermine graphiquement la médiane de la série.
\item Détermine graphiquement les quartiles $Q_1$ et $Q_3$ (rappel : $Q_1$ correspond à $25\,\%$ de l'effectif, $Q_3$ à $75\,\%$).
\end{enumerate}
\item Reconstitue le tableau des effectifs par classes $[0\,; 2[$, $[2\,; 4[$, \ldots, $[8\,; 10]$.
\item Vérifie par le calcul (interpolation linéaire) la valeur de la médiane trouvée graphiquement.
\item Calcule la moyenne et l'écart-type de la série (arrondis au centième).
\item Rédige une conclusion de cinq lignes environ, destinée au proviseur, présentant les résultats de l'enquête (niveau de lecture moyen, dispersion, proportion d'élèves lisant peu) et proposant une recommandation pour le club lecture.
\end{enumerate}
\end{exercice}

\newpage

\coursec{Corrigés des exercices}

\courssub{Corrigés des exercices d'application}

\begin{corrige}[Exercice 1 — QCM : les bons réflexes]
\begin{enumerate}
\item \textbf{c) $12,5$}. Le centre d'une classe $[a\,; b[$ est $\dfrac{a+b}{2}$. Ici $\dfrac{10+15}{2}=12,5$.
\item \textbf{b) le nombre d'individus dont la valeur du caractère est strictement inférieure à $10$}. L'effectif cumulé croissant d'une classe donne le nombre d'individus dont la valeur est \emph{strictement inférieure} à la borne supérieure de cette classe.
\item \textbf{b) la valeur du $51^{\text{e}}$ individu}. Pour un effectif total impair $N=101$, la médiane est la valeur de l'individu de rang $\dfrac{N+1}{2}=51$.
\item \textbf{c) la variance est nulle}. Si tous les écarts à la moyenne sont nuls, alors chaque $x_i=\overline{x}$, donc $\sigma^2=\frac{1}{N}\sum n_i(x_i-\overline{x})^2=0$.
\end{enumerate}
\end{corrige}

\begin{corrige}[Exercice 2 — Vrai ou faux, en justifiant]
\begin{enumerate}
\item \textbf{FAUX}. Deux séries peuvent avoir la même moyenne mais des dispersions très différentes. Exemple : série A : $\{10,10,10\}$ (moyenne $10$, $\sigma=0$) ; série B : $\{0,10,20\}$ (moyenne $10$, $\sigma\approx8,16$).
\item \textbf{VRAI}. Par définition, la fréquence d'une classe est $f_i=\frac{n_i}{N}$. La somme des fréquences est $\sum f_i = \frac{\sum n_i}{N}=1$ (ou $100\%$).
\item \textbf{VRAI}. Dans un histogramme à classes d'égales amplitudes, la largeur des rectangles est constante. L'aire d'un rectangle est $ \text{largeur} \times \text{hauteur} $. La hauteur est proportionnelle à l'effectif (ou à la fréquence), donc l'aire est proportionnelle à l'effectif.
\item \textbf{FAUX}. La formule $\overline{x} = \frac{\sum n_i x_i}{N}$ calcule la \emph{moyenne}, pas la médiane. La médiane se détermine à partir des effectifs cumulés.
\item \textbf{VRAI}. L'ajout d'une constante $c$ à toutes les valeurs décale la moyenne de $c$ ($\overline{x+c}=\overline{x}+c$) mais ne modifie pas les écarts à la moyenne, donc la variance et l'écart-type restent inchangés.
\end{enumerate}
\end{corrige}

\begin{corrige}[Exercice 3 — Définitions et vocabulaire]
\begin{enumerate}
\item \begin{itemize}
\item \cle{Population} : ensemble des individus (ou objets) sur lesquels porte l'étude statistique.
\item \cle{Caractère} : propriété qualitative ou quantitative observée sur chaque individu de la population.
\item \cle{Modalité} : valeur que peut prendre le caractère (pour un caractère quantitatif discret) ou intervalle de valeurs (classe pour un caractère continu regroupé).
\item \cle{Effectif cumulé croissant} : pour une modalité (ou une classe) donnée, somme des effectifs de cette modalité et de toutes celles qui la précèdent dans l'ordre croissant.
\end{itemize}
\item Un \cle{caractère discret} ne prend qu'un nombre fini ou dénombrable de valeurs isolées (ex. : nombre d'élèves par classe dans un établissement). Un \cle{caractère continu} peut prendre toute valeur dans un intervalle de $\mathbb{R}$ (ex. : taille en cm des élèves de Première A).
\item Pour une série regroupée en classes de centres $c_i$ et d'effectifs $n_i$ ($i=1,\dots,k$), la moyenne est
\[
\overline{x} = \frac{\sum_{i=1}^k n_i c_i}{N}, \qquad N=\sum_{i=1}^k n_i.
\]
$c_i$ est le centre de la classe $i$ (milieu de l'intervalle), $n_i$ son effectif, $N$ l'effectif total.
\item \begin{itemize}
\item \textbf{Formule de définition} : $\displaystyle \sigma^2 = \frac{1}{N}\sum_{i=1}^k n_i (x_i - \overline{x})^2$ (ou avec les centres $c_i$ pour les classes).
\item \textbf{Formule de König–Huygens} : $\displaystyle \sigma^2 = \frac{1}{N}\sum_{i=1}^k n_i x_i^2 - \overline{x}^2$ (ou $\frac{\sum n_i c_i^2}{N} - \overline{x}^2$ pour les classes).
\end{itemize}
La formule de König–Huygens est plus pratique en calcul car elle évite de calculer chaque écart $(x_i-\overline{x})$.
\end{enumerate}
\end{corrige}

\begin{corrige}[Exercice 4 — Calculs directs]
\begin{enumerate}
\item Effectif total : $N = 2+5+6+4+3 = 20$ matchs.
\item Moyenne :
\[
\overline{x} = \frac{0\times2 + 1\times5 + 2\times6 + 3\times4 + 4\times3}{20}
= \frac{0+5+12+12+12}{20} = \frac{41}{20} = 2,05 \text{ buts par match}.
\]
\item Tableau des effectifs cumulés croissants :
\begin{center}
\begin{tabular}{|c|ccccc|}
\hline
$ x_i $ & $0$ & $1$ & $2$ & $3$ & $4$ \\
\hline
$ n_i $ & $2$ & $5$ & $6$ & $4$ & $3$ \\
\hline
$ N_i $ (cumulé) & $2$ & $7$ & $13$ & $17$ & $20$ \\
\hline
\end{tabular}
\end{center}
$N=20$ (pair) $\Rightarrow$ la médiane est la moyenne des valeurs des $10^{\text{e}}$ et $11^{\text{e}}$ individus. Les deux se situent dans la modalité $2$ (car $N_1=7 < 10,11 \le N_2=13$). Donc $\text{Médiane} = 2$ buts.
\item Variance (formule de König–Huygens) :
\[
\sum n_i x_i^2 = 0^2\!\times2 + 1^2\!\times5 + 2^2\!\times6 + 3^2\!\times4 + 4^2\!\times3 = 0+5+24+36+48 = 113.
\]
\[
\overline{x}^2 = (2,05)^2 = 4,2025.
\]
\[
\sigma^2 = \frac{113}{20} - 4,2025 = 5,65 - 4,2025 = 1,4475.
\]
Écart-type : $\sigma = \sqrt{1,4475} \approx 1,2031 \approx 1,20$ buts (arrondi au centième).
\end{enumerate}
\end{corrige}

\begin{corrige}[Exercice 5 — Regroupement en classes]
\begin{enumerate}
\item Regroupement en classes d'amplitude $4$ à partir de $0$ :
\begin{center}
\begin{tabular}{|c|c|c|}
\hline
Classe & Effectif $n_i$ & Fréquence $f_i$ (\%)\\
\hline
$[0\,; 4[$ & $7$ & $23,33$ \\
$[4\,; 8[$ & $10$ & $33,33$ \\
$[8\,; 12[$ & $9$ & $30,00$ \\
$[12\,; 16[$ & $4$ & $13,33$ \\
\hline
Total & $30$ & $100$ \\
\hline
\end{tabular}
\end{center}
\item Tableau complet avec effectifs cumulés :
\begin{center}
\begin{tabular}{|c|c|c|c|c|}
\hline
Classe & $n_i$ & $f_i$ (\%) & $N_i$ (croissant) & $N_i'$ (décroissant) \\
\hline
$[0\,; 4[$ & $7$ & $23,33$ & $7$ & $30$ \\
$[4\,; 8[$ & $10$ & $33,33$ & $17$ & $23$ \\
$[8\,; 12[$ & $9$ & $30,00$ & $26$ & $13$ \\
$[12\,; 16[$ & $4$ & $13,33$ & $30$ & $4$ \\
\hline
\end{tabular}
\end{center}
\item Pourcentage d'appels $< 12$ min : classes $[0\,;4[$, $[4\,;8[$, $[8\,;12[$. Effectif cumulé croissant à la fin de $[8\,;12[$ est $26$. Pourcentage $= \frac{26}{30}\times100 \approx 86,67\%$.
\end{enumerate}
\end{corrige}

\begin{corrige}[Exercice 6 — Moyenne, médiane et dispersion d'une série en classes]
\begin{enumerate}
\item Centres des classes :
$c_1 = \frac{45+50}{2}=47,5$ ; $c_2=52,5$ ; $c_3=57,5$ ; $c_4=62,5$ ; $c_5=67,5$ (en kg).
\item Masse moyenne :
\[
\overline{x} = \frac{6\times47,5 + 12\times52,5 + 18\times57,5 + 10\times62,5 + 4\times67,5}{50}
= \frac{285 + 630 + 1035 + 625 + 270}{50} = \frac{2845}{50} = 56,9 \text{ kg}.
\]
\item Médiane : $N=50$. Pour une série regroupée en classes, la médiane se détermine par interpolation au rang $N/2 = 25$. Effectifs cumulés : $6,\;18,\;36,\;46,\;50$. La classe médiane est $[55\,;60[$ (car $18 < 25 \le 36$). Interpolation linéaire :
\[
\text{Médiane} = 55 + \frac{25 - 18}{18}\times 5 = 55 + \frac{7}{18}\times5 = 55 + \frac{35}{18} \approx 55 + 1,944 = 56,94 \text{ kg}.
\]
\item Variance (König–Huygens) :
\[
\sum n_i c_i^2 = 6\times2256,25 + 12\times2756,25 + 18\times3306,25 + 10\times3906,25 + 4\times4556,25 = 163412,5.
\]
\[
\frac{\sum n_i c_i^2}{N} = \frac{163412,5}{50} = 3268,25.
\]
\[
\sigma^2 = 3268,25 - (56,9)^2 = 3268,25 - 3237,61 = 30,64.
\]
Écart-type : $\sigma = \sqrt{30,64} \approx 5,54$ kg (arrondi au centième).
\emph{Interprétation} : Les masses des sacs s'écartent en moyenne d'environ $5,5$ kg de la masse moyenne $56,9$ kg.
\end{enumerate}
\end{corrige}

\begin{corrige}[Exercice 7 — Comparaison de deux classes]
\begin{enumerate}
\item \textbf{Non}, on ne peut pas dire que les deux classes ont le même niveau. Bien que les moyennes soient identiques ($10,4$), les dispersions sont très différentes : l'écart-type de la classe A2 ($4,6$) est plus du double de celui de la classe A1 ($1,8$). La classe A2 est beaucoup plus hétérogène : elle compte à la fois des élèves très forts et des élèves très faibles, alors que la classe A1 est plus homogène autour de la moyenne.
\item La classe \textbf{Première A1} est la plus homogène. L'écart-type mesure la dispersion des notes autour de la moyenne : plus il est petit, plus les valeurs sont regroupées près de la moyenne.
\item La classe \textbf{Première A2} fournira vraisemblablement plus d'élèves de très faible niveau. Avec un écart-type plus grand, la distribution est plus étalée ; la queue de distribution vers les notes basses est plus fournie. En supposant une forme de distribution similaire (ex. approximativement normale), la proportion d'élèves en dessous d'un seuil bas (ex. $5$) est plus grande pour un écart-type plus élevé.
\end{enumerate}
\end{corrige}

\begin{corrige}[Exercice 8 — Chasse aux erreurs]
\begin{enumerate}
\item \textbf{Erreurs du tableau} :
\begin{itemize}
\item Les classes sont mal définies : elles se chevauchent aux bornes ($[0;20]$ et $[20;40]$ contiennent toutes deux la valeur $20$). Il faut utiliser des intervalles semi-ouverts : $[0\,;20[$, $[20\,;40[$, $[40\,;60[$, $[60\,;80]$ (ou $[60\,;80[$).
\item Les effectifs cumulés croissants sont faux : 
$10+18=28$ (et non $30$) ; $28+20=48$ (correct) ; $48+12=60$ (et non $62$). Le total des effectifs est $60$, le dernier cumulé doit être $60$.
\end{itemize}
Tableau corrigé :
\begin{center}
\begin{tabular}{|l|c|c|}
\hline
\rowcolor{popblueL} Distance (km) & Effectif & Effectif cumulé croissant \\
\hline
$[0\,; 20[$ & $10$ & $10$ \\
$[20\,; 40[$ & $18$ & $28$ \\
$[40\,; 60[$ & $20$ & $48$ \\
$[60\,; 80]$ & $12$ & $60$ \\
\hline
\end{tabular}
\end{center}
\item \textbf{Formule de la variance corrigée} :
\[
\sigma^2 = \frac{1}{N}\sum n_i c_i^2 - \overline{x}^2 \qquad \text{(formule de König–Huygens)}.
\]
L'élève a oublié le carré sur $\overline{x}$.
\item Avec le tableau corrigé :
Centres : $c_1=10$, $c_2=30$, $c_3=50$, $c_4=70$.
Moyenne :
\[
\overline{x} = \frac{10\times10 + 18\times30 + 20\times50 + 12\times70}{60}
= \frac{100+540+1000+840}{60} = \frac{2480}{60} \approx 41,33 \text{ km}.
\]
Médiane : $N=60$, rangs $30$ et $31$. Effectifs cumulés : $10,\;28,\;48,\;60$. La classe médiane est $[40\,;60[$ ($28 < 30,5 \le 48$). Interpolation :
\[
\text{Médiane} = 40 + \frac{30 - 28}{20}\times 20 = 40 + 2 = 42 \text{ km}.
\]
\end{enumerate}
\end{corrige}

\begin{corrige}[Exercice 9 — Enquête CAMTEL — exercice type Probatoire]
\begin{enumerate}
\item Tableau complété :
\begin{center}
\begin{tabular}{|c|c|c|c|c|}
\hline
\rowcolor{popblueL} Temps (min) & Centre $c_i$ & Effectif $n_i$ & Fréquence $f_i$ & Effectif cumulé croissant $N_i$ \\
\hline
$[0\,;10[$ & $5$ & $12$ & $0,12$ & $12$ \\
$[10\,;20[$ & $15$ & $28$ & $0,28$ & $40$ \\
$[20\,;30[$ & $25$ & $35$ & $0,35$ & $75$ \\
$[30\,;40[$ & $35$ & $17$ & $0,17$ & $92$ \\
$[40\,;50[$ & $45$ & $8$ & $0,08$ & $100$ \\
\hline
Total & & $100$ & $1$ & \\
\hline
\end{tabular}
\end{center}
\item Temps d'attente moyen :
\[
\overline{x} = \frac{12\times5 + 28\times15 + 35\times25 + 17\times35 + 8\times45}{100}
= \frac{60+420+875+595+360}{100} = \frac{2310}{100} = 23,1 \text{ min}.
\]
\item \begin{enumerate}
\item Classe médiane : $N/2=50$. Effectifs cumulés : $12,\;40,\;75,\dots$ $\Rightarrow$ classe $[20\,;30[$.
\item Médiane par interpolation :
\[
\text{Médiane} = 20 + \frac{50-40}{35}\times 10 = 20 + \frac{100}{35} \approx 20 + 2,857 = 22,9 \text{ min (au dixième)}.
\]
\end{enumerate}
\item Variance et écart-type :
\[
\sum n_i c_i^2 = 12\times25 + 28\times225 + 35\times625 + 17\times1225 + 8\times2025 = 65500.
\]
\[
\frac{\sum n_i c_i^2}{N} = 655.
\]
\[
\sigma^2 = 655 - (23,1)^2 = 655 - 533,61 = 121,39.
\]
\[
\sigma = \sqrt{121,39} \approx 11,02 \text{ min (au centième)}.
\]
\item Histogramme : abscisse $1$ cm pour $10$ min $\Rightarrow$ largeur des rectangles $=1$ cm. Ordonnée $1$ cm pour $5$ clients $\Rightarrow$ hauteurs : $12/5=2,4$ cm ; $28/5=5,6$ cm ; $35/5=7$ cm ; $17/5=3,4$ cm ; $8/5=1,6$ cm.
\item \begin{enumerate}
\item Polygone des effectifs cumulés croissants : points $(0;0)$, $(10;12)$, $(20;40)$, $(30;75)$, $(40;92)$, $(50;100)$ reliés par segments.
\item Médiane graphique : abscisse du point d'ordonnée $50$ sur le polygone. Entre $(20;40)$ et $(30;75)$, on trouve $x \approx 22,9$ min, cohérent avec le calcul.
\end{enumerate}
\item Affirmation du chef : « plus de la moitié des clients attendent moins de $25$ min ». Sur le polygone, à $x=25$, l'ordonnée (effectif cumulé) est obtenue par interpolation entre $(20;40)$ et $(30;75)$ : $y = 40 + \frac{5}{10}\times35 = 57,5$. Donc $57,5$ clients sur $100$ attendent moins de $25$ min, soit $57,5\% > 50\%$. L'affirmation est \textbf{exacte}.
\end{enumerate}

\textbf{Barème indicatif (sur 20)} : Tableau complété (3 pts) ; Moyenne (2 pts) ; Classe médiane + interpolation (3 pts) ; Variance + écart-type (3 pts) ; Histogramme (3 pts) ; Polygone + médiane graphique (3 pts) ; Interprétation affirmation (3 pts).
\textbf{Points de vigilance} : Bien utiliser les centres des classes ; ne pas confondre effectifs et fréquences ; interpolation linéaire correcte (formule $\text{Médiane}=L+\frac{N/2-F_{\text{av}}}{f_{\text{med}}}\times h$) ; sur le polygone, les abscisses sont les bornes supérieures des classes ; vérifier que la somme des effectifs $=100$.
\end{corrige}

\begin{corrige}[Exercice 10 — Récoltes de cacao dans le Centre]
\begin{enumerate}
\item Effectif total $N=80$ : $8+20+a+18+6=52+a=80 \Rightarrow a=28$.
\item Tableau des fréquences et fréquences cumulées croissantes (en \%) :
\begin{center}
\begin{tabular}{|c|c|c|c|}
\hline
\rowcolor{popblueL} Classe (t) & Effectif $n_i$ & Fréquence $f_i$ (\%) & Fréquence cumulée croissante (\%) \\
\hline
$[0\,;2[$ & $8$ & $10,0$ & $10,0$ \\
$[2\,;4[$ & $20$ & $25,0$ & $35,0$ \\
$[4\,;6[$ & $28$ & $35,0$ & $70,0$ \\
$[6\,;8[$ & $18$ & $22,5$ & $92,5$ \\
$[8\,;10[$ & $6$ & $7,5$ & $100,0$ \\
\hline
Total & $80$ & $100$ & \\
\hline
\end{tabular}
\end{center}
\item Production moyenne : centres $1,\;3,\;5,\;7,\;9$.
\[
\overline{x} = \frac{8\times1 + 20\times3 + 28\times5 + 18\times7 + 6\times9}{80}
= \frac{8+60+140+126+54}{80} = \frac{388}{80} = 4,85 \text{ tonnes}.
\]
\item \begin{enumerate}
\item Médiane : $N/2=40$. Effectifs cumulés : $8,\;28,\;56,\dots$ $\Rightarrow$ classe médiane $[4\,;6[$. Interpolation (avec les effectifs) :
\[
\text{Médiane} = 4 + \frac{40 - 28}{28}\times 2 = 4 + \frac{12}{28}\times 2 = 4 + \frac{24}{28} = 4 + \frac{6}{7} \approx 4,857 \text{ tonnes}.
\]
\item \emph{Interprétation} : La moitié des planteurs de cette localité produisent moins de $4,86$ tonnes de cacao par an, et l'autre moitié produit plus.
\end{enumerate}
\item Variance et écart-type :
\[
\sum n_i c_i^2 = 8\times1 + 20\times9 + 28\times25 + 18\times49 + 6\times81 = 8+180+700+882+486 = 2256.
\]
\[
\frac{\sum n_i c_i^2}{N} = \frac{2256}{80} = 28,2.
\]
\[
\sigma^2 = 28,2 - (4,85)^2 = 28,2 - 23,5225 = 4,6775.
\]
\[
\sigma = \sqrt{4,6775} \approx 2,16 \text{ tonnes (au centième)}.
\]
\item Comparaison : Localité 1 : $\overline{x}=4,85$, $\sigma\approx2,16$. Localité 2 : $\overline{x}=4,8$, $\sigma=0,9$. L'écart-type de la localité 2 est beaucoup plus petit, donc les productions y sont \textbf{plus homogènes}. Pour l'acheteur, la localité 2 est \textbf{plus sûre} : les quantités livrées sont plus prévisibles, moins de risque de fortes variations d'un planteur à l'autre.
\end{enumerate}
\end{corrige}

\begin{corrige}[Exercice 11 — Lecture graphique et rédaction d'un rapport]
\begin{enumerate}
\item Graphique : polygone des effectifs cumulés croissants passant par $(0;0)$, $(2;15)$, $(4;45)$, $(6;90)$, $(8;110)$, $(10;120)$. Unités : $1$ cm pour $2$ livres (abscisse), $1$ cm pour $10$ élèves (ordonnée).
\item \begin{enumerate}
\item Médiane graphique : $N=120$, rang $60$. Sur le polygone, l'ordonnée $60$ se situe entre $(4;45)$ et $(6;90)$. Par interpolation : $x = 4 + \frac{60-45}{90-45}\times2 = 4 + \frac{15}{45}\times2 = 4 + \frac{2}{3} \approx 4,67$ livres.
\item Quartiles : $Q_1$ (rang $30$) entre $(2;15)$ et $(4;45)$ : $x = 2 + \frac{30-15}{30}\times2 = 2+1 = 3$ livres. $Q_3$ (rang $90$) : point $(6;90)$ exactement $\Rightarrow Q_3 = 6$ livres.
\end{enumerate}
\item Tableau des effectifs par classes (différences des cumulés) :
\begin{center}
\begin{tabular}{|c|c|}
\hline
\rowcolor{popblueL} Classe (livres) & Effectif $n_i$ \\
\hline
$[0\,;2[$ & $15$ \\
$[2\,;4[$ & $30$ \\
$[4\,;6[$ & $45$ \\
$[6\,;8[$ & $20$ \\
$[8\,;10]$ & $10$ \\
\hline
Total & $120$ \\
\hline
\end{tabular}
\end{center}
\item Vérification médiane par interpolation (classe $[4\,;6[$) :
$L=4$, $F_{\text{av}}=45$, $f_{\text{med}}=45$, $h=2$.
\[
\text{Médiane} = 4 + \frac{60-45}{45}\times2 = 4 + \frac{15}{45}\times2 = 4 + \frac{2}{3} = 4,67 \text{ livres}.
\]
\item Moyenne et écart-type : centres $1,\;3,\;5,\;7,\;9$.
\[
\overline{x} = \frac{15\times1 + 30\times3 + 45\times5 + 20\times7 + 10\times9}{120}
= \frac{15+90+225+140+90}{120} = \frac{560}{120} = \frac{14}{3} \approx 4,67 \text{ livres}.
\]
\[
\sum n_i c_i^2 = 15\times1 + 30\times9 + 45\times25 + 20\times49 + 10\times81 = 3200.
\]
\[
\sigma^2 = \frac{3200}{120} - \left(\frac{14}{3}\right)^2 = \frac{80}{3} - \frac{196}{9} = \frac{240-196}{9} = \frac{44}{9} \approx 4,89.
\]
\[
\sigma = \sqrt{\frac{44}{9}} = \frac{2}{3}\sqrt{11} \approx 2,21 \text{ livres (au centième)}.
\]
\item \textbf{Conclusion pour le proviseur} :
L'enquête montre que les élèves lisent en moyenne $4,67$ livres par trimestre, avec un écart-type de $2,21$ livres, ce qui indique une dispersion modérée. La médiane ($4,67$ livres) confirme que la moitié des élèves lit moins de $5$ livres. Les quartiles $Q_1=3$ et $Q_3=6$ révèlent que $25\%$ des élèves lisent au maximum $3$ livres (lecteurs faibles) et $25\%$ en lisent au moins $6$ (bons lecteurs). Il est recommandé au club lecture de cibler prioritairement le quartile inférieur en proposant des animations (lectures à voix haute, défis lecture) et en diversifiant les ouvrages accessibles pour stimuler les élèves les plus en retrait.
\end{enumerate}
\end{corrige}

\newpage

\coursec{Fiche bilan}

\begin{popfigure}
\begin{tikzpicture}[
    mindmap,
    grow cyclic,
    every node/.style={concept, execute at begin node=\hskip0pt},
    concept color=popblue,
    level 1/.append style={level distance=4.5cm, sibling angle=360/6},
    level 2/.append style={level distance=3cm, sibling angle=45},
    texte court/.style={text width=2.8cm, align=center, font=\small\bfseries},
    branch1/.style={concept color=poporange},
    branch2/.style={concept color=popgreen},
    branch3/.style={concept color=poppurple},
    branch4/.style={concept color=poppink},
    branch5/.style={concept color=popgold},
    branch6/.style={concept color=popblue!70!black}
]
\node[texte court, minimum size=1.8cm, align=center] {Statistiques\\ Organisation \& Gestion\\ des Données}
    child[branch1] { node[texte court, align=center] {1. Organisation\\ des données\\ (Classes, Effectifs,\\ Fréquences, Cumulés)} }
    child[branch2] { node[texte court, align=center] {2. Position :\\ Moyenne\\ (Formules, Propriétés)} }
    child[branch3] { node[texte court, align=center] {3. Position :\\ Médiane\\ (Tableau, Polygone,\\ Interpolation)} }
    child[branch4] { node[texte court, align=center] {4. Dispersion :\\ Variance \&\\ Écart-type\\ (Koenig-Huygens)} }
    child[branch5] { node[texte court, align=center] {5. Représentations\\ Graphiques\\ (Histogramme,\\ Polygones cumulés)} }
    child[branch6] { node[texte court, align=center] {6. Interprétation\\ \& Communication\\ (Contexte réel)} };
\end{tikzpicture}
\legende{Carte mentale de la leçon NA06 : Statistiques (1ère A)}
\end{popfigure}

\begin{bilan}

\textbf{Définitions clés}
\begin{itemize}
    \item \cle{Population} : ensemble des individus étudiés (taille $N$).
    \item \cle{Caractère statistique} : variable observée (quantitatif discret/continu, qualitatif).
    \item \cle{Modalité} $x_i$ : valeur prise par le caractère ; \cle{Effectif} $n_i$ : nombre d'individus pour $x_i$.
    \item \cle{Fréquence} $f_i = \frac{n_i}{N}$ ; $\sum n_i = N$, $\sum f_i = 1$.
    \item \cle{Classes d'égales amplitudes} : intervalles $[a_k; b_k[$ de même longueur $h = b_k - a_k$ (dernière classe fermée). \cle{Modalité de classe} = centre $c_k = \frac{a_k+b_k}{2}$.
    \item \cle{Effectifs cumulés croissants} $N_k = \sum_{i=1}^k n_i$ ; \cle{décroissants} $N'_k = \sum_{i=k}^m n_i$. Idem pour fréquences $F_k, F'_k$.
\end{itemize}

\textbf{Formules à connaître par cœur}
\begin{center}
\begin{tabularx}{\linewidth}{|l|X|X|}\hline
\rowcolor{popblueL} \textbf{Paramètre} & \textbf{Série simple (modalités $x_i$, effectifs $n_i$)} & \textbf{Série groupée (classes centres $c_k$, effectifs $n_k$)} \\ \hline
Moyenne $\bar{x}$ & $\displaystyle \bar{x} = \frac{\sum n_i x_i}{N}$ & $\displaystyle \bar{x} \approx \frac{\sum n_k c_k}{N}$ \\ \hline
Variance $V$ (ou $\sigma^2$) & $\displaystyle V = \frac{\sum n_i (x_i-\bar{x})^2}{N} = \frac{\sum n_i x_i^2}{N} - \bar{x}^2$ & $\displaystyle V \approx \frac{\sum n_k c_k^2}{N} - \bar{x}^2$ \\ \hline
Écart-type $\sigma$ & $\sigma = \sqrt{V}$ & $\sigma = \sqrt{V}$ \\ \hline
Médiane $M_e$ & Valeur centrale (ou moyenne des 2 centrales) si $N$ pair/impair & \textbf{Interpolation linéaire} dans la classe médiane : \\
& & $M_e = a + \frac{\frac{N}{2} - N_{\text{avant}}}{n_{\text{méd}}} \times h$ \\ \hline
\end{tabularx}
\end{center}

\textbf{Méthodes express}
\begin{enumerate}
    \item \textbf{Regrouper en classes} : Amplitude $h \approx \frac{\text{Étendue}}{\text{Nb classes souhaité (5 à 15)}} \to$ arrondir $h$ à une valeur « propre » (1, 2, 5, 10, 20...). Définir la 1ère borne $a_1 \le \min$. Construire le tableau.
    \item \textbf{Moyenne} : Tableau auxiliaire ($x_i$, $n_i$, $n_i x_i$, $n_i x_i^2$) $\to$ Somme $\to$ Formules.
    \item \textbf{Médiane} : Calculer $N/2$. Chercher la 1ère valeur où $N_k \ge N/2$ (croissant) ou tracer le polygone croissant $\to$ lire l'abscisse de l'ordonnée $N/2$.
    \item \textbf{Variance} : Toujours utiliser \textbf{Koenig-Huygens} : $V = \frac{\sum n_i x_i^2}{N} - \bar{x}^2$ (moins d'erreurs de calcul).
    \item \textbf{Histogramme} : Axe des abscisses = modalités/classes (échelle réelle). Axe des ordonnées = \cle{Densité} $d_k = \frac{n_k}{N \times h}$ (ou effectifs $n_k$ si $h$ constant). \textbf{Aire du rectangle = Effectif (ou Fréquence)}.
    \item \textbf{Polygones cumulés} : Points $(x_i, N_i)$ croissant ou $(x_i, N'_i)$ décroissant. Relier par segments. Médiane = abscisse du point d'ordonnée $N/2$.
\end{enumerate}

\end{bilan}

\begin{aretenir}[Checklist avant le Probatoire]
\begin{itemize}
    \item $\square$ Je sais définir : population, caractère, modalité, effectif, fréquence, effectif total $N$.
    \item $\square$ Je sais construire un tableau d'effectifs et de fréquences (croissantes \textbf{et} décroissantes) pour une série brute ou groupée.
    \item $\square$ Je sais choisir une amplitude $h$ « propre » et regrouper une série continue en classes d'égales amplitudes.
    \item $\square$ Je sais calculer la \textbf{moyenne} $\bar{x}$ (série simple et groupée) et j'en connais les propriétés (translation, homothétie).
    \item $\square$ Je sais déterminer la \textbf{médiane} $M_e$ par le tableau des effectifs cumulés (formule d'interpolation pour classes) \textbf{et} par lecture graphique sur le polygone cumulé croissant.
    \item $\square$ Je sais calculer la \textbf{variance} $V$ et l'\textbf{écart-type} $\sigma$ en utilisant \textbf{exclusivement} la formule de Koenig-Huygens ($V = \frac{\sum n x^2}{N} - \bar{x}^2$).
    \item $\square$ Je sais construire un \textbf{histogramme} rigoureux : axes gradués, densité (si amplitudes inégales) ou effectifs (si égales), aires proportionnelles aux effectifs.
    \item $\square$ Je sais tracer les \textbf{polygones des effectifs cumulés} (croissant et décroissant) et en lire la médiane (et les quartiles).
    \item $\square$ Je sais \textbf{interpréter} : comparer deux séries via moyennes (niveau) et écarts-types (dispersion/homogénéité) dans un énoncé concret (notes, salaires, rendements agricoles...).
    \item $\square$ Je vérifie toujours : $\sum n_i = N$, $\sum f_i = 1$, $V \ge 0$, médiane dans la classe médiane, somme des aires de l'histogramme = $N$ (ou 1).
    \item $\square$ Je soigne ma rédaction : phrases complètes, unités, conclusion interprétée (« La série B est plus homogène car $\sigma_B < \sigma_A$ »).
\end{itemize}
\end{aretenir}

\findecours

\end{document}
